% Acides et bases en solution aqueuse : eau, autoprotolyse et pH — Chimie Tle TI — Cours signé OnBuch+
% Programme officiel MINESEC, Terminale TI. Compiler avec : tectonic L06-acides-bases-solution-aqueuse.tex
\newcommand{\DOCMATIERE}{Chimie}
\newcommand{\DOCNIVEAU}{Tle TI}
\newcommand{\DOCMODULE}{Module 2 · Acides et bases}
\newcommand{\DOCLECON}{06}
\newcommand{\DOCTITRE}{Acides et bases en solution aqueuse : eau, autoprotolyse et pH}
% OnBuch+ — préambule « cours signés » ENRICHI (compilé avec tectonic / XeLaTeX)
% Système visuel « Pop » d'OnBuch+ : fond crème (jamais blanc), bordures encre
% épaisses, ombres « dures » décalées, titres Archivo Black en capitales.
% Version enrichie pour les cours de chimie : chimie (mhchem, chemfig),
% graphes (pgfplots), boîtes pédagogiques supplémentaires (définition,
% propriété, méthode, attention, le savais-tu, exercice résolu, exercice,
% TP, bilan), page de garde refaite et signature OnBuch+ ancrée en bas.
%
% Macros attendues dans main.tex AVANT \input{preamble} :
%   \DOCMATIERE  \DOCNIVEAU  \DOCMODULE  \DOCLECON (numéro)  \DOCTITRE
\documentclass[11pt]{article}

\usepackage{fontspec}
\usepackage[french]{babel}
\usepackage[a4paper,margin=2cm,top=3.4cm,bottom=2.8cm,headheight=1.4cm,footskip=1.3cm]{geometry}
\usepackage{amsmath,amssymb,amsfonts}
\usepackage[table]{xcolor} % [table] : fournit aussi \rowcolors (alternance de lignes)
\usepackage{pagecolor}
\usepackage{tikz}
\usetikzlibrary{arrows.meta,positioning,calc,shapes.geometric,shapes.misc,shapes.arrows,decorations.pathmorphing,decorations.markings,patterns,shadows,mindmap,trees,fit,backgrounds,matrix,intersections}
\usepackage{pgfplots}
\pgfplotsset{compat=1.18}
\usepackage[version=4]{mhchem}
\usepackage{chemfig}
\usepackage{enumitem}
\usepackage{fancyhdr}
% [most] charge toutes les bibliothèques tcolorbox (listings, minted, raster,
% documentation...) et gonfle inutilement la mémoire TeX sur un document long
% (~300 boîtes) : on ne charge que ce qui est réellement utilisé.
\usepackage[skins,breakable]{tcolorbox}
\usepackage{graphicx}
\usepackage{colortbl}
\usepackage{booktabs}
\usepackage{tabularx}
\usepackage{array}
\usepackage{multicol}
\usepackage{siunitx}
% parse-numbers=false : accepte aussi \SI{2,5 \times 10^{-3}}{...}
\sisetup{locale=FR,output-decimal-marker={,},group-minimum-digits=4,parse-numbers=false}
\DeclareSIUnit\ohm{\ensuremath{\symup{\Omega}}} % U+2126 absent de la police
\usepackage{qrcode}
% \degree : commande fréquemment utilisée par les modèles pour un angle ou une
% température en dehors de \SI{}{} (non fournie par les paquets chargés).
\providecommand{\degree}{^{\circ}}
\usepackage{lastpage}
\usepackage{hyperref}
\hypersetup{hidelinks}

% ============================================================
% Palette « Pop » OnBuch+ — mêmes hex que la classe P (plus_ui.dart)
% ============================================================
\definecolor{poporange}{HTML}{F59321}
\definecolor{popdark}{HTML}{E07A0C}
\definecolor{popcream}{HTML}{FFF6E8}
\definecolor{popink}{HTML}{1C1714}
\definecolor{poppurple}{HTML}{7A5AE0}
\definecolor{popgreen}{HTML}{2E9E6A}
\definecolor{poppink}{HTML}{E0457B}
\definecolor{popblue}{HTML}{2F7DE1}
\definecolor{popgold}{HTML}{C98A12}
\definecolor{popmuted}{HTML}{8A7F76}
\definecolor{popline}{HTML}{EADFD2}
\definecolor{popgray}{HTML}{8A7F76}
\colorlet{popgrey}{popgray}
% Alias tolérants (le modèle nomme parfois les couleurs différemment)
\colorlet{popwhite}{white}
\colorlet{popblack}{popink}
\colorlet{popred}{poppink}
\colorlet{popyellow}{popgold}
% Teintes claires pour fonds de tableaux / zones
\colorlet{popblueL}{popblue!10!white}
\colorlet{poporangeL}{poporange!14!white}
\colorlet{popgreenL}{popgreen!12!white}
\colorlet{poppurpleL}{poppurple!12!white}
\colorlet{poppinkL}{poppink!12!white}
\colorlet{popgoldL}{popgold!14!white}

\pagecolor{popcream}

% ============================================================
% Typographie
% ============================================================
\defaultfontfeatures{Ligatures=TeX}
\setmainfont{PlusJakartaSans-Medium.ttf}[Path=../../../fonts/,BoldFont=PlusJakartaSans-Bold.ttf]
% Maths en sans-serif (Fira Math) avec chiffres et lettres droites de Plus
% Jakarta, pour que les formules se fondent dans le texte.
\usepackage[math-style=ISO,bold-style=ISO]{unicode-math}
\setmathfont{FiraMath-Regular.otf}[Path=../../../fonts/]
\setmathfont{PlusJakartaSans-Medium.ttf}[Path=../../../fonts/,range={up/num,up/{Latin,latin},\mathup}]
\usepackage{icomma} % virgule décimale sans espace en mode math
% Caractères absents de la police de texte (sinon carré vide) : rendus par
% leur équivalent mathématique, en texte comme en formule.
\usepackage{newunicodechar}
\newunicodechar{α}{\ensuremath{\alpha}}
\newunicodechar{Α}{\ensuremath{\alpha}}
\newunicodechar{β}{\ensuremath{\beta}}
\newunicodechar{δ}{\ensuremath{\delta}}
\newunicodechar{✓}{\popcheck{popgreen}}
\newfontfamily\popdisplay{ArchivoBlack-Regular.ttf}[Path=../../../fonts/]
\newfontfamily\poplogo{SpaceGrotesk-Bold.ttf}[Path=../../../fonts/]
\newfontfamily\popsemi{PlusJakartaSans-SemiBold.ttf}[Path=../../../fonts/]
\newfontfamily\popxbold{PlusJakartaSans-ExtraBold.ttf}[Path=../../../fonts/]
\newfontfamily\popxboldls{PlusJakartaSans-ExtraBold.ttf}[Path=../../../fonts/,LetterSpace=3.0]

\setlist[enumerate]{leftmargin=1.7em,itemsep=5pt,topsep=4pt}
\setlist[itemize]{leftmargin=1.7em,itemsep=4pt,topsep=4pt,label={\color{poporange}\textbullet}}
\setlength{\parindent}{0pt}
\setlength{\parskip}{5pt}
\renewcommand{\arraystretch}{1.25}

% chemfig : liaisons un peu plus épaisses, encre Pop
\setchemfig{atom sep=2em,bond style={line width=0.9pt,popink},cram width=3pt}

% pgfplots : style Pop par défaut
\pgfplotsset{
  every axis/.append style={axis line style={popink,line width=1pt},tick style={popink},
    grid style={popline,line width=0.5pt},label style={font=\small},tick label style={font=\footnotesize}},
  popcurve/.style={poporange,line width=1.8pt,no markers,smooth},
}

% ============================================================
% Pill / squiggle
% ============================================================
\newcommand{\poppill}[3]{%
  \tikz[baseline=-0.55ex]\node[draw=popink,line width=1.2pt,rounded corners=2.6pt,%
    fill=#2,inner xsep=6.5pt,inner ysep=3pt]{%
    \popxboldls\fontsize{7}{7}\selectfont\color{#3}\MakeUppercase{#1}};%
}
\newcommand{\popsquiggle}[1]{%
  \tikz[baseline=-0.4ex]{
    \draw[#1,line width=2.6pt,line cap=round]
      plot [smooth] coordinates {(0,0.09) (0.34,0.01) (0.68,0.21) (1.02,0.01) (1.36,0.10)};
  }%
}
% Mot-clé surligné (marqueur orange sous le texte)
\newcommand{\cle}[1]{{\popxbold\color{popdark}#1}}

% ============================================================
% En-tête / pied
% ============================================================
\pagestyle{fancy}
\fancyhf{}
\renewcommand{\headrulewidth}{0pt}
\renewcommand{\footrulewidth}{0pt}
\lhead{\raisebox{-0.15ex}{\poplogo\fontsize{13}{13}\selectfont\color{popink}OnBuch\color{poporange}+}}
\rhead{\poppill{\DOCMATIERE\ \textbullet\ \DOCNIVEAU\ \textbullet\ Leçon \DOCLECON}{white}{popink}}
\lfoot{\popxbold\fontsize{7.6}{7.6}\selectfont\color{popmuted}\MakeUppercase{Cours signé OnBuch+}\ \textbullet\ {\popsemi\DOCTITRE}}
\rfoot{\tikz[baseline=-0.6ex]\node[rounded corners=8.5pt,fill=popink,minimum height=17pt,minimum width=17pt,inner xsep=5pt,inner sep=0]{\popxbold\fontsize{8}{8}\selectfont\color{popcream}\thepage\,/\,\pageref*{LastPage}};}

% ============================================================
% Titres
% ============================================================
\newcommand{\chaptitre}[2]{%
  \begin{tcolorbox}[enhanced,colback=poporange,colframe=popink,boxrule=2.6pt,arc=11pt,
      drop shadow=popink,left=15pt,right=15pt,top=12pt,bottom=11pt,before skip=3pt,after skip=18pt]
    \poppill{#2}{popink}{popcream}\par\vspace{9pt}
    {\raggedright\hyphenpenalty=10000\exhyphenpenalty=10000\popdisplay\fontsize{22}{24}\selectfont\color{popcream}\MakeUppercase{#1}\par}
  \end{tcolorbox}
}
% Section numérotée : pastille orange avec le numéro + titre Archivo
\newcounter{popsec}
\newcommand{\coursec}[1]{%
  \stepcounter{popsec}\setcounter{popsub}{0}%
  \par\vspace{16pt}\noindent
  \begin{minipage}{\linewidth}
  \tikz[baseline=-0.7ex]\node[draw=popink,line width=1.6pt,fill=poporange,rounded corners=4pt,minimum size=22pt,inner sep=0]{\popdisplay\fontsize{12}{12}\selectfont\color{popcream}\thepopsec};%
  \hspace{8pt}{\popdisplay\fontsize{15}{17}\selectfont\color{popink}\MakeUppercase{#1}}\par
  \vspace{4pt}\noindent{\color{poporange}\rule{36pt}{3.6pt}}
  \end{minipage}\par\nopagebreak\vspace{8pt}
}
\newcounter{popsub}
\newcommand{\courssub}[1]{%
  \stepcounter{popsub}%
  \par\vspace{9pt}\noindent{\popxbold\fontsize{11.5}{13}\selectfont\color{popink}{\color{poporange}\thepopsec.\thepopsub}\hspace{6pt}#1}\par\nopagebreak\vspace{4pt}
}

% ============================================================
% Boîtes pédagogiques — même recette : fond blanc, bordure épaisse colorée,
% ombre dure encre, badge plein. Titre optionnel [#1].
% ============================================================
\tcbset{popbase/.style={breakable,enhanced,colback=white,boxrule=2.3pt,arc=9pt,
  shadow={3pt}{-3pt}{0pt}{popink},left=14pt,right=14pt,top=11pt,bottom=11pt,before skip=11pt,after skip=15pt,
  fontupper=\popsemi\fontsize{10.6}{14.5}\selectfont\color{popink},
  fontlower=\popsemi\fontsize{10.6}{14.5}\selectfont\color{popink}}}
% Coche dessinée (le glyphe ✓ n'existe pas dans les polices du document)
\newcommand{\popcheck}[1]{\tikz[baseline=-0.5ex]\draw[#1,line width=1.6pt,line cap=round,line join=round](-0.13,0.0)--(-0.04,-0.09)--(0.13,0.1);}
\newcommand{\popboxhead}[3]{\poppill{#1}{#2}{white}\ifx\relax#3\relax\else\hspace{6pt}{\popxbold\fontsize{10.6}{12}\selectfont\color{popink}#3}\fi\par\vspace{7pt}}

\newtcolorbox{aretenir}[1][]{popbase,colframe=popgreen,before upper={\popboxhead{À retenir}{popgreen}{#1}}}
\newtcolorbox{exemplebox}[1][]{popbase,colframe=poporange,before upper={\popboxhead{Exemple}{poporange}{#1}}}
\newtcolorbox{definition}[1][]{popbase,colframe=popblue,before upper={\popboxhead{Définition}{popblue}{#1}}}
\newtcolorbox{propriete}[1][]{popbase,colframe=poppurple,before upper={\popboxhead{Propriété}{poppurple}{#1}}}
\newtcolorbox{methode}[1][]{popbase,colframe=popgold,before upper={\popboxhead{Méthode}{popgold}{#1}}}
\newtcolorbox{attention}[1][]{popbase,colframe=poppink,before upper={\popboxhead{Attention}{poppink}{#1}}}
\newtcolorbox{experience}[1][]{popbase,colframe=popblue,colback=popblueL,before upper={\popboxhead{Expérience}{popblue}{#1}}}
% « Le savais-tu ? » : carte violette pleine (contexte, culture, Cameroun)
\newtcolorbox{savaistu}[1][]{popbase,colback=poppurple,colframe=popink,
  fontupper=\popsemi\fontsize{10.4}{14.2}\selectfont\color{white},
  before upper={\poppill{Le savais-tu\,?}{white}{poppurple}\ifx\relax#1\relax\else\hspace{6pt}{\popxbold\color{white}#1}\fi\par\vspace{7pt}}}
% Exercice résolu : énoncé en haut, \tcblower puis la solution sur fond crème
\newcounter{popexo}\newcounter{popexr}
\newtcolorbox{exoresolu}[1][]{popbase,colframe=poporange,colbacklower=popcream,
  segmentation style={popink,line width=1.2pt,dashed},
  before upper={\stepcounter{popexr}\popboxhead{Exercice résolu \thepopexr}{poporange}{#1}},
  before lower={\poppill{Solution}{popink}{popcream}\par\vspace{6pt}}}
% Exercice (non résolu en place) — la correction est dans la partie Corrigés
\newtcolorbox{exercice}[1][]{popbase,colframe=popink,
  before upper={\stepcounter{popexo}\popboxhead{Exercice \thepopexo}{popink}{#1}}}
% Corrigé
\newtcolorbox{corrige}[1][]{popbase,colframe=popgreen,colback=popgreenL,
  before upper={\popboxhead{Corrigé}{popgreen}{#1}}}
% Objectifs / prérequis / bilan
\newtcolorbox{objectifs}{popbase,colframe=popink,colback=poporangeL,
  before upper={\popboxhead{Objectifs de la leçon}{popink}{}}}
\newtcolorbox{prerequis}{popbase,colframe=popink,colback=popblueL,
  before upper={\popboxhead{Prérequis}{popblue}{}}}
\newtcolorbox{bilan}{popbase,colframe=popink,colback=white,boxrule=2.8pt,
  before upper={\popboxhead{Fiche bilan}{poporange}{L'essentiel à connaître pour l'examen}}}
% Figure encadrée avec légende
\newtcolorbox{popfigure}[1][]{enhanced,breakable=false,colback=white,colframe=popline,boxrule=1.4pt,arc=8pt,
  left=10pt,right=10pt,top=10pt,bottom=8pt,before skip=10pt,after skip=12pt,halign=center,
  lower separated=false,halign lower=center,
  fontlower=\popsemi\fontsize{9}{11.5}\selectfont\color{popmuted},
  #1}
\newcommand{\legende}[1]{\par\vspace{6pt}{\popsemi\fontsize{9}{11.5}\selectfont\color{popmuted}#1}}

% ============================================================
% Signature OnBuch+
% ============================================================
\newcommand{\poplogoinline}[1]{{\poplogo\fontsize{#1}{#1}\selectfont\color{popink}OnBuch\color{poporange}+}}
\newcommand{\signatureonbuch}{%
  \begin{tcolorbox}[enhanced,colback=popink,colframe=popink,boxrule=2.2pt,arc=10pt,
      drop shadow=poporange,left=14pt,right=14pt,top=10pt,bottom=10pt,before skip=0pt,after skip=0pt]
    \tikz[baseline=-0.7ex]\node[circle,fill=popgreen,draw=popcream,line width=1.3pt,minimum size=22pt,inner sep=0]{\popcheck{white}};%
    \hspace{9pt}\begin{minipage}[c]{0.62\linewidth}
      {\popxbold\fontsize{12}{13}\selectfont\color{popcream}Signé\ \textbullet\ {\poplogo OnBuch\color{poporange}+}}\par\vspace{2pt}
      {\raggedright\popsemi\fontsize{8}{10.5}\selectfont\color{popline}\MakeUppercase{Équipe pédagogique OnBuch+}\\\MakeUppercase{Conforme au programme officiel MINESEC}\par}
    \end{minipage}\hfill
    {\poplogo\fontsize{22}{22}\selectfont\color{popcream}OnBuch\color{poporange}+}
  \end{tcolorbox}%
}

% ============================================================
% Page de garde
% #1 titre · #2 matière · #3 niveau · #4 module · #5 n° leçon · #6 durée ·
% #7 description / objectifs (texte ou itemize)
% ============================================================
\newcommand{\pagegarde}[7]{%
  \thispagestyle{empty}
  \newgeometry{a4paper,margin=1.7cm,top=1.6cm,bottom=1.4cm}%
  \begin{tikzpicture}[remember picture,overlay]
    \fill[poporange!18] ([xshift=-3.2cm,yshift=-3.6cm]current page.north east) circle (4.6cm);
    \fill[poppurple!14] ([xshift=2.4cm,yshift=7.6cm]current page.south west) circle (3.8cm);
  \end{tikzpicture}%
  {\poplogo\fontsize{30}{30}\selectfont\color{popink}OnBuch\color{poporange}+}\hfill
  \raisebox{0.6ex}{\poppill{Cours signé \textbullet\ Premium}{popink}{popcream}}\par
  \vspace{1pt}\popsquiggle{poppurple}\par
  \vspace{8pt}
  \begin{tcolorbox}[enhanced,colback=poporange,colframe=popink,boxrule=2.8pt,arc=13pt,
      drop shadow=popink,left=18pt,right=18pt,top=12pt,bottom=13pt,after skip=10pt]
    \poppill{#2 \textbullet\ #3}{popink}{popcream}\hspace{5pt}\poppill{#4}{white}{popink}\par\vspace{8pt}
    {\popdisplay\fontsize{14}{14}\selectfont\color{popink}LEÇON #5}\par\vspace{4pt}
    {\raggedright\hyphenpenalty=10000\exhyphenpenalty=10000\popdisplay\fontsize{25}{27}\selectfont\color{popcream}\MakeUppercase{#1}\par}
    \vspace{8pt}
    {\popxbold\fontsize{9.5}{9.5}\selectfont\color{popink}\MakeUppercase{Durée conseillée : #6}}
  \end{tcolorbox}
  \begin{tcolorbox}[enhanced,colback=white,colframe=popink,boxrule=2.2pt,arc=10pt,
      drop shadow=popink,left=16pt,right=16pt,top=10pt,bottom=10pt,after skip=8pt]
    \poppill{Dans cette leçon}{poppurple}{white}\par\vspace{5pt}
    \popsemi\fontsize{9.3}{12.6}\selectfont\color{popink}#7
  \end{tcolorbox}
  \noindent\begin{minipage}[t]{0.487\linewidth}
    \begin{tcolorbox}[enhanced,colback=popgreen,colframe=popink,boxrule=2.2pt,arc=10pt,
        drop shadow=popink,left=11pt,right=11pt,top=10pt,bottom=10pt,height=3.1cm,valign=center]
      \tcbox[colback=white,colframe=popink,boxrule=1.3pt,arc=5pt,boxsep=0pt,nobeforeafter,
        left=3pt,right=3pt,top=3pt,bottom=3pt,valign=center]{\qrcode[height=1.9cm]{https://play.google.com/store/apps/details?id=com.luvvix.onbuch&hl=fr}}%
      \hspace{8pt}\begin{minipage}[c]{0.5\linewidth}
        {\popdisplay\fontsize{11}{12.5}\selectfont\color{white}\MakeUppercase{Télécharge OnBuch}}\par\vspace{4pt}
        {\raggedright\popsemi\fontsize{8.4}{11}\selectfont\color{white}Gratuit \textbullet\ Google Play\\Tuteur IA, annales, concours, résultats\par}
      \end{minipage}
    \end{tcolorbox}
  \end{minipage}\hfill
  \begin{minipage}[t]{0.487\linewidth}
    \begin{tcolorbox}[enhanced,colback=poppurple,colframe=popink,boxrule=2.2pt,arc=10pt,
        drop shadow=popink,left=11pt,right=11pt,top=10pt,bottom=10pt,height=3.1cm,valign=center]
      \tikz[baseline=-0.7ex]\node[circle,fill=white,draw=popink,line width=1.3pt,minimum size=1.6cm,inner sep=0]{\popdisplay\fontsize{24}{24}\selectfont\color{poppurple}+};%
      \hspace{8pt}\begin{minipage}[c]{0.6\linewidth}
        {\popdisplay\fontsize{11}{12.5}\selectfont\color{white}\MakeUppercase{Passe à OnBuch+}}\par\vspace{4pt}
        {\raggedright\popsemi\fontsize{8.4}{11}\selectfont\color{white}Tous les cours signés, Tuteur IA illimité, fiches PDF — onglet OnBuch+ de l'app\par}
      \end{minipage}
    \end{tcolorbox}
  \end{minipage}
  \vspace{8pt}
  \popcontenu
  \vfill
  \signatureonbuch
  \restoregeometry
  \newpage
}

% Rangée « Ce cours contient » (page de garde)
\newcommand{\poptile}[3]{% #1 couleur #2 gros texte #3 légende
  \begin{tcolorbox}[enhanced,colback=#1,colframe=popink,boxrule=2pt,arc=9pt,shadow={2.5pt}{-2.5pt}{0pt}{popink},
      left=6pt,right=6pt,top=9pt,bottom=9pt,halign=center,valign=center,height=1.75cm,width=\linewidth,nobeforeafter]
    {\popdisplay\fontsize{12.5}{13}\selectfont\color{white}\MakeUppercase{#2}}\par\vspace{3pt}
    {\centering\popsemi\fontsize{7.8}{9.5}\selectfont\color{white}#3\par}
  \end{tcolorbox}}
\newcommand{\popcontenu}{%
  \par\noindent{\popxboldls\fontsize{7.5}{7.5}\selectfont\color{popmuted}CE COURS SIGNÉ CONTIENT}\par\vspace{7pt}
  \noindent\begin{minipage}[t]{0.235\linewidth}\poptile{popblue}{Cours}{complet, illustré, pas à pas}\end{minipage}\hfill
  \begin{minipage}[t]{0.235\linewidth}\poptile{poporange}{Méthodes}{et exercices résolus}\end{minipage}\hfill
  \begin{minipage}[t]{0.235\linewidth}\poptile{poppink}{Exercices}{type examen + corrigés}\end{minipage}\hfill
  \begin{minipage}[t]{0.235\linewidth}\poptile{popgreen}{Fiche bilan}{l'essentiel à retenir}\end{minipage}\par}

% Sommaire
\newtcolorbox{sommaire}{enhanced,colback=white,colframe=popink,boxrule=2.2pt,arc=10pt,
  drop shadow=popink,left=18pt,right=18pt,top=13pt,bottom=13pt,before skip=6pt,after skip=20pt,
  before upper={\poppill{Sommaire}{poppurple}{white}\par\vspace{9pt}\popsemi\fontsize{10.6}{14.5}\selectfont\color{popink}}}

% Signature de fin de cours — poussée tout en bas de la dernière page
\newcommand{\findecours}{%
  \par\vfill\nopagebreak
  \begin{center}\popsquiggle{poporange}\end{center}\vspace{4pt}
  \signatureonbuch
}
% Glyphes absents de Fira Math : redéfinis à partir de glyphes disponibles
\AtBeginDocument{%
  \def\setminus{\mathbin{\backslash}}\let\smallsetminus\setminus
  \def\vdots{\mathord{\vbox{\baselineskip3.2pt\lineskiplimit0pt\kern4pt\hbox{.}\hbox{.}\hbox{.}}}}%
  \def\ddots{\mathinner{\mkern1mu\raise6.5pt\hbox{.}\mkern2mu\raise3.6pt\hbox{.}\mkern2mu\raise0.7pt\hbox{.}\mkern1mu}}%
  \def\triangle{\mathord{\scalebox{0.9}{$\symup{\Delta}$}}}%
  \def\nabla{\mathord{\raisebox{\depth}{\scalebox{1}[-1]{$\symup{\Delta}$}}}}%
  \def\checkmark{\popcheck{popdark}}%
  \def\star{\mathbin{\ast}}\def\bigstar{\mathord{\ast}}%
  \def\diamond{\mathbin{\scalebox{0.6}{$\Diamond$}}}%
  \def\bigcup{\mathop{\mathchoice{\raisebox{-0.3em}{\scalebox{1.7}{$\cup$}}}{\scalebox{1.25}{$\cup$}}{\cup}{\cup}}\displaylimits}%
  \def\bigcap{\mathop{\mathchoice{\raisebox{-0.3em}{\scalebox{1.7}{$\cap$}}}{\scalebox{1.25}{$\cap$}}{\cap}{\cap}}\displaylimits}%
  \def\lesseqqgtr{\mathrel{\lessgtr}}\def\bigoplus{\mathop{\oplus}\displaylimits}%
}
\providecommand{\ssi}{si et seulement si} % abréviation fréquente

%% FIGFIX-BEGIN v5 — correctifs globaux des figures (docs/onbuch-generation/tools/figfix)
\usepackage{adjustbox}\usepackage{etoolbox}\tcbuselibrary{raster}
% toute figure TikZ/pgfplots plus large que la ligne est réduite à la largeur disponible
\BeforeBeginEnvironment{tikzpicture}{\begin{adjustbox}{max width=\linewidth}}
\AfterEndEnvironment{tikzpicture}{\end{adjustbox}}
% légendes pgfplots lisibles par-dessus les courbes
\pgfplotsset{every axis/.append style={legend style={fill=white,fill opacity=0.92,text opacity=1,draw=black!15,inner sep=3pt}}}
% étiquettes d'arêtes (syntaxe edge["..."]) avec fond blanc
\tikzset{every edge quotes/.append style={fill=white,inner sep=1.2pt}}
%% FIGFIX-END
\begin{document}

\pagegarde{Acides et bases en solution aqueuse : eau, autoprotolyse et pH}{Chimie}{Tle TI}{Module 2 · Acides et bases}{06}{6 h (cours + TP)}{Cette leçon étudie l'eau comme solvant et comme ampholyte : autoprotolyse, produit ionique Ke et définition du pH. Elle apprend à calculer le pH et les concentrations ioniques des solutions aqueuses, à prévoir leur caractère acide, neutre ou basique, et à mesurer le pH au laboratoire.
\vspace{4pt}
\begin{multicols}{2}\small
\begin{itemize}
\item À la fin de cette leçon, je sais expliquer la structure polaire de la molécule d'eau et son pouvoir solvant et dissociant.
\item À la fin de cette leçon, je sais reconnaître un acide et une base selon Bronsted et identifier les couples acide/base.
\item À la fin de cette leçon, je sais écrire l'équation-bilan de l'autoprotolyse de l'eau et définir le produit ionique Ke.
\item À la fin de cette leçon, je sais définir le pH et le calculer à partir de la concentration en ions H3O+.
\item À la fin de cette leçon, je sais calculer les concentrations [H3O+] et [HO-] d'une solution connaissant son pH.
\item À la fin de cette leçon, je sais déduire le caractère acide, neutre ou basique d'une solution à partir de son pH à une température donnée.
\end{itemize}\end{multicols}}

\begin{sommaire}
\begin{multicols}{2}
\begin{enumerate}
\item L'eau : structure et pouvoir solvant
\item Acides et bases selon Bronsted
\item Autoprotolyse de l'eau et produit ionique Ke
\item Le pH : définition et mesure
\item Calculs de pH et concentrations ioniques
\item Préparation de solutions par dilution
\item Travaux pratiques
\item Méthodes et exercices résolus
\item Exercices
\item Corrigés des exercices
\item Fiche bilan
\end{enumerate}
\end{multicols}
\end{sommaire}

\begin{objectifs}
\begin{itemize}
\item À la fin de cette leçon, je sais expliquer la structure polaire de la molécule d'eau et son pouvoir solvant et dissociant.
\item À la fin de cette leçon, je sais reconnaître un acide et une base selon Bronsted et identifier les couples acide/base.
\item À la fin de cette leçon, je sais écrire l'équation-bilan de l'autoprotolyse de l'eau et définir le produit ionique Ke.
\item À la fin de cette leçon, je sais définir le pH et le calculer à partir de la concentration en ions H3O+.
\item À la fin de cette leçon, je sais calculer les concentrations [H3O+] et [HO-] d'une solution connaissant son pH.
\item À la fin de cette leçon, je sais déduire le caractère acide, neutre ou basique d'une solution à partir de son pH à une température donnée.
\item À la fin de cette leçon, je sais effectuer les approximations (espèces majoritaires, minoritaires, ultra-minoritaires) pour simplifier un calcul de pH.
\item À la fin de cette leçon, je sais préparer une solution diluée à partir d'une solution concentrée (calcul et protocole).
\item À la fin de cette leçon, je sais mesurer un pH avec le papier pH et utiliser les indicateurs colorés.
\end{itemize}
\end{objectifs}

\begin{prerequis}
\begin{itemize}
\item Notion d'ion, de molécule et de liaison covalente (Seconde / Première)
\item Concentration molaire, quantité de matière et dilution (Première)
\item Écriture et équilibrage d'une équation-bilan de réaction chimique
\item Fonction logarithme décimal et puissances de 10 (mathématiques de Première)
\item Réaction acido-basique vue en introduction du Module 2 (transfert de proton H+)
\end{itemize}
\end{prerequis}

\newpage

\coursec{L'eau : structure et pouvoir solvant}

Amina, élève de Terminale TI à Yaoundé, accompagne sa mère au centre de santé : le bilan sanguin indique un pH de 7,38, jugé « normal » par le médecin. De retour à la maison, elle mesure avec du papier pH le jus de citron (pH ≈ 2), l'eau de forage du quartier (pH ≈ 6,5) et la lessive de savon (pH ≈ 10). Pourquoi l'eau pure, qui semble « sans rien », contient-elle pourtant des ions \ce{H3O+} et \ce{HO-} ? Comment un simple nombre, le pH, permet-il de comparer l'acidité du sang, d'un jus de fruit ou d'un produit ménager ? Cette leçon donne les clés pour lire et interpréter ces valeurs du quotidien et des analyses médicales.

\courssub{Géométrie et polarité de la molécule d'eau}

La molécule d'eau, \ce{H2O}, est le socle de toute chimie en solution aqueuse. Sa géométrie n'est pas linéaire : les deux atomes d'hydrogène forment un angle au voisinage de l'atome d'oxygène.

\begin{popfigure}
\begin{tikzpicture}[scale=1.6, every node/.style={font=\small}]
% Atomes
\node[draw, circle, fill=popblue!80, text=white, minimum size=1.4cm, inner sep=1pt] (O) at (0,0) {\Large O};
\node[draw, circle, fill=poporange!80, text=white, minimum size=0.8cm, inner sep=1pt] (H1) at (1.2,0.85) {\large H};
\node[draw, circle, fill=poporange!80, text=white, minimum size=0.8cm, inner sep=1pt] (H2) at (1.2,-0.85) {\large H};

% Liaisons
\draw[thick, popdark] (O) -- (H1) node[midway, above left, sloped] {$\delta^+$};
\draw[thick, popdark] (O) -- (H2) node[midway, below left, sloped] {$\delta^+$};

% Doublets non liants (représentation schématique)
\draw[thick, poppurple, dashed] (O) ++(135:0.7) arc (135:165:0.7) node[midway, above left, poppurple] {$e^-$};
\draw[thick, poppurple, dashed] (O) ++(225:0.7) arc (225:195:0.7) node[midway, below left, poppurple] {$e^-$};

% Dipôles de liaison
\draw[->, thick, popblue] (O) -- ++(-1.3,-0.9) node[left] {$\vec{\mu}_{\ce{O-H}}$};
\draw[->, thick, popblue] (O) -- ++(-1.3,0.9) node[left] {$\vec{\mu}_{\ce{O-H}}$};

% Dipôle moléculaire résultant
\draw[->, very thick, popred] (0,-1.8) -- (0,-2.8) node[below] {$\vec{\mu}_{\text{mol}}$};

% Angle
\draw[<->, thick] (0.7,0.45) arc (33:-33:0.7) node[midway, right=3mm] {$104,5^\circ$};

% Charges partielles
\node[popblue, font=\small] at (-1.5, 0) {$\delta^-$};
\end{tikzpicture}
\legende{Géométrie coudée de la molécule d'eau : angle \ce{H-O-H} ≈ 104,5°, charges partielles $\delta^-$ sur O et $\delta^+$ sur H, dipôles de liaison et dipôle moléculaire résultant.}
\end{popfigure}

\begin{exemplebox}[Données caractéristiques de la molécule d'eau]
\textbf{Angle \ce{H-O-H}} : $104,5^\circ$ (molécule coudée).\\
\textbf{Longueur de liaison \ce{O-H}} : $\approx \SI{96}{pm}$ ($0,96\,\text{\AA}$).\\
\textbf{Moment dipolaire} : $\mu = \SI{1,85}{D}$ (Debye).\\
\textbf{Électronégativités (échelle Pauling)} : $\chi_{\ce{O}} = 3,44$ ; $\chi_{\ce{H}} = 2,20$. $\Delta\chi = 1,24$ → liaison polarisée covalente.
\end{exemplebox}

\textbf{Pourquoi cet angle ?}\par
L'atome d'oxygène possède 6 électrons de valence. Il forme deux liaisons covalentes \ce{O-H} et conserve deux doublets non liants. Selon la méthode VSEPR, les quatre domaines électroniques (2 liaisons + 2 doublets) s'orientent vers les sommets d'un tétraèdre. Les doublets non liants, plus répulsifs, « compressent » l'angle de liaison : l'angle tétraédrique idéal ($109,5^\circ$) diminue à $104,5^\circ$.

\textbf{Polarité des liaisons et de la molécule.}\par
L'oxygène, plus électronégatif, attire le nuage électronique des liaisons \ce{O-H} vers lui : chaque liaison porte un dipôle $\vec{\mu}_{\ce{O-H}}$ orienté de H vers O. Comme la molécule est coudée, la somme vectorielle des deux dipôles de liaison n'est pas nulle : \textbf{la molécule d'eau est polaire}, elle possède un dipôle permanent $\vec{\mu}_{\text{mol}}$ orienté de l'atome d'hydrogène vers l'atome d'oxygène (du centre des charges positives vers le centre des charges négatives).

\begin{attention}[Piège classique : la molécule d'eau linéaire]
Ne dessine \textbf{jamais} la molécule d'eau sous forme linéaire (\ce{H-O-H} à $180^\circ$). Si elle était linéaire, les deux dipôles de liaison s'annuleraient : $\vec{\mu}_{\text{mol}} = \vec{0}$. L'eau serait apolaire, n'aurait pas de liaisons hydrogène, bouillerait vers $-\SI{80}{\celsius}$ et la vie telle que nous la connaissons serait impossible. \textbf{L'angle coudé est la clé de tout.}
\end{attention}

\courssub{Liaisons hydrogène : la « colle » de l'eau}

La polarité de l'eau et la petite taille de l'atome d'hydrogène (dépourvu de coquille électronique interne quand il est lié à O) permettent la formation de \textbf{liaisons hydrogène} entre molécules d'eau voisines.

\begin{definition}[Liaison hydrogène]
Une liaison hydrogène est une interaction électrostatique attractive entre un atome d'hydrogène \textbf{lié à un atome très électronégatif} (O, N, F) portant une charge partielle $\delta^+$ et un \textbf{doublet non liant} d'un atome électronégatif (O, N, F) portant une charge partielle $\delta^-$ sur une molécule voisine. Elle est notée par des pointillés ($\cdots$).
\end{definition}

Dans l'eau liquide, chaque molécule \ce{H2O} peut former jusqu'à \textbf{quatre liaisons hydrogène} : deux comme donneur (via ses atomes H) et deux comme accepteur (via les doublets de l'O). Ce réseau tridimensionnel dynamique confère à l'eau des propriétés exceptionnelles.

\begin{propriete}[Conséquences macroscopiques des liaisons hydrogène]
\begin{itemize}
    \item \textbf{Températures de changement d'état élevées :} $T_{\text{éb}} = \SI{100}{\celsius}$ (au lieu de $\approx -\SI{80}{\celsius}$ sans liaisons H), $T_{\text{fus}} = \SI{0}{\celsius}$.
    \item \textbf{Chaleur massique élevée :} $c_{\text{eau}} = \SI{4,18}{J.g^{-1}.K^{-1}}$ → l'eau tamponne les variations de température (climat, organismes vivants).
    \item \textbf{Anomalie de la densité :} maximum de densité à $\SI{4}{\celsius}$ → la glace flotte, isolant l'eau liquide en dessous (survie de la faune aquatique au Lac Nyos ou dans les lacs de montagne en saison sèche).
    \item \textbf{Tension de surface élevée :} déplacement de l'eau dans les xylems des plantes (cacao, bananier), formation de gouttes sphériques.
\end{itemize}
\end{propriete}

\begin{savaistu}[L'eau, cette anomalie qui fait la vie]
Sans liaisons hydrogène, l'eau bouillerait à environ $-\SI{80}{\celsius}$ (extrapolation depuis \ce{H2S}, \ce{H2Se}, \ce{H2Te}). À température ambiante, ce serait un gaz ! Les océans, le sang, la sève, le vin de palme, le bili-bili n'existeraient pas sous forme liquide. Le lac Nyos, au Nord-Ouest du Cameroun, garde sa vie aquatique grâce à la stratification thermique rendue possible par l'anomalie de densité de l'eau à $\SI{4}{\celsius}$. Les liaisons hydrogène ne sont pas un détail : elles sont la condition \emph{sine qua non} de la vie sur Terre.
\end{savaistu}

\courssub{Pouvoir solvant de l'eau : dissolution des solides ioniques et composés polaires}

L'eau est souvent appelée « solvant universel » (à tort, car elle ne dissout pas tout). C'est le \textbf{solvant polaire} par excellence.

\begin{definition}[Solvant polaire et solvatation (hydratation)]
Un \textbf{solvant polaire} est un solvant dont les molécules possèdent un moment dipolaire permanent non nul (ex. : \ce{H2O}, \ce{CH3OH}, \ce{NH3}).
La \textbf{solvatation} est l'entourage des particules du soluté (ions ou molécules) par les molécules du solvant. Quand le solvant est l'eau, on parle d'\textbf{hydratation}. Les ions en solution aqueuse sont notés avec l'indice \ce{(aq)} (ex. \ce{Na+(aq)}, \ce{Cl-(aq)}).
\end{definition}

\textbf{La règle de similitude : « Les semblables dissolvent les semblables ».}\par
\begin{itemize}
    \item \textbf{Composés ioniques (sels) :} L'eau dissout la plupart des sels (NaCl, \ce{KNO3}, \ce{CuSO4}). L'énergie libérée par l'hydratation des ions compense l'énergie de réseau du cristal.
    \item \textbf{Composés polaires covalents :} L'eau dissout l'éthanol (\ce{CH3CH2OH}), le glucose (\ce{C6H12O6}), l'acide acétique (\ce{CH3COOH}), l'urée, le sucre de canne. Les groupes \ce{-OH}, \ce{-COOH}, \ce{-NH2} forment des liaisons hydrogène avec l'eau.
    \item \textbf{Composés apolaires :} L'eau \textbf{ne dissout pas} (ou très peu) l'huile de palme, l'essence, le hexane, le benzène, le dioxygène \ce{O2}, le diazote \ce{N2}. Il n'y a pas d'interactions électrostatiques favorables entre molécules d'eau polaires et molécules apolaires.
\end{itemize}

\begin{popfigure}
\begin{tikzpicture}[scale=0.85, every node/.style={font=\small}]
% Étape 1 : Cristal ionique
\begin{scope}[xshift=0cm]
\node[draw, rounded corners, fill=popblueL, inner sep=4pt] at (0, 3.3) {\textbf{Étape 1 : Cristal ionique solide}};
\foreach \x in {0,1,2} {
  \foreach \y in {0,1,2} {
    \pgfmathsetmacro{\col}{mod(\x+\y,2)==0 ? "poporange!80" : "popblue!80"};
    \pgfmathsetmacro{\lbl}{mod(\x+\y,2)==0 ? "Na$^+$" : "Cl$^-$"};
    \node[draw, circle, fill=\col, text=white, minimum size=0.6cm, inner sep=0] at (\x*0.9, \y*0.9) {\lbl};
  }
}
\draw[->, thick, popdark] (3.5, 1.35) -- (4.5, 1.35) node[midway, above] {Eau};
\end{scope}

% Étape 2 : Solvatation / Rupture du réseau
\begin{scope}[xshift=6.5cm]
\node[draw, rounded corners, fill=poporangeL, inner sep=4pt] at (0, 3.3) {\textbf{Étape 2 : Solvatation (Hydratation)}};
% Ions au centre
\node[draw, circle, fill=poporange!80, text=white, minimum size=0.7cm] (Na) at (-1, 1.35) {Na$^+$};
\node[draw, circle, fill=popblue!80, text=white, minimum size=0.7cm] (Cl) at (2, 1.35) {Cl$^-$};
% Molécules d'eau autour de Na+
\foreach \angle/\pos in {45/0.65, 135/0.65, 225/0.65, 315/0.65, 90/1.05, 270/1.05} {
  \node[draw, circle, fill=popblue!30, minimum size=0.5cm] at ($(-1,1.35)+(\angle:\pos)$) {H$_2$O};
  \draw[->, thin, popblue] ($(-1,1.35)+(\angle:\pos)$) -- ($(-1,1.35)+(\angle:0.35)$);
}
% Molécules d'eau autour de Cl-
\foreach \angle/\pos in {45/0.65, 135/0.65, 225/0.65, 315/0.65, 90/1.05, 270/1.05} {
  \node[draw, circle, fill=popblue!30, minimum size=0.5cm] at ($(2,1.35)+(\angle:\pos)$) {H$_2$O};
  \draw[->, thin, poporange] ($(2,1.35)+(\angle:0.35)$) -- ($(2,1.35)+(\angle:\pos)$);
}
\end{scope}

% Étape 3 : Solution aqueuse (Dispersion)
\begin{scope}[xshift=14cm]
\node[draw, rounded corners, fill=popgreenL, inner sep=4pt] at (0, 3.3) {\textbf{Étape 3 : Solution aqueuse (Dispersion)}};
\foreach \i in {1,...,8} {
  \pgfmathsetmacro{\x}{rand*3};
  \pgfmathsetmacro{\y}{rand*1.8+0.3};
  \pgfmathsetmacro{\col}{mod(\i,2)==0 ? "poporange!80" : "popblue!80"};
  \pgfmathsetmacro{\lbl}{mod(\i,2)==0 ? "Na$^+$" : "Cl$^-$"};
  \node[draw, circle, fill=\col, text=white, minimum size=0.5cm] at (\x, \y) {\lbl};
  % Quelques eaux autour
  \foreach \a in {0, 90, 180, 270} {
    \node[circle, fill=popblue!20, minimum size=0.3cm] at ($(\x,\y)+(\a:0.4)$) {\tiny H$_2$O};
  }
}
\end{scope}
\end{tikzpicture}
\legende{Mécanisme de dissolution du chlorure de sodium \ce{NaCl} dans l'eau : (1) Réseau cristallin ordonné ; (2) Attaque par les dipôles de l'eau : orientation des molécules d'eau (O vers Na$^+$, H vers Cl$^-$), rupture du réseau ; (3) Ions hydratés dispersés uniformément dans le volume.}
\end{popfigure}

\textbf{Mécanisme microscopique de la dissolution d'un cristal ionique (ex. \ce{NaCl}).}\par
\begin{enumerate}
    \item \textbf{Approche :} Les molécules d'eau polaires s'orientent à la surface du cristal : l'atome O ($\delta^-$) pointe vers les cations \ce{Na+}, les atomes H ($\delta^+$) pointent vers les anions \ce{Cl-}.
    \item \textbf{Solvatation (Hydratation) :} Les interactions ion-dipôle entre \ce{Na+}/\ce{Cl-} et \ce{H2O} sont assez fortes pour arracher les ions du réseau cristallin (énergie de réseau). Chaque ion se trouve au centre d'une \textbf{sphère d'hydratation} (généralement 4 à 6 molécules d'eau pour \ce{Na+}, 6 pour \ce{Cl-}).
    \item \textbf{Dispersion :} Les ions hydratés, stabilisés, diffusent dans le volume de la solution par agitation thermique. La solution est homogène, transparente, conductrice de l'électricité.
\end{enumerate}

\begin{aretenir}[Équation de dissolution du chlorure de sodium]
L'équation-bilan de la dissolution (changement d'état, pas de réaction chimique, les ions existent déjà dans le cristal) s'écrit :
\[ \ce{NaCl(s) ->[H2O] Na+(aq) + Cl-(aq)} \]
\begin{itemize}
    \item \ce{(s)} : état solide (cristal).
    \item \ce{(aq)} : état aqueux (ion hydraté, solvaté par l'eau).
    \item La flèche \ce{->} indique que la dissolution est souvent totale (sel soluble) à concentration raisonnable. Pour un sel peu soluble, on utiliserait \ce{<=>}.
\end{itemize}
\end{aretenir}

\textbf{Pouvoir dissociant et constante diélectrique.}\par
La capacité de l'eau à séparer les ions d'un cristal (à « dissocier » le réseau) provient de sa \textbf{constante diélectrique relative} (permittivité) très élevée : $\varepsilon_r \approx 80$ (à $\SI{25}{\celsius}$).
La force électrostatique de Coulomb entre deux charges $q_1$ et $q_2$ à distance $r$ dans un milieu de permittivité $\varepsilon = \varepsilon_0 \varepsilon_r$ est :
\[ F = \frac{q_1 q_2}{4\pi\varepsilon_0\varepsilon_r r^2} \]
Dans l'eau ($\varepsilon_r \approx 80$), la force d'attraction entre \ce{Na+} et \ce{Cl-} est \textbf{80 fois plus faible} que dans le vide (ou l'air). Cela facilite énormément la séparation des ions par l'agitation thermique et la solvatation.

\courssub{Limites du pouvoir solvant : les composés apolaires}

\begin{experience}[Miscibilité eau / huile de palme (ou huile de cuisine)]
\textbf{Matériel :} Tube à essai, eau, huile de palme rouge (ou huile d'arachide), colorant alimentaire (facultatif).
\textbf{Protocole :}
\begin{enumerate}
    \item Verser \SI{5}{mL} d'eau dans le tube.
    \item Ajouter \SI{5}{mL} d'huile de palme.
    \item Boucher, agiter vigoureusement pendant 30 s.
    \item Laisser reposer \SI{2}{min}.
\end{enumerate}
\textbf{Observations :} Lors de l'agitation, une émulsion trouble se forme (gouttelettes d'huile dans l'eau). Au repos, deux phases distinctes se reforment rapidement : l'huile (moins dense, $\rho \approx 0,92$) remonte en surface, l'eau reste en bas. La limite de séparation est nette.
\textbf{Interprétation :} L'eau (polaire, liaisons H) et l'huile (apolaire, liaisons de Van der Waals uniquement) n'ont pas d'affinité intermoléculaire. L'énergie nécessaire pour créer une cavité dans le réseau d'eau pour y loger une molécule d'huile n'est pas compensée par des interactions eau-huile. \textbf{L'eau ne solvate pas les molécules apolaires.}
\end{experience}

\begin{attention}[Erreur fréquente : « L'eau dissout tout »]
Faux. L'eau ne dissout \textbf{pas} les corps gras (huiles, graisses, cires, essence, gazole, polymères comme le polyéthylène). Pour les dissoudre, il faut un solvant apolaire (hexane, éther, tétrachlorure de carbone \ce{CCl4} — toxique) ou un tensioactif (savon, détergent) qui forme des micelles. C'est le principe du lavage du linge avec du savon de ménage (fabriqué artisanalement à partir d'huile de palme et de soude au Cameroun) : le savon a une tête polaire (aimant l'eau) et une queue apolaire (aimant la graisse).
\end{attention}

\courssub{Exemple résolu : Préparation d'une solution aqueuse de glucose}

\begin{exoresolu}[Préparation de \SI{250,0}{mL} de solution de glucose à \SI{50,0}{g.L^{-1}}]
\textbf{Énoncé :} On souhaite préparer \SI{250,0}{mL} d'une solution aqueuse de glucose (\ce{C6H12O6}, $M = \SI{180,16}{g.mol^{-1}}$) de concentration massique $C_m = \SI{50,0}{g.L^{-1}}$. Calculer la masse de glucose à peser et décrire la procédure opératoire au laboratoire.
\tcblower
\textbf{Données :} $V_{\text{sol}} = \SI{250,0}{mL} = \SI{0,2500}{L}$ ; $C_m = \SI{50,0}{g.L^{-1}}$.
\textbf{Calcul de la masse $m$ :}
$C_m = \frac{m}{V_{\text{sol}}} \implies m = C_m \times V_{\text{sol}} = \SI{50,0}{g.L^{-1}} \times \SI{0,2500}{L} = \SI{12,50}{g}$.
\textbf{Procédure opératoire (verrerie classe A) :}
\begin{enumerate}
    \item Peser exactement \SI{12,50}{g} de glucose solide dans une capsule de pesée (balance au centième ou millième de gramme).
    \item Transvaser le glucose dans un bécher propre (ex. \SI{100}{mL}) avec un peu d'eau distillée (\SI{\approx 50}{mL}). Mélanger à la baguette de verre jusqu'à dissolution complète (le glucose est polaire, très soluble).
    \item Transférer quantitativement la solution dans une fiole jaugée de \SI{250,0}{mL} (entonnoir, rinçage du bécher et de la baguette avec de l'eau distillée, verser les rinçages dans la fiole).
    \item Compléter à l'eau distillée jusqu'au trait de jauge (bas du ménisque au niveau du trait, œil au niveau du trait). Bouchonner, homogénéiser par retournements répétés.
\end{enumerate}
\textbf{Résultat :} On obtient \SI{250,0}{mL} de solution de glucose \SI{50,0}{g.L^{-1}} (soit $C \approx \SI{0,277}{mol.L^{-1}}$).
\end{exoresolu}

\courssub{Exercices d'application}

\begin{exercice}[Géométrie et polarité]
\begin{enumerate}
    \item Représenter la molécule d'eau \ce{H2O} en indiquant : l'angle \ce{H-O-H}, les charges partielles $\delta^+$ et $\delta^-$, le dipôle de chaque liaison \ce{O-H} et le dipôle moléculaire résultant.
    \item Expliquer pourquoi l'angle \ce{H-O-H} vaut $104,5^\circ$ et non $109,5^\circ$ (tétraèdre) ni $180^\circ$ (linéaire).
    \item Prédire si les molécules suivantes peuvent former des liaisons hydrogène \textbf{entre elles} : \ce{CH4}, \ce{NH3}, \ce{HF}, \ce{CH3OH}, \ce{CH3CH3}. Justifier.
\end{enumerate}
\end{exercice}

\begin{exercice}[Dissolution et solvatation]
\begin{enumerate}
    \item Écrire l'équation de dissolution du sulfate de cuivre(II) \ce{CuSO4} dans l'eau. Préciser l'état physique de chaque espèce.
    \item Décrire l'orientation des molécules d'eau autour de l'ion \ce{Cu^{2+}(aq)} et de l'ion \ce{SO4^{2-}(aq)}.
    \item On verse \SI{5}{mL} d'huile de palme dans \SI{50}{mL} d'eau. Après agitation et repos, on observe deux phases. Expliquer cette observation à l'échelle microscopique en utilisant la notion d'interactions intermoléculaires.
\end{enumerate}
\end{exercice}

\begin{corrige}[Exercice 1]
\begin{enumerate}
    \item Voir figure du cours. Dipôles \ce{O-H} orientés H $\to$ O. Dipôle moléculaire orienté du milieu des H vers O. Angle $104,5^\circ$.
    \item 4 domaines électroniques (2 liaisons, 2 doublets non liants) $\to$ géométrie électronique tétraédrique. Répulsion doublet non liant / doublet de liaison > répulsion liaison / liaison $\to$ angle compressé à $104,5^\circ$. Linéaire impliquerait 2 domaines (sp) ou absence de doublets non liants, ce qui n'est pas le cas de l'O.
    \item \ce{NH3} (OUI, H sur N + doublet sur N), \ce{HF} (OUI, H sur F + doublet sur F), \ce{CH3OH} (OUI, groupe \ce{-OH}). \ce{CH4} (NON, H sur C pas assez $\delta^+$, pas de doublet sur C), \ce{CH3CH3} (NON, pas d'atome H lié à O/N/F, pas de doublet sur C).
\end{enumerate}
\end{corrige}

\begin{corrige}[Exercice 2]
\begin{enumerate}
    \item $\ce{CuSO4(s) ->[H2O] Cu^{2+}(aq) + SO4^{2-}(aq)}$
    \item Autour de \ce{Cu^{2+}} : atomes O ($\delta^-$) des molécules d'eau pointent vers le cation (coordination souvent 4 ou 6, géométrie plane carrée ou octaédrique, couleur bleue caractéristique). Autour de \ce{SO4^{2-}} : atomes H ($\delta^+$) pointent vers les atomes O de l'anion.
    \item L'eau est polaire (liaisons H fortes entre molécules d'eau). L'huile de palme est apolaire (liaisons de Van der Waals). Pour mélanger l'huile à l'eau, il faudrait « casser » le réseau de liaisons H de l'eau pour y insérer des molécules d'huile. L'énergie gagnée par les nouvelles interactions eau-huile (très faibles, dipôle-induit) ne compense pas l'énergie perdue par la rupture des liaisons H eau-eau. Le système minimise son énergie en séparant les deux phases.
\end{enumerate}
\end{corrige}

\coursec{Acides et bases selon Brønsted}

Dans la section précédente, nous avons vu que l'eau est une molécule \cle{polaire}, coudée, dont les doublets non liants de l'oxygène et les atomes d'hydrogène partiellement chargés positivement expliquent son extraordinaire \cle{pouvoir solvant} : elle sépare les ions des cristaux ioniques et solvate les molécules polaires. Mais l'eau ne se contente pas de dissoudre : elle \emph{participe} aux réactions. Quand tu verses du vinaigre dans l'eau, ou quand tu dilues de l'ammoniaque pour laver les sols, il se passe un échange microscopique d'une particule minuscule mais essentielle : le \cle{proton} \ce{H+}. Cet échange est le cœur de toute la chimie acido-basique. C'est le chimiste danois Johannes \textbf{Brønsted} qui, en 1923, en a proposé la définition moderne que nous allons construire pas à pas.

\courssub{La définition de Brønsted : un échange de proton}

\textbf{Intuition d'abord.} Imagine un match de football : il y a celui qui \emph{passe} le ballon et celui qui le \emph{reçoit}. En chimie acido-basique, le ballon, c'est le proton \ce{H+}. Un acide est le passeur, une base est le receveur. Sans receveur, pas de passe possible : un acide ne peut manifester son acidité qu'en présence d'une base, et réciproquement.

\begin{definition}[Acide et base selon Brønsted]
\begin{itemize}
\item Un \cle{acide} est une espèce chimique (molécule ou ion) capable de \textbf{céder un ou plusieurs protons} \ce{H+}.
\item Une \cle{base} est une espèce chimique (molécule ou ion) capable de \textbf{capter un ou plusieurs protons} \ce{H+}.
\end{itemize}
Cette capacité s'exprime par la demi-équation formelle :
\[ \text{acide} \;=\; \text{base} \;+\; \ce{H+} \]
\end{definition}

\begin{attention}[Deux pièges à détruire immédiatement]
\begin{enumerate}
\item \textbf{Ne jamais inverser les rôles} : l'acide \emph{donne} \ce{H+}, la base \emph{accepte} \ce{H+}. Beaucoup d'élèves confondent au baccalauréat. Astuce mnémotechnique : \textbf{A}cide = \textbf{A}bandonne ; \textbf{B}ase = \textbf{B}oiseau qui attrape.
\item Le proton \ce{H+} \textbf{n'existe jamais seul en solution aqueuse}. Nu, c'est un ion extrêmement réactif (c'est un noyau d'hydrogène sans électron !). Il se fixe immédiatement sur une molécule d'eau grâce à un doublet non liant de l'oxygène, pour former l'ion \cle{oxonium} \ce{H3O+} :
\[ \ce{H+ + H2O -> H3O+} \]
Quand on écrit \ce{H+} dans une demi-équation, c'est une \emph{commodité d'écriture} ; en solution, c'est toujours \ce{H3O+} que l'on manipule.
\end{enumerate}
\end{attention}

\courssub{Couples acide/base conjugués}

\textbf{Intuition.} Quand un acide a cédé son proton, il ne disparaît pas : il devient une espèce qui, elle, \emph{pourrait capter} un proton. Il s'est transformé en sa base « partenaire ». Les deux espèces forment un couple inséparable, comme les deux faces d'une pièce.

\begin{definition}[Couple acide/base conjugués]
Deux espèces chimiques forment un \cle{couple acide/base} (ou couple conjugué) si elles s'échangent par transfert d'un proton :
\[ \text{acide} \;=\; \text{base} \;+\; \ce{H+} \]
On note le couple \textbf{acide/base}. La base est dite \emph{base conjuguée} de l'acide, et l'acide est l'\emph{acide conjugué} de la base. La demi-équation s'écrit toujours avec l'\textbf{acide à gauche}.
\end{definition}

\textbf{Exemples fondamentaux à connaître par cœur :}

\begin{itemize}
\item Couple de l'ion oxonium : \ce{H3O+}/\ce{H2O} \quad demi-équation : $\ce{H3O+} = \ce{H2O} + \ce{H+}$
\item Couple de l'ion hydroxyde : \ce{H2O}/\ce{HO-} \quad demi-équation : $\ce{H2O} = \ce{HO-} + \ce{H+}$
\item Couple de l'acide éthanoïque (le vinaigre !) : \ce{CH3COOH}/\ce{CH3COO-} \quad $\ce{CH3COOH} = \ce{CH3COO-} + \ce{H+}$
\item Couple de l'ion ammonium : \ce{NH4+}/\ce{NH3} \quad $\ce{NH4+} = \ce{NH3} + \ce{H+}$
\item Couple du chlorure d'hydrogène : \ce{HCl}/\ce{Cl-} \quad $\ce{HCl} = \ce{Cl-} + \ce{H+}$
\end{itemize}

\begin{attention}[Vérifier un couple en 5 secondes]
Pour vérifier que deux espèces forment bien un couple, compte les atomes et les charges : on doit passer de l'une à l'autre en ajoutant ou retirant \emph{exactement un} \ce{H+}. Exemple : \ce{CH3COOH} et \ce{CH3COO-} diffèrent d'un H et d'une charge $+1$ : c'est un couple. En revanche \ce{H2SO4} et \ce{SO4^2-} ne forment \textbf{pas} un couple direct : ils diffèrent de \emph{deux} protons ; il faut passer par l'intermédiaire \ce{HSO4-}.
\end{attention}

\courssub{La réaction acido-basique : un transfert de proton}

\textbf{Intuition.} Revenons au match : la passe n'existe que s'il y a un passeur \emph{et} un receveur. De même, une réaction acido-basique met en jeu \textbf{deux couples} : l'acide du couple 1 cède son proton à la base du couple 2.

\begin{definition}[Réaction acido-basique]
Une \cle{réaction acido-basique} est un transfert de proton \ce{H+} de l'acide d'un couple vers la base d'un autre couple :
\[ \text{acide}_1 + \text{base}_2 \;\ce{->}\; \text{base}_1 + \text{acide}_2 \]
\end{definition}

\begin{methode}[Écrire l'équation-bilan d'une réaction acido-basique en 3 étapes]
\begin{enumerate}
\item \textbf{Identifier les deux couples} mis en jeu et repérer qui joue le rôle d'acide (celui qui possède le H à céder) et qui joue le rôle de base. Écrire les couples sous la forme acide/base.
\item \textbf{Écrire les deux demi-équations} : celle de l'acide qui cède \ce{H+} (acide $\to$ base + \ce{H+}) et celle de la base qui capte \ce{H+} (base + \ce{H+} $\to$ acide).
\item \textbf{Additionner} les deux demi-équations et \textbf{simplifier} : le proton \ce{H+} ne doit \emph{jamais} apparaître dans l'équation-bilan finale. Vérifier la conservation des atomes et des charges.
\end{enumerate}
\end{methode}

\begin{exemplebox}[L'acide du vinaigre réagit avec l'eau]
Le vinaigre consommé à Douala ou à Bafoussam est une solution d'\textbf{acide éthanoïque} \ce{CH3COOH} (environ \SI{1}{mol.L^{-1}} pour un vinaigre à 6 degrés). Appliquons la méthode :
\begin{enumerate}
\item Couples : \ce{CH3COOH}/\ce{CH3COO-} (l'acide éthanoïque cède son proton) et \ce{H3O+}/\ce{H2O} (l'eau joue ici le rôle de \emph{base} : elle capte le proton).
\item Demi-équations :
\[ \ce{CH3COOH} = \ce{CH3COO-} + \ce{H+} \qquad ; \qquad \ce{H2O} + \ce{H+} = \ce{H3O+} \]
\item Somme et simplification de \ce{H+} :
\[ \ce{CH3COOH + H2O <=> CH3COO- + H3O+} \]
\end{enumerate}
Commentaire : cette réaction est \textbf{partielle} (d'où la double flèche \ce{<=>}) : à \SI{25}{\celsius}, seule une faible fraction des molécules d'acide éthanoïque cède leur proton. C'est ce qui fait du vinaigre un acide \emph{doux}, sans danger en cuisine.
\end{exemplebox}

\begin{popfigure}
\begin{center}
\begin{tikzpicture}[scale=1.05, every node/.style={font=\small}]
% Molécule acide éthanoïque
\node[font=\large\bfseries, popblue] (ac) at (0,0) {\chemfig{CH_3-C(=[1]O)-[7]OH}};
\node[below=0.15cm of ac, popblue] {acide éthanoïque (acide)};
% Molécule d'eau
\node[font=\large\bfseries, poporange] (eau) at (6.2,0) {\chemfig{H-[1]O-[7]H}};
\node[below=0.15cm of eau, poporange] {eau (base)};
% Flèche de transfert du proton
\draw[->, very thick, poppink, decorate, decoration={snake, amplitude=0.6mm}] (2.1,0.75) to[bend left=25] node[midway, above, font=\small\bfseries, poppink]{transfert de \ce{H+}} (5.3,0.75);
% Doublet de l'oxygène qui capte
\draw[->, thick, popgreen] (5.15,-0.55) to[bend right=15] (4.6,-0.15);
\node[popgreen, font=\footnotesize] at (6.4,-0.95) {doublet non liant de O};
% Flèche de réaction globale
\draw[->, ultra thick, popdark] (3.1,-1.7) -- (3.1,-2.5);
% Produits
\node[font=\large\bfseries, poppurple] (prod1) at (0.3,-3.3) {\ce{CH3COO-}};
\node[below=0.1cm of prod1, poppurple, font=\footnotesize] {ion éthanoate};
\node[font=\large] at (2.6,-3.3) {$+$};
\node[font=\large\bfseries, popgold] (prod2) at (5.6,-3.3) {\ce{H3O+}};
\node[below=0.1cm of prod2, popgold, font=\footnotesize] {ion oxonium};
\end{tikzpicture}
\end{center}
\legende{Transfert du proton \ce{H+} de l'acide éthanoïque vers une molécule d'eau : l'acide devient ion éthanoate, l'eau devient ion oxonium.}
\end{popfigure}

\courssub{Les ampholytes : des espèces à double personnalité}

\textbf{Intuition.} Certains joueurs de football peuvent être passeurs \emph{ou} receveurs selon l'action. L'eau en est l'exemple parfait : face à l'acide éthanoïque, elle capte un proton (elle est base) ; face à l'ammoniac, elle cède un proton (elle est acide).

\begin{definition}[Ampholyte (espèce amphotère)]
Un \cle{ampholyte} (ou espèce \cle{amphotère}) est une espèce chimique qui peut se comporter \textbf{soit comme un acide, soit comme une base} selon le partenaire rencontré. Elle appartient donc à \textbf{deux couples} : elle est la base de l'un et l'acide de l'autre.
\end{definition}

\textbf{Exemple 1 : l'eau.} Elle appartient aux couples \ce{H3O+}/\ce{H2O} (elle y est la base) et \ce{H2O}/\ce{HO-} (elle y est l'acide). Cette double appartenance explique l'\emph{autoprotolyse} de l'eau, que nous étudierons dans la section suivante.

\textbf{Exemple 2 : l'ion hydrogénocarbonate} \ce{HCO3-}, présent dans le sang et dans le bicarbonate de soude utilisé par les mamans contre les brûlures d'estomac. Il appartient à deux couples :
\begin{itemize}
\item \ce{CO2, H2O}/\ce{HCO3-} : il y est la \emph{base} : $\ce{HCO3-} + \ce{H+} = \ce{CO2, H2O}$
\item \ce{HCO3-}/\ce{CO3^2-} : il y est l'\emph{acide} : $\ce{HCO3-} = \ce{CO3^2-} + \ce{H+}$
\end{itemize}

\begin{savaistu}[Le tampon sanguin : pourquoi ton analyse médicale affiche pH = 7,4]
Quand tu fais une prise de sang à l'hôpital de district, l'un des paramètres surveillés est le pH sanguin, qui doit rester entre 7,35 et 7,45. Or ton métabolisme produit en permanence des acides (acide lactique lors d'un effort, \ce{CO2} de la respiration). Comment le sang résiste-t-il ? Grâce au couple \ce{CO2, H2O}/\ce{HCO3-}, appelé \textbf{tampon bicarbonate} : si un excès d'acide arrive, \ce{HCO3-} le neutralise ; si un excès de base arrive, \ce{CO2, H2O} la neutralise. Un pH sanguin inférieur à 7,35 signale une \emph{acidose}, supérieur à 7,45 une \emph{alcalose} : deux urgences médicales. La chimie de Brønsted veille sur toi à chaque seconde !
\end{savaistu}

\courssub{Acides forts, acides faibles : première approche}

Tous les acides ne cèdent pas leur proton avec la même facilité. En solution aqueuse, on distingue :

\begin{itemize}
\item Les \cle{acides forts} : leur réaction avec l'eau est \textbf{totale} (flèche simple \ce{->}). Exemple : le chlorure d'hydrogène :
\[ \ce{HCl + H2O -> H3O+ + Cl-} \]
Il ne reste pratiquement plus de molécules \ce{HCl} en solution. Sa base conjuguée \ce{Cl-} est dite \emph{indifférente} (elle ne capte plus de proton).
\item Les \cle{acides faibles} : leur réaction avec l'eau est \textbf{partielle} (double flèche \ce{<=>}), un équilibre s'établit. Exemple : l'acide éthanoïque du vinaigre.
\item Les \cle{bases fortes} : réaction totale avec l'eau. Exemple : l'ion hydroxyde \ce{HO-} (soude, potasse). Elles captent le proton de l'eau totalement.
\item Les \cle{bases faibles} : réaction partielle avec l'eau. Exemple : l'ammoniac \ce{NH3}.
\end{itemize}

Cette distinction sera précisée quantitativement dans les leçons suivantes (constantes d'acidité, avancement final) ; pour l'instant, retiens le critère : \textbf{totale = fort, partielle = faible}.

\begin{exoresolu}[L'ammoniaque des produits ménagers est une base faible]
\textbf{Énoncé.} L'ammoniaque vendue au marché pour le nettoyage est une solution d'ammoniac \ce{NH3} de concentration $C = \SI{1,0e-1}{mol.L^{-1}}$. L'ammoniac réagit partiellement avec l'eau. On mesure à l'équilibre une concentration en ions hydroxyde $[\ce{HO-}] = \SI{1,3e-3}{mol.L^{-1}}$.
\begin{enumerate}
\item Identifier les deux couples mis en jeu et écrire l'équation-bilan de la réaction.
\item Calculer le taux d'avancement final $\tau$ de la réaction et conclure sur la force de la base.
\end{enumerate}
\tcblower
\textbf{Solution détaillée.}
\begin{enumerate}
\item \ce{NH3} capte un proton : c'est la base du couple \ce{NH4+}/\ce{NH3}. L'eau cède un proton : c'est l'acide du couple \ce{H2O}/\ce{HO-}.
Demi-équations : $\ce{NH3} + \ce{H+} = \ce{NH4+}$ et $\ce{H2O} = \ce{HO-} + \ce{H+}$.
En additionnant et en simplifiant \ce{H+} :
\[ \ce{NH3 + H2O <=> NH4+ + HO-} \]
\item Dressons le tableau d'avancement (par litre de solution), en négligeant les \ce{HO-} venant de l'autoprotolyse de l'eau :
\begin{center}
\begin{tabular}{lcccc}
\toprule
 & \ce{NH3} & \ce{H2O} & \ce{NH4+} & \ce{HO-} \\
\midrule
Initial (mol) & $0,10$ & excès & $0$ & $\approx 0$ \\
Final (mol) & $0,10 - x_f$ & excès & $x_f$ & $x_f$ \\
\bottomrule
\end{tabular}
\end{center}
Si la réaction était totale, $x_{max} = C \times V = \SI{0,10}{mol}$. Or $x_f = [\ce{HO-}] \times V = \SI{1,3e-3}{mol}$. Donc :
\[ \tau = \frac{x_f}{x_{max}} = \frac{\num{1,3e-3}}{\num{0,10}} = \num{1,3e-2} \approx 1,3\,\% \]
$\tau \ll 100\,\%$ : la réaction est très limitée, l'ammoniac est bien une \textbf{base faible}. C'est pourquoi l'ammoniaque, bien que corrosive à forte dose, peut être manipulée diluée pour le ménage, contrairement à la soude qui est une base forte.
\end{enumerate}
\end{exoresolu}

\courssub{Tableau de synthèse : les couples usuels au quotidien}

\begin{center}
\begin{tabularx}{\linewidth}{|l|l|X|}
\hline
\rowcolor{popblueL} \textbf{Couple acide/base} & \textbf{Demi-équation} & \textbf{Exemple camerounais} \\
\hline
\ce{H3O+}/\ce{H2O} & $\ce{H3O+} = \ce{H2O} + \ce{H+}$ & Ion responsable de l'acidité du jus de bissap ou du tamarin \\
\hline
\ce{H2O}/\ce{HO-} & $\ce{H2O} = \ce{HO-} + \ce{H+}$ & L'eau, ampholyte universel \\
\hline
\ce{CH3COOH}/\ce{CH3COO-} & $\ce{CH3COOH} = \ce{CH3COO-} + \ce{H+}$ & Vinaigre de table ; vin de palme qui « tourne » à l'air libre \\
\hline
\ce{NH4+}/\ce{NH3} & $\ce{NH4+} = \ce{NH3} + \ce{H+}$ & Ammoniaque des détergents ménagers \\
\hline
\ce{HCl}/\ce{Cl-} & $\ce{HCl} = \ce{Cl-} + \ce{H+}$ & Acide chlorhydrique (détartrant), acide fort \\
\hline
\ce{CO2, H2O}/\ce{HCO3-} & $\ce{CO2, H2O} = \ce{HCO3-} + \ce{H+}$ & Tampon du sang ; eaux gazeuses \\
\hline
\ce{HCO3-}/\ce{CO3^2-} & $\ce{HCO3-} = \ce{CO3^2-} + \ce{H+}$ & Bicarbonate contre l'acidité gastrique \\
\hline
\end{tabularx}
\end{center}

\begin{aretenir}[L'essentiel de la section]
\begin{itemize}
\item Selon Brønsted : un \cle{acide} cède \ce{H+}, une \cle{base} capte \ce{H+}.
\item Un \cle{couple acide/base} relie deux espèces échangeant un proton : acide = base + \ce{H+}.
\item Une \cle{réaction acido-basique} est un transfert de proton entre l'acide d'un couple et la base d'un autre couple ; \ce{H+} n'apparaît jamais dans le bilan (en solution, il est fixé sous forme \ce{H3O+}).
\item Un \cle{ampholyte} (eau, \ce{HCO3-}) est à la fois acide et base.
\item Acide ou base \textbf{fort(e)} : réaction totale avec l'eau ; \textbf{faible} : réaction partielle, équilibre.
\end{itemize}
\end{aretenir}

\begin{exercice}[À toi de jouer]
L'acide lactique \ce{CH3-CHOH-COOH}, produit dans les muscles lors de l'effort, appartient au couple \ce{CH3-CHOH-COOH}/\ce{CH3-CHOH-COO-}.
\begin{enumerate}
\item Écris la demi-équation de ce couple.
\item Écris l'équation-bilan de la réaction de l'acide lactique avec l'ion hydrogénocarbonate \ce{HCO3-} (base du couple \ce{CO2, H2O}/\ce{HCO3-}).
\item L'ion lactate est-il un acide ou une base ? Justifie.
\end{enumerate}
\end{exercice}

\begin{corrige}[Exercice : acide lactique]
\begin{enumerate}
\item $\ce{CH3-CHOH-COOH} = \ce{CH3-CHOH-COO-} + \ce{H+}$.
\item Demi-équation de la base : $\ce{HCO3-} + \ce{H+} = \ce{CO2, H2O}$. En additionnant et en simplifiant \ce{H+} :
\[ \ce{CH3-CHOH-COOH + HCO3- -> CH3-CHOH-COO- + CO2, H2O} \]
\item L'ion lactate \ce{CH3-CHOH-COO-} est la \textbf{base conjuguée} de l'acide lactique : il est capable de capter un proton pour régénérer l'acide.
\end{enumerate}
\end{corrige}

Maintenant que tu maîtrises les couples et les transferts de proton, une question surgit naturellement : l'eau étant à la fois acide et base, peut-elle réagir \emph{avec elle-même} ? C'est l'objet de la section suivante : l'autoprotolyse de l'eau et son produit ionique $K_e$, clé de tous les calculs de pH.

\coursec{Autoprotolyse de l'eau et produit ionique \texorpdfstring{$K_e$}{Ke}}

Nous venons de voir qu'au sens de Brønsted, l'eau est un cas remarquable : elle peut \emph{donner} un proton (elle se comporte alors comme un acide) ou en \emph{capter} un (elle se comporte alors comme une base). On dit qu'elle est \cle{amphotère} (ou ampholyte). Cette double personnalité va avoir une conséquence capitale : l'eau peut réagir \emph{avec elle-même}, une molécule jouant le rôle d'acide et l'autre celui de base. C'est cette réaction, discrète mais fondamentale, que nous allons étudier maintenant. Elle gouverne le comportement de \emph{toutes} les solutions aqueuses : ton sang, le jus de bissap, l'eau de pluie ou la soude de la savonnerie.

\courssub{Observation : l'eau pure conduit (très faiblement) le courant}

\begin{experience}[L'eau pure est-elle vraiment isolante ?]
\textbf{Dispositif.} On plonge deux électrodes reliées à un générateur et à un ampèremètre très sensible (microampèremètre) dans un bécher d'eau distillée. On compare avec une solution de chlorure de sodium.
\par\medskip
\textbf{Observations.}
\begin{itemize}
\item Avec l'eau distillée : l'ampèremètre détecte un courant \emph{extrêmement faible} (quelques microampères), mais non nul.
\item Avec l'eau salée : le courant est intense, la lampe du circuit brille franchement.
\end{itemize}
\textbf{Interprétation.} Un courant électrique dans une solution ne peut circuler que grâce au déplacement d'\cle{ions}. L'eau pure conduit très faiblement le courant : elle contient donc des ions, mais en \emph{quantité infime}. D'où viennent-ils ? Ils ne peuvent provenir que de l'eau elle-même.
\end{experience}

\begin{attention}[Erreur fréquente n\textsuperscript{o} 1]
Beaucoup d'élèves affirment : « l'eau pure ne contient aucun ion, c'est un isolant parfait ». \textbf{C'est faux !} L'eau pure contient toujours des ions \ce{H3O+} et \ce{HO-}, issus de sa propre dissociation. À \SI{25}{\celsius}, leur concentration est de $10^{-7}$~\si{mol.L^{-1}} : c'est très faible devant les \SI{55,5}{mol.L^{-1}} de molécules d'eau, mais ce n'est \emph{pas zéro}. Toute la chimie du pH repose sur cette présence infime mais réelle.
\end{attention}

\courssub{L'équation-bilan de l'autoprotolyse}

Puisque l'eau est amphotère, imaginons deux molécules d'eau qui se rencontrent : l'une cède un proton \ce{H+} (rôle d'acide), l'autre le capte (rôle de base). Il se forme un ion \cle{oxonium} \ce{H3O+} et un ion \cle{hydroxyde} \ce{HO-} :

\[ \ce{H2O + H2O <=> H3O+ + HO-} \quad \text{ou plus simplement} \quad \ce{2 H2O <=> H3O+ + HO-} \]

\begin{definition}[Autoprotolyse de l'eau]
L'\cle{autoprotolyse} de l'eau est la réaction de transfert de proton \emph{entre deux molécules d'eau} :
\[ \ce{2 H2O <=> H3O+ + HO-} \]
Le préfixe \emph{auto-} signifie que l'eau réagit sur elle-même ; \emph{protolyse} signifie « coupure par transfert de proton ». Cette réaction est :
\begin{itemize}
\item \textbf{très limitée} (équilibre très déplacé vers la gauche) : l'eau reste majoritairement sous forme de molécules \ce{H2O} ;
\item \textbf{rapide} et \textbf{réversible} (d'où la double flèche \ce{<=>}) ;
\item \textbf{endothermique} : elle absorbe de la chaleur, donc elle est favorisée par une élévation de température.
\end{itemize}
\end{definition}

\begin{popfigure}
\centering
\begin{tikzpicture}[scale=1.0, every node/.style={font=\small}]
% Molécule 1 (acide) : H2O qui donne H+
\draw[fill=popblue!25, draw=popblue] (0,0) circle (0.42) node {O};
\draw[fill=popblue!10, draw=popblue] (-0.55,0.45) circle (0.22) node {\scriptsize H};
\draw[fill=popblue!10, draw=popblue] (0.55,0.45) circle (0.22) node {\scriptsize H};
\draw[popblue, thick] (-0.36,0.28) -- (-0.18,0.12);
\draw[popblue, thick] (0.36,0.28) -- (0.18,0.12);
\node[below, popblue] at (0,-0.5) {\ce{H2O} (acide)};
% Molécule 2 (base) : H2O qui capte H+
\draw[fill=poporange!30, draw=poporange] (3.2,0) circle (0.42) node {O};
\draw[fill=poporange!15, draw=poporange] (3.75,0.45) circle (0.22) node {\scriptsize H};
\draw[fill=poporange!15, draw=poporange] (2.65,0.45) circle (0.22) node {\scriptsize H};
\draw[poporange, thick] (2.84,0.28) -- (3.02,0.12);
\draw[poporange, thick] (3.56,0.28) -- (3.38,0.12);
\node[below, poporange] at (3.2,-0.5) {\ce{H2O} (base)};
% Flèche transfert de proton
\draw[-{Stealth[length=3mm]}, poppurple, very thick, dashed] (0.6,0.55) .. controls (1.6,1.4) .. (2.55,0.6);
\node[poppurple, above] at (1.6,1.15) {transfert de \ce{H+}};
% Flèche d'équilibre
\draw[-{Stealth[length=3mm]}, popdark, very thick] (4.6,0.1) -- (5.8,0.1);
\draw[-{Stealth[length=3mm]}, popdark, very thick] (5.8,-0.25) -- (4.6,-0.25);
% Produits : H3O+
\draw[fill=popgreen!30, draw=popgreen] (7.0,0) circle (0.42) node {O};
\draw[fill=popgreen!15, draw=popgreen] (6.45,0.45) circle (0.22) node {\scriptsize H};
\draw[fill=popgreen!15, draw=popgreen] (7.55,0.45) circle (0.22) node {\scriptsize H};
\draw[fill=popgreen!15, draw=popgreen] (7.0,0.62) circle (0.22) node {\scriptsize H};
\draw[popgreen, thick] (6.64,0.28) -- (6.82,0.12);
\draw[popgreen, thick] (7.36,0.28) -- (7.18,0.12);
\draw[popgreen, thick] (7.0,0.42) -- (7.0,0.2);
\node[popgreen] at (7.6,-0.15) {$+$};
\node[below, popgreen] at (7.0,-0.5) {\ce{H3O+}};
% Produit : HO-
\draw[fill=poppink!30, draw=poppink] (9.2,0) circle (0.42) node {O};
\draw[fill=poppink!15, draw=poppink] (9.75,0.45) circle (0.22) node {\scriptsize H};
\draw[poppink, thick] (9.56,0.28) -- (9.38,0.12);
\node[poppink] at (9.75,-0.2) {$-$};
\node[below, poppink] at (9.2,-0.5) {\ce{HO-}};
\node[popdark] at (8.1,0.55) {$+$};
\end{tikzpicture}
\legende{Schéma de l'autoprotolyse : une molécule d'eau (acide) cède un proton \ce{H+} à une autre molécule d'eau (base). Il se forme un ion oxonium \ce{H3O+} et un ion hydroxyde \ce{HO-}.}
\end{popfigure}

\begin{exemplebox}[À quel point la réaction est-elle limitée ?]
Dans l'eau pure à \SI{25}{\celsius}, la concentration en ions oxonium est $[\ce{H3O+}] = 10^{-7}$~\si{mol.L^{-1}}, alors que la concentration en molécules d'eau vaut environ \SI{55,5}{mol.L^{-1}}. La proportion de molécules ionisées est donc :
\[ \frac{10^{-7}}{55,5} \approx 1,8 \times 10^{-9} \]
Autrement dit, \textbf{seulement 2 molécules d'eau sur 555 millions sont ionisées} à un instant donné. L'autoprotolyse est bien une réaction \emph{très limitée} : l'eau est et reste essentiellement constituée de molécules \ce{H2O}.
\end{exemplebox}

\courssub{Le produit ionique de l'eau : \texorpdfstring{$K_e$}{Ke} et \texorpdfstring{$pK_e$}{pKe}}

L'autoprotolyse est un équilibre chimique. Comme tout équilibre, il est caractérisé par une constante. L'eau étant le solvant (sa « concentration » est constante et n'apparaît pas dans l'expression), la constante d'équilibre s'écrit simplement :

\begin{definition}[Produit ionique de l'eau]
Le \cle{produit ionique de l'eau}, noté \cle{$K_e$}, est la constante d'équilibre de l'autoprotolyse :
\[ K_e = [\ce{H3O+}] \times [\ce{HO-}] \]
C'est une grandeur \textbf{sans unité} (les concentrations sont rapportées à la concentration de référence $c^0 = 1$~\si{mol.L^{-1}}). À \SI{25}{\celsius} :
\[ K_e = 1,0 \times 10^{-14} \]
On définit aussi le \cle{$pK_e$} :
\[ pK_e = -\log K_e \qquad \text{soit à \SI{25}{\celsius} : } pK_e = 14 \]
\end{definition}

\begin{aretenir}[La relation fondamentale]
Dans \textbf{toute solution aqueuse}, à \textbf{toute température}, qu'elle soit acide, neutre ou basique, les ions \ce{H3O+} et \ce{HO-} sont \emph{toujours présents simultanément} et leurs concentrations vérifient :
\[ [\ce{H3O+}] \times [\ce{HO-}] = K_e \]
À \SI{25}{\celsius} : $[\ce{H3O+}] \times [\ce{HO-}] = 1,0 \times 10^{-14}$.
\par\medskip
Conséquence immédiate : les deux concentrations sont \emph{liées}. Si l'on ajoute un acide, $[\ce{H3O+}]$ augmente, donc $[\ce{HO-}]$ diminue automatiquement pour que le produit reste égal à $K_e$. On ne peut jamais avoir une solution aqueuse sans ion \ce{HO-}, pas même un acide concentré !
\end{aretenir}

\courssub{Influence de la température sur \texorpdfstring{$K_e$}{Ke}}

L'autoprotolyse est \cle{endothermique} : il faut fournir de la chaleur pour « casser » des molécules d'eau en ions. D'après le principe de Le Chatelier (loi de modération), une élévation de température favorise le sens endothermique, donc le sens de formation des ions : \textbf{$K_e$ augmente quand la température augmente}.

\begin{popfigure}
\centering
\begin{tikzpicture}
\begin{axis}[
width=0.85\linewidth, height=6.2cm,
xlabel={Température $\theta$ (\si{\celsius})},
ylabel={$pK_e$},
xmin=0, xmax=100, ymin=12.8, ymax=15.2,
grid=major, grid style={popline!40},
xtick={0,25,37,50,75,100},
label style={font=\small}, ticklabel style={font=\small},
]
\addplot[popcurve, domain=0:100, samples=100] {14.94 - 0.0375*x + 0.000075*x*x};
\addplot[only marks, mark=*, poporange, mark size=2pt] coordinates {(0,14.94) (25,14.0) (37,13.62) (50,13.26) (100,12.25)};
\node[poporange, font=\scriptsize, anchor=south west] at (axis cs:27,14.0) {25 °C : $pK_e = 14$};
\node[poporange, font=\scriptsize, anchor=south] at (axis cs:37,13.72) {37 °C : $pK_e \approx 13{,}6$};
\node[popmuted, font=\scriptsize, anchor=north west] at (axis cs:2,15.1) {$pK_e$ diminue (donc $K_e$ augmente) quand $\theta$ augmente};
\end{axis}
\end{tikzpicture}
\legende{Évolution du $pK_e$ de l'eau en fonction de la température. L'autoprotolyse étant endothermique, $K_e$ augmente (et $pK_e$ diminue) quand la température s'élève.}
\end{popfigure}

\begin{center}
\begin{tabularx}{\linewidth}{|l|X|X|X|}
\hline
\rowcolor{popblueL} $\theta$ (\si{\celsius}) & $K_e$ & $pK_e$ & Neutralité ($[\ce{H3O+}] = \sqrt{K_e}$) \\
\hline
0 & $1,1 \times 10^{-15}$ & 14,94 & $[\ce{H3O+}] = 3,3 \times 10^{-8}$~\si{mol.L^{-1}} \\
\hline
25 & $1,0 \times 10^{-14}$ & 14,00 & $[\ce{H3O+}] = 1,0 \times 10^{-7}$~\si{mol.L^{-1}} \\
\hline
37 & $2,4 \times 10^{-14}$ & 13,62 & $[\ce{H3O+}] = 1,5 \times 10^{-7}$~\si{mol.L^{-1}} \\
\hline
100 & $5,5 \times 10^{-13}$ & 12,26 & $[\ce{H3O+}] = 7,4 \times 10^{-7}$~\si{mol.L^{-1}} \\
\hline
\end{tabularx}
\end{center}

\begin{attention}[La neutralité dépend de la température !]
Dans l'eau \emph{pure}, l'autoprotolyse produit autant de \ce{H3O+} que de \ce{HO-} : $[\ce{H3O+}] = [\ce{HO-}] = \sqrt{K_e}$. À \SI{25}{\celsius}, cela donne $10^{-7}$~\si{mol.L^{-1}}, d'où le fameux « pH neutre = 7 ». Mais à \SI{100}{\celsius}, la neutralité correspond à $[\ce{H3O+}] = 7,4 \times 10^{-7}$~\si{mol.L^{-1}}, soit un pH d'environ 6,1 ! \textbf{Le pH 7 n'est synonyme de neutralité qu'à \SI{25}{\celsius}.} Retiens : la valeur $K_e = 10^{-14}$ et $pK_e = 14$ n'est valable qu'à \SI{25}{\celsius}, température de référence de tous les exercices du baccalauréat sauf indication contraire.
\end{attention}

\begin{savaistu}[Le pH de ton sang et la température du corps]
Le corps humain fonctionne à \SI{37}{\celsius}, pas à \SI{25}{\celsius}. À cette température, $K_e \approx 2,4 \times 10^{-14}$, donc $pK_e \approx 13,6$ et la neutralité se situe à $pH \approx 6,8$ et non 7 ! C'est un point crucial pour interpréter les analyses médicales : le sang artériel normal a un pH compris entre 7,38 et 7,42, ce qui le rend \emph{légèrement basique} par rapport à la neutralité à \SI{37}{\celsius}. Lorsque le laboratoire d'analyses de l'hôpital mesure ton pH sanguin, l'appareil est étalonné en tenant compte de la température. Un pH sanguin inférieur à 7,35 signale une \emph{acidose}, supérieur à 7,45 une \emph{alcalose} : deux situations dangereuses que le médecin doit corriger rapidement.
\end{savaistu}

\courssub{Exploitation : calculer une concentration ionique à partir de l'autre}

La relation $[\ce{H3O+}] \times [\ce{HO-}] = K_e$ permet, connaissant l'une des deux concentrations, de trouver l'autre immédiatement.

\begin{methode}[{Calculer [\ce{HO-}] à partir de [\ce{H3O+}] (et inversement)}]
\begin{enumerate}
\item Vérifier la température et noter la valeur de $K_e$ correspondante (à \SI{25}{\celsius} : $K_e = 1,0 \times 10^{-14}$).
\item Écrire la relation fondamentale : $[\ce{H3O+}] \times [\ce{HO-}] = K_e$.
\item Isoler l'inconnue :
\[ [\ce{HO-}] = \frac{K_e}{[\ce{H3O+}]} \qquad \text{ou} \qquad [\ce{H3O+}] = \frac{K_e}{[\ce{HO-}]} \]
\item Effectuer le calcul en veillant aux puissances de 10, puis conclure avec l'unité \si{mol.L^{-1}}.
\item Vérifier la cohérence : si $[\ce{H3O+}] > 10^{-7}$~\si{mol.L^{-1}} (milieu acide), alors $[\ce{HO-}]$ doit être \emph{inférieure} à $10^{-7}$~\si{mol.L^{-1}}.
\end{enumerate}
\end{methode}

\begin{exoresolu}[Concentrations ioniques dans un jus de bissap]
\textbf{Énoncé.} Le jus de bissap (infusion de fleurs d'hibiscus) est naturellement acide. Une mesure donne $[\ce{H3O+}] = 2,5 \times 10^{-4}$~\si{mol.L^{-1}} à \SI{25}{\celsius}.
\begin{enumerate}
\item Calculer la concentration en ions hydroxyde $[\ce{HO-}]$ de ce jus.
\item Comparer les deux concentrations et conclure sur les espèces majoritaires.
\end{enumerate}
\tcblower
\textbf{Solution détaillée.}
\begin{enumerate}
\item À \SI{25}{\celsius}, $K_e = 1,0 \times 10^{-14}$. La relation fondamentale s'écrit :
\[ [\ce{HO-}] = \frac{K_e}{[\ce{H3O+}]} = \frac{1,0 \times 10^{-14}}{2,5 \times 10^{-4}} = 4,0 \times 10^{-11}~\si{mol.L^{-1}} \]
\item Comparons :
\[ \frac{[\ce{H3O+}]}{[\ce{HO-}]} = \frac{2,5 \times 10^{-4}}{4,0 \times 10^{-11}} = 6,3 \times 10^{6} \]
Il y a plus de \textbf{6 millions} de fois plus d'ions oxonium que d'ions hydroxyde ! L'ion \ce{H3O+} est l'espèce \emph{majoritaire} (c'est lui qui impose le caractère acide du jus), l'ion \ce{HO-} est \emph{ultra-minoritaire} mais toujours présent. Vérification : $[\ce{H3O+}] > 10^{-7}$~\si{mol.L^{-1}}, le milieu est bien acide, et $[\ce{HO-}] < 10^{-7}$~\si{mol.L^{-1}}, ce qui est cohérent.
\end{enumerate}
\end{exoresolu}

\begin{exoresolu}[Et dans une solution basique ?]
\textbf{Énoncé.} Une savonnerie artisanale prépare une lessive de soude dont la concentration en ions hydroxyde est $[\ce{HO-}] = 5,0 \times 10^{-3}$~\si{mol.L^{-1}} à \SI{25}{\celsius}. Calculer $[\ce{H3O+}]$ dans cette solution.
\tcblower
\textbf{Solution détaillée.}
\[ [\ce{H3O+}] = \frac{K_e}{[\ce{HO-}]} = \frac{1,0 \times 10^{-14}}{5,0 \times 10^{-3}} = 2,0 \times 10^{-12}~\si{mol.L^{-1}} \]
La concentration en ions oxonium est infime ($[\ce{H3O+}] \ll 10^{-7}$~\si{mol.L^{-1}}) : l'ion \ce{H3O+} est ultra-minoritaire, l'ion \ce{HO-} est majoritaire. La solution est bien basique, comme attendu pour une lessive de soude utilisée en saponification.
\end{exoresolu}

\begin{attention}[Pièges classiques à éviter]
\begin{itemize}
\item \textbf{Oublier que $K_e$ dépend de la température.} Si l'énoncé précise une température différente de \SI{25}{\celsius}, il doit te donner la valeur de $K_e$ : utilise-la, et non $10^{-14}$ !
\item \textbf{Croire qu'une solution acide ne contient pas d'ions \ce{HO-}.} Faux : il y en a toujours, en quantité infime, imposée par $K_e$.
\item \textbf{Se tromper dans les puissances de 10.} Rappel : $\dfrac{10^{-14}}{10^{-4}} = 10^{-14-(-4)} = 10^{-10}$. Pose le calcul proprement.
\item \textbf{Confondre $K_e$ et $pK_e$.} $K_e = 10^{-14}$ (sans unité) ; $pK_e = 14$ (sans unité non plus). La relation est $pK_e = -\log K_e$, soit $K_e = 10^{-pK_e}$.
\end{itemize}
\end{attention}

\begin{exercice}[À toi de jouer]
À \SI{25}{\celsius}, on considère trois solutions aqueuses :
\begin{itemize}
\item Solution A (eau de source) : $[\ce{H3O+}] = 3,2 \times 10^{-7}$~\si{mol.L^{-1}} ;
\item Solution B (vinaigre) : $[\ce{H3O+}] = 1,6 \times 10^{-3}$~\si{mol.L^{-1}} ;
\item Solution C (eau de Javel diluée) : $[\ce{HO-}] = 4,0 \times 10^{-4}$~\si{mol.L^{-1}}.
\end{itemize}
\begin{enumerate}
\item Calculer la concentration manquante ($[\ce{HO-}]$ ou $[\ce{H3O+}]$) pour chaque solution.
\item Préciser, pour chaque solution, l'ion majoritaire et le caractère acide ou basique.
\end{enumerate}
\end{exercice}

\begin{corrige}[Exercice : À toi de jouer]
\begin{enumerate}
\item On applique $[\ce{H3O+}] \times [\ce{HO-}] = 1,0 \times 10^{-14}$ :
\begin{itemize}
\item Solution A : $[\ce{HO-}] = \dfrac{1,0 \times 10^{-14}}{3,2 \times 10^{-7}} = 3,1 \times 10^{-8}$~\si{mol.L^{-1}} ;
\item Solution B : $[\ce{HO-}] = \dfrac{1,0 \times 10^{-14}}{1,6 \times 10^{-3}} = 6,3 \times 10^{-12}$~\si{mol.L^{-1}} ;
\item Solution C : $[\ce{H3O+}] = \dfrac{1,0 \times 10^{-14}}{4,0 \times 10^{-4}} = 2,5 \times 10^{-11}$~\si{mol.L^{-1}}.
\end{itemize}
\item Solution A : $[\ce{H3O+}] > [\ce{HO-}]$, légèrement acide, \ce{H3O+} majoritaire. Solution B : $[\ce{H3O+}] \gg [\ce{HO-}]$, nettement acide, \ce{H3O+} majoritaire et \ce{HO-} ultra-minoritaire. Solution C : $[\ce{HO-}] \gg [\ce{H3O+}]$, basique, \ce{HO-} majoritaire et \ce{H3O+} ultra-minoritaire.
\end{enumerate}
\end{corrige}

\begin{aretenir}[L'essentiel de la section]
\begin{itemize}
\item L'eau pure conduit très faiblement le courant : elle contient des ions \ce{H3O+} et \ce{HO-} issus de l'\cle{autoprotolyse} : $\ce{2 H2O <=> H3O+ + HO-}$, réaction très limitée et endothermique.
\item Le \cle{produit ionique} $K_e = [\ce{H3O+}] \times [\ce{HO-}]$ vaut $1,0 \times 10^{-14}$ à \SI{25}{\celsius} ($pK_e = 14$) et \textbf{augmente avec la température}.
\item Dans \emph{toute} solution aqueuse, $[\ce{H3O+}] \times [\ce{HO-}] = K_e$ : connaître l'une des concentrations, c'est connaître l'autre.
\item Dans l'eau pure à \SI{25}{\celsius} : $[\ce{H3O+}] = [\ce{HO-}] = 10^{-7}$~\si{mol.L^{-1}} (2 molécules ionisées sur 555 millions).
\end{itemize}
Cette relation fondamentale est la clé de la suite : nous allons maintenant définir une échelle pratique pour exprimer $[\ce{H3O+}]$ sans manipuler des puissances de dix : le fameux \cle{pH}.
\end{aretenir}

\coursec{Le pH : définition et mesure}

Dans la section précédente, nous avons découvert que l'eau pure, grâce à l'\cle{autoprotolyse}, contient toujours des ions \ce{H3O+} et \ce{HO-}, et que leur produit $K_e = [\ce{H3O+}]\times[\ce{HO-}]$ est constant à température donnée ($K_e = 10^{-14}$ à \SI{25}{\celsius}). Nous savons donc que toute solution aqueuse contient des ions oxonium \ce{H3O+} : c'est leur concentration qui fait la différence entre le jus de bissap bien acide, l'eau de forage de ton quartier et l'eau de lessive de maman. Mais voilà : ces concentrations sont minuscules et s'expriment avec des puissances de dix négatives peu parlantes ($10^{-2}$, $10^{-7}$, $10^{-11}$...). Les chimistes ont donc inventé une \emph{échelle pratique}, simple et universelle, pour comparer l'acidité des solutions : le \cle{pH}. C'est l'objet de cette section, capitale pour le Baccalauréat TI et pour la lecture de tes futures analyses médicales.

\courssub{Définition du pH}

\textbf{Pourquoi une échelle logarithmique ?}\par
La concentration en ions oxonium $[\ce{H3O+}]$ varie énormément d'une solution à l'autre : de quelques moles par litre dans un acide concentré à $10^{-14}$ mol/L dans une base très concentrée. Manier des nombres aussi étalés est inconfortable. L'idée géniale du Danois Sørensen (1909) fut d'appliquer la fonction \emph{logarithme décimal} pour « compresser » cette échelle en nombres simples, généralement compris entre 0 et 14.

\begin{definition}[Définition du pH]
Le \cle{pH} (potentiel Hydrogène) d'une solution aqueuse diluée est défini par :
\[ \boxed{\;\mathrm{pH} = -\log\,[\ce{H3O+}]\;} \]
où $[\ce{H3O+}]$ est la concentration molaire des ions oxonium, exprimée en \si{mol.L^{-1}}. Le pH est un nombre \textbf{sans unité}.

Relation inverse, tout aussi importante :
\[ \boxed{\;[\ce{H3O+}] = 10^{-\mathrm{pH}}\;} \]
\end{definition}

\begin{attention}[Conditions de validité]
La relation $\mathrm{pH} = -\log[\ce{H3O+}]$ n'est rigoureusement valable que pour des solutions \textbf{diluées} ($[\ce{H3O+}] \leq \SI{1}{mol.L^{-1}}$ environ). Pour des solutions concentrées, il faut utiliser la notion d'\emph{activité}, qui dépasse le programme de Terminale TI. Au Baccalauréat, applique la formule sans hésiter dès que la solution est diluée.
\end{attention}

\textbf{Deux conséquences immédiates à maîtriser :}
\begin{itemize}
\item Plus $[\ce{H3O+}]$ est \textbf{grande}, plus le pH est \textbf{petit} : le pH varie en sens inverse de la concentration en ions oxonium. Une solution très acide a un pH faible.
\item Quand $[\ce{H3O+}]$ est \textbf{multipliée par 10}, le pH \textbf{diminue de 1 unité} : en effet, $-\log(10\times[\ce{H3O+}]) = \mathrm{pH} - 1$. Inversement, diluer 10 fois un acide fort augmente son pH de 1.
\end{itemize}

\begin{exemplebox}[Le jus de citron, 100 000 fois plus acide que l'eau pure]
Le jus de citron a un pH voisin de 2. Sa concentration en ions oxonium vaut :
\[ [\ce{H3O+}] = 10^{-\mathrm{pH}} = 10^{-2}\ \si{mol.L^{-1}} \]
L'eau pure à \SI{25}{\celsius} a un pH de 7, donc $[\ce{H3O+}] = 10^{-7}\ \si{mol.L^{-1}}$.
Le rapport des concentrations est :
\[ \frac{10^{-2}}{10^{-7}} = 10^{5} = 100\,000 \]
Le jus de citron contient donc \textbf{100 000 fois plus} d'ions oxonium que l'eau pure, alors que leurs pH ne diffèrent que de 5 unités. C'est toute la puissance (et le piège !) de l'échelle logarithmique.
\end{exemplebox}

\begin{attention}[Erreur fréquente : le pH n'est pas forcément entre 0 et 14]
L'échelle 0--14 n'est qu'\textbf{usuelle} : elle correspond aux solutions diluées. Pour des solutions \textbf{concentrées}, le pH peut sortir de cet intervalle :
\begin{itemize}
\item une solution d'acide chlorhydrique à \SI{2}{mol.L^{-1}} a un pH voisin de $-\log 2 \approx -0{,}3$ : un pH \textbf{négatif} est possible ;
\item une solution d'hydroxyde de sodium (soude) à \SI{5}{mol.L^{-1}} a un pH supérieur à 14.
\end{itemize}
Ne dis jamais « un pH ne peut pas dépasser 14 » : dis « pour les solutions diluées usuelles, le pH est compris entre 0 et 14 ».
\end{attention}

\courssub{Solutions neutres, acides et basiques}

Rappelons la conséquence de l'autoprotolyse : dans \emph{toute} solution aqueuse, $[\ce{H3O+}]\times[\ce{HO-}] = K_e$. Comparer les deux concentrations permet alors de classer la solution.

\begin{definition}[Neutralité, acidité, basicité]
Une solution aqueuse est dite :
\begin{itemize}
\item \cle{neutre} si $[\ce{H3O+}] = [\ce{HO-}]$ ;
\item \cle{acide} si $[\ce{H3O+}] > [\ce{HO-}]$ ;
\item \cle{basique} si $[\ce{H3O+}] < [\ce{HO-}]$.
\end{itemize}
À \SI{25}{\celsius}, $K_e = 10^{-14}$, donc dans une solution neutre : $[\ce{H3O+}] = \sqrt{K_e} = 10^{-7}\ \si{mol.L^{-1}}$, c'est-à-dire $\mathrm{pH} = 7$. D'où le critère pratique à \SI{25}{\celsius} :
\begin{center}
acide : $\mathrm{pH} < 7$ \qquad neutre : $\mathrm{pH} = 7$ \qquad basique : $\mathrm{pH} > 7$
\end{center}
\end{definition}

\begin{attention}[La neutralité dépend de la température]
Le produit ionique $K_e$ \textbf{augmente avec la température} (l'autoprotolyse est endothermique). Le pH de neutralité vaut $\mathrm{pH}_{\text{neutre}} = \dfrac{\mathrm{p}K_e}{2}$ :
\begin{itemize}
\item à \SI{0}{\celsius} : $\mathrm{p}K_e = 14{,}9$, donc $\mathrm{pH}_{\text{neutre}} \approx 7{,}45$ ;
\item à \SI{25}{\celsius} : $\mathrm{p}K_e = 14$, donc $\mathrm{pH}_{\text{neutre}} = 7$ ;
\item à \SI{60}{\celsius} : $\mathrm{p}K_e = 13$, donc $\mathrm{pH}_{\text{neutre}} = 6{,}5$.
\end{itemize}
Ainsi, une eau pure à \SI{60}{\celsius} a un pH de 6,5 tout en restant \textbf{parfaitement neutre} ($[\ce{H3O+}] = [\ce{HO-}]$). Le fameux « pH = 7 » n'est valable qu'à \SI{25}{\celsius} : précise toujours la température dans tes copies !
\end{attention}

\begin{popfigure}
\begin{tikzpicture}[x=0.62cm,y=1cm]
% Bandeau coloré de l'échelle de pH (couleurs déjà éclaircies à 75 % dans la liste)
\foreach \i/\col in {
  0/poppink!75,
  1/poppink!75,
  2/poporange!75,
  3/poporange!75,
  4/popgold!75,
  5/popgold!75,
  6/popgold!75,
  7/popgreen!75,
  8/popgreen!70!popblue,
  9/popblue!60,
  10/popblue!75,
  11/popblue!75,
  12/poppurple!75,
  13/poppurple!75,
  14/poppurple!75
}{
  \fill[\col] (\i,0) rectangle (\i+1,0.9);
}
\draw[thick,popdark] (0,0) rectangle (15,0.9);
\foreach \i in {0,...,15}{\draw[popdark] (\i,0) -- (\i,-0.12);}
\foreach \i in {0,...,14}{\node[below,font=\small\bfseries,text=popdark] at (\i+0.5,-0.15) {\i};}
% Zones acide / neutre / basique
\draw[decorate,decoration={brace,amplitude=5pt},popink,thick] (0.05,1.05) -- (6.95,1.05) node[midway,above=6pt,font=\small\bfseries,text=popink]{Zone acide : $\mathrm{pH}<7$};
\node[font=\small\bfseries,text=popgreen!60!popdark] at (7.5,1.55) {neutre};
\draw[->,thick,popgreen!60!popdark] (7.5,1.35) -- (7.5,1.0);
\draw[decorate,decoration={brace,amplitude=5pt},popblue,thick] (8.05,1.05) -- (14.95,1.05) node[midway,above=6pt,font=\small\bfseries,text=popblue]{Zone basique : $\mathrm{pH}>7$};
% Exemples de produits usuels
\node[below,font=\footnotesize,text=popdark,align=center] at (2.5,-0.85) {jus de citron\\ vinaigre};
\draw[->,popmuted] (2.5,-0.75) -- (2.5,-0.35);
\node[below,font=\footnotesize,text=popdark,align=center] at (5.5,-0.85) {jus de bissap};
\draw[->,popmuted] (5.5,-0.75) -- (5.5,-0.35);
\node[below,font=\footnotesize,text=popdark,align=center] at (7.5,-0.85) {eau de forage\\ (pure)};
\draw[->,popmuted] (7.5,-0.75) -- (7.5,-0.35);
\node[below,font=\footnotesize,text=popdark,align=center] at (9.5,-0.85) {savon de\\ Marseille};
\draw[->,popmuted] (9.5,-0.75) -- (9.5,-0.35);
\node[below,font=\footnotesize,text=popdark,align=center] at (13,-0.85) {lessive\\ eau de Javel};
\draw[->,popmuted] (13,-0.75) -- (13,-0.35);
\end{tikzpicture}
\legende{Échelle usuelle de pH à \SI{25}{\celsius} : plus on va vers la gauche, plus la solution est acide ; plus on va vers la droite, plus elle est basique. Quelques produits usuels du quotidien camerounais y sont placés.}
\end{popfigure}

\courssub{Mesure du pH}

Comment connaître concrètement le pH d'une solution ? Deux grandes familles de méthodes existent : les méthodes \emph{colorimétriques} (papier pH, indicateurs colorés), rapides mais approximatives, et la méthode \emph{électrochimique} (pH-mètre), précise.

\begin{methode}[Mesurer un pH : papier pH ou pH-mètre ?]
\textbf{Avec le papier pH (mesure rapide, approximative) :}
\begin{enumerate}
\item Déposer une goutte de la solution sur une bandelette de papier pH à l'aide d'une baguette de verre (ne jamais tremper la bandelette dans le flacon, pour ne pas contaminer la solution).
\item Comparer la couleur obtenue à l'échelle de teintes imprimée sur la boîte.
\item Lire le pH : la précision est de l'ordre de $\pm 0{,}5$ à $\pm 1$ unité de pH.
\end{enumerate}
\textbf{Avec le pH-mètre (mesure précise) :}
\begin{enumerate}
\item \textbf{Étalonner} l'appareil : plonger l'électrode de verre successivement dans deux \cle{solutions tampons} étalons (par exemple pH = 7 puis pH = 4) et régler l'appareil pour qu'il affiche ces valeurs. Cette étape est \textbf{indispensable} avant toute mesure.
\item Rincer l'électrode à l'eau distillée et l'essuyer délicatement.
\item Plonger l'électrode dans la solution à étudier, agiter doucement, attendre la stabilisation de l'affichage.
\item Lire le pH : précision de $\pm 0{,}05$ voire $\pm 0{,}01$ unité.
\end{enumerate}
Le pH-mètre mesure en réalité une tension électrique entre l'électrode de verre et une électrode de référence, tension proportionnelle au pH.
\end{methode}

\begin{attention}[Bonnes pratiques et pièges de mesure]
\begin{itemize}
\item Un pH-mètre non étalonné donne des valeurs \textbf{fausses} : l'étalonnage n'est pas optionnel.
\item La mesure dépend de la température : les pH-mètres modernes possèdent une sonde de température de compensation.
\item Le papier pH ne doit pas être mouillé à l'avance avec de l'eau distillée : cela diluerait la goutte testée et fausserait la mesure.
\item Les solutions très colorées (jus de bissap, vin de palme) rendent la lecture au papier pH délicate : préférer le pH-mètre.
\end{itemize}
\end{attention}

\courssub{Les indicateurs colorés acido-basiques}

Un \cle{indicateur coloré} acido-basique est un couple acide/base dont la forme acide et la forme basique ont des \textbf{couleurs différentes}. Sa teinte dépend donc du pH de la solution dans laquelle on le verse. On l'utilise pour estimer le pH ou, surtout, pour repérer l'équivalence lors d'un dosage (leçons suivantes).

\begin{definition}[Zone de virage]
La \cle{zone de virage} d'un indicateur coloré est l'intervalle de pH (en général d'environ 2 unités, centré sur le $\mathrm{p}K_a$ du couple indicateur) dans lequel sa couleur change progressivement de la teinte acide à la teinte basique. En dehors de cette zone, l'indicateur prend franchement l'une des deux teintes.
\end{definition}

\begin{center}
\begin{tabularx}{\linewidth}{|l|X|X|X|}
\hline
\rowcolor{popblueL} \textbf{Indicateur} & \textbf{Teinte acide} & \textbf{Zone de virage} & \textbf{Teinte basique} \\
\hline
Hélianthine (méthylorange) & rouge & 3,1 -- 4,4 (orange) & jaune \\
\hline
Bleu de bromothymol (BBT) & jaune & 6,0 -- 7,6 (vert) & bleu \\
\hline
Phénolphtaléine & incolore & 8,2 -- 10,0 (rose pâle) & rose fuchsia \\
\hline
\end{tabularx}
\end{center}

\textbf{Comment choisir un indicateur coloré ?} On choisit l'indicateur dont la \emph{zone de virage encadre le pH attendu} de la solution étudiée. Par exemple :
\begin{itemize}
\item pour vérifier qu'un jus de citron est acide (pH $\approx 2$), l'hélianthine vire au rouge : elle convient ;
\item pour tester la neutralité d'une eau de forage (pH $\approx 7$), seul le BBT, dont la zone de virage contient 7, est adapté : il donnera une teinte verte ;
\item pour caractériser une solution de savon (pH $\approx 9$--10), la phénolphtaléine vire au rose.
\end{itemize}

\begin{exoresolu}[Choisir le bon indicateur]
Énoncé. On dispose de trois solutions A, B, C de pH respectifs 3, 7 et 10 (à \SI{25}{\celsius}). On y ajoute quelques gouttes de bleu de bromothymol. Quelle teinte prend chaque solution ? Peut-on distinguer A de B avec l'hélianthine ?
\tcblower
Solution détaillée. La zone de virage du BBT est 6,0 -- 7,6 (jaune en milieu acide, bleu en milieu basique, vert dans la zone).
\begin{itemize}
\item Solution A (pH = 3 $<$ 6) : le BBT est sous sa forme acide $\Rightarrow$ teinte \textbf{jaune}.
\item Solution B (pH = 7, dans la zone de virage) : teinte \textbf{verte}.
\item Solution C (pH = 10 $>$ 7,6) : le BBT est sous sa forme basique $\Rightarrow$ teinte \textbf{bleue}.
\end{itemize}
Avec l'hélianthine (zone de virage 3,1 -- 4,4) : la solution A (pH = 3) est juste en dessous de la zone, donc rouge ; la solution B (pH = 7) est bien au-dessus, donc jaune. On \textbf{peut} distinguer A de B avec l'hélianthine. En revanche, elle ne distinguerait pas B de C (toutes deux jaunes, pH $>$ 4,4). Conclusion : on choisit toujours l'indicateur dont la zone de virage contient le pH à mettre en évidence.
\end{exoresolu}

\courssub{Application : lire un résultat d'analyse médicale}

Le pH n'est pas qu'une affaire de laboratoire de chimie : c'est une \textbf{constante biologique surveillée en médecine}. Le sang humain est maintenu par des systèmes tampons dans une fourchette très étroite.

\begin{definition}[pH sanguin et déséquilibres acido-basiques]
Le pH du sang artériel normal est compris entre \textbf{7,35 et 7,45} (sang légèrement basique). On parle de :
\begin{itemize}
\item \cle{acidose} si $\mathrm{pH} < 7{,}35$ (excès d'ions \ce{H3O+} : diabète décompensé, insuffisance respiratoire, effort intense...) ;
\item \cle{alcalose} si $\mathrm{pH} > 7{,}45$ (déficit en ions \ce{H3O+} : hyperventilation, vomissements répétés...).
\end{itemize}
Un pH sanguin inférieur à 7,0 ou supérieur à 7,8 met la vie en danger : d'où l'importance des appareils de mesure précis dans les hôpitaux.
\end{definition}

\begin{exemplebox}[Interpréter un bulletin d'analyse]
Un patient reçoit à l'hôpital de district un bulletin indiquant : « gazométrie sanguine : pH = 7,29 ». Comme $7{,}29 < 7{,}35$, le patient est en \textbf{acidose} (légère). En concentration : $[\ce{H3O+}] = 10^{-7{,}29} \approx 5{,}1\times 10^{-8}\ \si{mol.L^{-1}}$, contre $4{,}0\times 10^{-8}\ \si{mol.L^{-1}}$ à pH 7,40. Un écart de pH de 0,11 unité correspond déjà à une augmentation de près de 30 \% de la concentration en ions oxonium : voilà pourquoi le corps tolère si mal les variations de pH.
\end{exemplebox}

\begin{savaistu}[Du jus d'orange aux hôpitaux de Yaoundé et Douala]
Le premier pH-mètre électronique commercial a été mis au point en 1934 par Arnold Beckman, un chimiste américain... pour répondre à une demande de l'industrie californienne du \textbf{jus d'orange}, qui avait besoin de contrôler précisément l'acidité de ses fruits ! L'invention a connu un tel succès qu'elle a financé toute une industrie des instruments de mesure. Aujourd'hui, les pH-mètres et gazomètres dérivés de cet appareil équipent les laboratoires d'analyses médicales de Yaoundé et de Douala, où ils permettent chaque jour de mesurer le pH sanguin des patients, de contrôler la qualité de l'eau potable distribuée par la Camwater, ou de suivre la fermentation du vin de palme et de la bière dans les unités de production locales.
\end{savaistu}

\begin{aretenir}[L'essentiel de la section]
\begin{itemize}
\item $\mathrm{pH} = -\log[\ce{H3O+}]$ et $[\ce{H3O+}] = 10^{-\mathrm{pH}}$ (solutions diluées, $[\ce{H3O+}]$ en \si{mol.L^{-1}}).
\item À \SI{25}{\celsius} : acide si pH $< 7$, neutre si pH $= 7$, basique si pH $> 7$. Le pH neutre vaut $\mathrm{p}K_e/2$ et dépend de la température.
\item Multiplier $[\ce{H3O+}]$ par 10 fait \textbf{baisser} le pH de 1 unité.
\item L'échelle 0--14 est usuelle : un pH peut être négatif ou supérieur à 14 pour des solutions concentrées.
\item Papier pH : rapide, précision $\pm 0{,}5$ à $\pm 1$. pH-mètre : précis ($\pm 0{,}05$), mais exige un \textbf{étalonnage} avec des solutions tampons.
\item Indicateurs colorés : hélianthine (3,1--4,4), BBT (6,0--7,6), phénolphtaléine (8,2--10,0) ; on choisit celui dont la zone de virage contient le pH attendu.
\item pH sanguin normal : 7,35--7,45 ; acidose en dessous, alcalose au-dessus.
\end{itemize}
\end{aretenir}

\begin{exercice}[À toi de jouer]
\begin{enumerate}
\item Le jus de bissap préparé par ta voisine a un pH de 3,2 à \SI{25}{\celsius}. Calcule $[\ce{H3O+}]$ et $[\ce{HO-}]$ dans ce jus. Est-il acide, neutre ou basique ?
\item On dilue 10 fois une solution de pH = 4 dont l'acidité est due à un acide fort. Quel est le pH de la solution diluée ?
\item Quel indicateur coloré choisirais-tu pour distinguer une eau de forage (pH $\approx 7$) d'une eau de lessive (pH $\approx 12$) ? Justifie.
\end{enumerate}
\end{exercice}

\begin{corrige}[Exercice : À toi de jouer]
\begin{enumerate}
\item $[\ce{H3O+}] = 10^{-3{,}2} \approx 6{,}3\times 10^{-4}\ \si{mol.L^{-1}}$. D'après le produit ionique : $[\ce{HO-}] = \dfrac{K_e}{[\ce{H3O+}]} = \dfrac{10^{-14}}{6{,}3\times 10^{-4}} \approx 1{,}6\times 10^{-11}\ \si{mol.L^{-1}}$. Comme pH $= 3{,}2 < 7$ (ou $[\ce{H3O+}] > [\ce{HO-}]$), le jus est \textbf{acide}.
\item Diluer 10 fois un acide fort divise $[\ce{H3O+}]$ par 10, donc le pH augmente de 1 : pH $= 5$.
\item Le BBT : sa zone de virage (6,0--7,6) contient le pH de l'eau de forage, qui apparaîtra verte, tandis que l'eau de lessive (pH 12 $>$ 7,6) le fera virer au bleu. La phénolphtaléine conviendrait aussi (incolore à pH 7, rose à pH 12), mais pas l'hélianthine, jaune dans les deux cas.
\end{enumerate}
\end{corrige}

\coursec{Calculs de pH et concentrations ioniques}

Dans la section précédente, tu as appris à mesurer le pH avec un pH-mètre ou du papier pH. Mais un chimiste digne de ce nom ne se contente pas de \emph{mesurer} : il doit être capable de \emph{prévoir}. Si un technicien de la SONARA prépare une solution d'acide chlorhydrique à \SI{2,0e-3}{mol.L^{-1}} pour détartrer une canalisation, quel pH obtiendra-t-il ? Si une laborantine d'un hôpital de Yaoundé dilue dix fois une solution de soude, comment le pH va-t-il évoluer ? C'est toute l'ambition de cette section : à partir de la seule concentration apportée en soluté, calculer le pH et les concentrations de toutes les espèces présentes, en toute rigueur, grâce à deux outils que tu vas apprendre à manier comme un professionnel : l'\cle{électroneutralité} et les \cle{approximations justifiées}.

\courssub{Les relations maîtresses}

Tout ce qui suit repose sur trois relations que tu dois connaître par c\oe ur, car elles sont la clé de tous les exercices du baccalauréat TI.

\begin{aretenir}[Les trois relations fondamentales]
À \SI{25}{\celsius}, pour toute solution aqueuse :
\[ \text{pH} = -\log\,[\ce{H3O+}] \qquad \Longleftrightarrow \qquad [\ce{H3O+}] = 10^{-\text{pH}} \]
\[ [\ce{HO-}] = \frac{K_e}{[\ce{H3O+}]} = 10^{\text{pH} - \text{p}K_e} \qquad \text{avec } \text{p}K_e = 14{,}0 \]
Ces relations sont \textbf{toujours vraies}, quelle que soit la solution (acide, base, sel, eau pure).
\end{aretenir}

L'idée intuitive est simple : le pH est une \og échelle compressée \fg{} de la concentration en ions oxonium. Comme $[\ce{H3O+}]$ varie typiquement de \SI{10}{mol.L^{-1}} à \SI{e-15}{mol.L^{-1}}, le logarithme la ramène à une échelle confortable de 0 à 14. Retenir le sens de variation : \textbf{plus la concentration en \ce{H3O+} est grande, plus le pH est petit}.

\begin{attention}[Erreur fréquente : le passage pH $\to$ concentration]
Pour retrouver la concentration à partir du pH, on utilise $[\ce{H3O+}] = 10^{-\text{pH}}$, et \textbf{surtout pas} $-\text{pH}$ ni $10^{\text{pH}}$. Exemple : si pH $= 3{,}0$, alors $[\ce{H3O+}] = 10^{-3} = \SI{1,0e-3}{mol.L^{-1}}$. Sur ta calculatrice, tape \texttt{10\^{}(-3)} ou la touche $10^x$. Une erreur de signe ici fausse tout l'exercice.
\end{attention}

\courssub{Électroneutralité et bilan des espèces}

\courssub{Le bilan des espèces : qui est présent dans le bécher ?}

Avant tout calcul, le chimiste fait l'\emph{inventaire} de ce qui nage réellement dans la solution. Une solution aqueuse est électriquement neutre : elle contient donc forcément des cations et des anions.

\begin{definition}[Électroneutralité]
Toute solution aqueuse est électriquement neutre : la somme des concentrations des charges positives est égale à la somme des concentrations des charges négatives, chaque concentration étant multipliée par la valeur absolue de la charge de l'ion.
\end{definition}

Prenons l'exemple d'une solution de chlorure de sodium et d'acide chlorhydrique. Les espèces présentes sont : \ce{Na+}, \ce{Cl-}, \ce{H3O+} et \ce{HO-} (ces deux derniers étant toujours présents dans l'eau, à cause de l'autoprotolyse). L'électroneutralité s'écrit :
\[ [\ce{Na+}] + [\ce{H3O+}] = [\ce{Cl-}] + [\ce{HO-}] \]

\begin{attention}[Piège classique]
N'oublie jamais les ions \ce{H3O+} et \ce{HO-} dans le bilan des espèces : ils sont présents dans \textbf{toute} solution aqueuse, même s'ils ne sont pas apportés par le soluté. Oublier \ce{HO-} dans l'électroneutralité d'une solution acide, c'est se priver de la relation qui permet de vérifier la cohérence du calcul.
\end{attention}

\courssub{Les approximations : majoritaire, minoritaire, ultra-minoritaire}

Voici l'outil intellectuel le plus précieux de cette leçon. Dans une solution, les concentrations des espèces sont souvent très différentes : il est alors légitime de \emph{négliger} les plus petites devant les plus grandes pour simplifier les équations.

\begin{definition}[Hiérarchie des espèces]
\begin{itemize}
\item Une espèce est dite \cle{majoritaire} si sa concentration est la plus élevée (ou parmi les plus élevées) de la solution.
\item Une espèce est dite \cle{minoritaire} si sa concentration est nettement plus faible que celle de l'espèce majoritaire, mais pas négligeable.
\item Une espèce est dite \cle{ultra-minoritaire} si sa concentration est négligeable devant celle de l'espèce de référence.
\end{itemize}
\end{definition}

\begin{propriete}[Règle pratique de négligence]
Une espèce A est négligeable devant une espèce B si :
\[ [\text{A}] \leq \frac{[\text{B}]}{10} \]
c'est-à-dire si sa concentration est au moins \textbf{dix fois plus faible} (un ordre de grandeur d'écart). On peut alors supprimer A des bilans de matière et de l'électroneutralité.
\end{propriete}

L'intuition : si tu partages un sac de riz entre 1000 personnes et que l'une d'elles ne reçoit qu'un grain, tu peux ignorer ce grain dans tes comptes. En chimie, ignorer une espèce ultra-minoritaire transforme une équation compliquée en une formule immédiate. Attention toutefois : \textbf{toute approximation doit être vérifiée a posteriori}. On calcule, puis on contrôle que l'espèce négligée est bien au moins dix fois moins concentrée.

\courssub{Cas 1 : solution d'acide fort}

\begin{exemplebox}[Dissolution du chlorure d'hydrogène]
Le chlorure d'hydrogène \ce{HCl} est un acide fort : sa réaction avec l'eau est \textbf{totale} :
\[ \ce{HCl + H2O -> H3O+ + Cl-} \]
(réaction totale, rapide, exothermique ; l'eau joue le rôle de base). Une solution de concentration apportée $C$ contient donc $[\ce{H3O+}]_{\text{apportés}} = C$ et $[\ce{Cl-}] = C$.
\end{exemplebox}

Appliquons la démarche rigoureuse. Espèces présentes : \ce{H3O+}, \ce{Cl-}, \ce{HO-}. Électroneutralité :
\[ [\ce{H3O+}] = [\ce{Cl-}] + [\ce{HO-}] = C + [\ce{HO-}] \]
Si la solution n'est pas trop diluée ($C \gg 10^{-7}$), les ions \ce{HO-} issus de l'autoprotolyse sont ultra-minoritaires : $[\ce{HO-}] \ll C$. Il reste $[\ce{H3O+}] \simeq C$, d'où :

\begin{aretenir}[pH d'un acide fort]
Pour une solution d'acide fort de concentration $C$ (avec $10^{-6} \lesssim C \lesssim \SI{1}{mol.L^{-1}}$) :
\[ \boxed{\text{pH} = -\log C} \]
Vérification de cohérence obligatoire : à \SI{25}{\celsius}, le pH d'un acide doit être \textbf{inférieur à 7}.
\end{aretenir}

\begin{exemplebox}[Acide chlorhydrique à \SI{2,0e-3}{mol.L^{-1}}]
$C = \SI{2,0e-3}{mol.L^{-1}}$.
\[ \text{pH} = -\log(2{,}0 \times 10^{-3}) = 3 - \log 2{,}0 = 2{,}7 \]
\[ [\ce{HO-}] = \frac{K_e}{[\ce{H3O+}]} = \frac{10^{-14}}{2{,}0 \times 10^{-3}} = \SI{5,0e-12}{mol.L^{-1}} \]
Vérifications : pH $= 2{,}7 < 7$ (cohérent pour un acide) ; $[\ce{HO-}] = \SI{5,0e-12}{} \ll \SI{2,0e-3}{}$ : l'approximation est excellente (huit ordres de grandeur d'écart !). Bilan : \ce{Cl-} et \ce{H3O+} majoritaires (\SI{2,0e-3}{mol.L^{-1}}), \ce{HO-} ultra-minoritaire.
\end{exemplebox}

\begin{popfigure}
\begin{center}
\begin{tikzpicture}
\begin{semilogyaxis}[
  width=0.85\linewidth, height=6cm,
  ybar, bar width=18pt,
  xtick={1,2,3},
  xticklabels={$\ce{Cl-}$,$\ce{H3O+}$,$\ce{HO-}$},
  ymin=1e-13, ymax=1e-1,
  ylabel={Concentration (mol.L$^{-1}$)},
  ylabel style={font=\small},
  x tick label style={font=\small},
  ymajorgrids=true, grid style={popline!40},
  every axis plot/.append style={fill=popblue!70, draw=popblue},
  nodes near coords={
    \pgfmathprintnumber[sci,sci precision=1]{\pgfplotspointmeta}
  },
  every node near coord/.append style={font=\scriptsize, rotate=90, anchor=west, yshift=2pt},
]
\addplot coordinates {(1,2e-3) (2,2e-3) (3,5e-12)};
\end{semilogyaxis}
\end{tikzpicture}
\end{center}
\legende{Concentrations des espèces dans une solution de \ce{HCl} à $C = \SI{2,0e-3}{mol.L^{-1}}$ (échelle logarithmique) : \ce{Cl-} et \ce{H3O+} sont majoritaires, \ce{HO-} est ultra-minoritaire, invisible à l'échelle ordinaire.}
\end{popfigure}

\courssub{Cas 2 : solution de base forte}

\begin{exemplebox}[Dissolution de l'hydroxyde de sodium]
La soude \ce{NaOH} est une base forte : sa dissolution dans l'eau est \textbf{totale} :
\[ \ce{NaOH ->[eau] Na+ + HO-} \]
(dissociation totale, exothermique : c'est pourquoi l'eau chauffe quand on y dissout de la soude, comme le savent les savonnières artisanales qui préparent la lessive de soude pour saponifier l'huile de palme). Une solution de concentration $C$ contient $[\ce{Na+}] = C$ et $[\ce{HO-}]_{\text{apportés}} = C$.
\end{exemplebox}

Électroneutralité : $[\ce{Na+}] + [\ce{H3O+}] = [\ce{HO-}]$, soit $[\ce{HO-}] = C + [\ce{H3O+}]$. Si la solution n'est pas trop diluée, $[\ce{H3O+}] \ll C$, donc $[\ce{HO-}] \simeq C$. Alors :
\[ [\ce{H3O+}] = \frac{K_e}{[\ce{HO-}]} = \frac{K_e}{C} \quad\Longrightarrow\quad \text{pH} = -\log\frac{K_e}{C} = \text{p}K_e + \log C \]

\begin{aretenir}[pH d'une base forte]
Pour une solution de base forte de concentration $C$ (avec $10^{-6} \lesssim C \lesssim \SI{1}{mol.L^{-1}}$), à \SI{25}{\celsius} :
\[ \boxed{\text{pH} = \text{p}K_e + \log C = 14 + \log C} \]
Vérification de cohérence : le pH d'une base doit être \textbf{supérieur à 7}.
\end{aretenir}

\begin{exemplebox}[Soude à \SI{1,0e-2}{mol.L^{-1}}]
$C = \SI{1,0e-2}{mol.L^{-1}}$.
\[ \text{pH} = 14 + \log(10^{-2}) = 14 - 2 = 12{,}0 \]
\[ [\ce{H3O+}] = 10^{-12} = \SI{1,0e-12}{mol.L^{-1}} \]
Vérifications : pH $= 12 > 7$ (cohérent pour une base) ; $[\ce{H3O+}] \ll C$ : approximation validée. Bilan : \ce{Na+} et \ce{HO-} majoritaires, \ce{H3O+} ultra-minoritaire.
\end{exemplebox}

\courssub{Effet de la dilution sur le pH}

\begin{propriete}[Dilution d'un acide ou d'une base forte]
Diluer $10$ fois une solution d'acide fort (concentration divisée par 10) \textbf{augmente} son pH de $1$ unité :
\[ \text{pH}' = -\log\frac{C}{10} = -\log C + 1 = \text{pH} + 1 \]
De même, diluer $10$ fois une solution de base forte \textbf{diminue} son pH de $1$ unité. La dilution rapproche toujours le pH de $7$, sans jamais pouvoir le franchir.
\end{propriete}

C'est logique : diluer, c'est diminuer la concentration en \ce{H3O+} (pour un acide), donc augmenter le pH. Mais on ne pourra jamais rendre une solution acide basique en ajoutant de l'eau : quand $C$ s'approche de $10^{-7}$, l'autoprotolyse de l'eau reprend ses droits et le pH plafonne vers 7.

\begin{attention}[Erreur fréquente : appliquer les formules hors de leur domaine]
La formule pH $= -\log C$ n'est valable que pour un \textbf{acide fort} en solution \textbf{suffisamment concentrée} ($C \gtrsim \SI{e-6}{mol.L^{-1}}$).
\begin{itemize}
\item Pour un \textbf{acide faible}, la réaction avec l'eau est partielle : $[\ce{H3O+}] < C$ et pH $> -\log C$. Appliquer pH $= -\log C$ à l'acide éthanoïque du vinaigre est faux.
\item Pour une solution \textbf{très diluée} ($C$ proche de $10^{-7}$), les \ce{H3O+} de l'autoprotolyse ne sont plus négligeables : une solution de \ce{HCl} à \SI{e-8}{mol.L^{-1}} n'a \textbf{pas} un pH de 8 (une solution acide ne peut pas être basique !) ; son pH réel est voisin de 6,98.
\end{itemize}
\end{attention}

\courssub{Préparation d'une solution par dilution : mode opératoire (savoir-faire)}

Le savoir-faire \og Préparer une solution diluée à partir d'une solution concentrée \fg{} est évalué au Bac TI. La relation $C_m V_m = C_f V_f$ (conservation de la quantité de matière) guide le calcul, mais la rigueur expérimentale garantit le résultat.

\begin{methode}[Préparer une solution fille par dilution d'une solution mère]
\textbf{Matériel :} pipette jaugée (ou graduée selon la précision demandée), fiole jaugée de volume $V_f$, bécher, entonnoir, eau distillée, baguette de verre.
\begin{enumerate}
\item \textbf{Calculer} le volume $V_m$ de solution mère à prélever : $C_m V_m = C_f V_f$.
\item \textbf{Rincer} la pipette avec la solution mère (2-3 fois, en faisant tourner pour mouiller toutes les parois).
\item \textbf{Prélever} $V_m$ de solution mère au-dessus de la jauge, ajuster le ménisque au trait de jauge (œil au niveau du ménisque), laisser s'écouler dans la fiole jaugée (ne pas souffler la dernière goutte, la pipette est étalonnée \og à livrer \fg{}).
\item \textbf{Rincer} le bécher et l'entonnoir utilisés pour le transvasement avec un peu d'eau distillée, verser les rinçages dans la fiole (récupération quantitative).
\item \textbf{Compléter} au trait de jauge avec de l'eau distillée (au compte-gouttes près de la jauge), boucher, homogénéiser énergiquement (retourner la fiole plusieurs fois).
\item \textbf{Étiqueter} la fiole (nature du soluté, concentration, date, nom).
\end{enumerate}
\end{methode}

\begin{attention}[Pièges au laboratoire]
\begin{itemize}
\item Ne jamais verser l'eau dans l'acide concentré (éclaboussures dangereuses par échauffement brutal) : on verse la solution mère \emph{dans} l'eau. Pour les acides concentrés du commerce (ex. \ce{HCl} \SI{12}{mol.L^{-1}}, \ce{H2SO4} \SI{18}{mol.L^{-1}}), on met d'abord \SI{\sim 50}{\%} d'eau dans la fiole, puis l'acide, puis on complète.
\item La pipette jaugée ne se chauffe pas, ne passe pas au four (dilatation = perte de justesse).
\item La fiole jaugée ne sert qu'à préparer la solution, pas à la stocker longtemps (transvaser dans un flacon en PEHD bouché).
\end{itemize}
\end{attention}

\courssub{Organigramme de décision et méthode générale}

\begin{popfigure}
\begin{center}
\begin{tikzpicture}[
  node distance=0.7cm and 1.1cm,
  startstop/.style={rectangle, rounded corners, draw=popdark, fill=popblueL, text width=4.2cm, align=center, font=\small},
  decision/.style={diamond, draw=popdark, fill=popgoldL, aspect=2.2, align=center, font=\small, inner sep=1pt},
  process/.style={rectangle, draw=popdark, fill=popgreenL, text width=3.6cm, align=center, font=\small},
  warn/.style={rectangle, draw=popdark, fill=poppinkL, text width=3.4cm, align=center, font=\small},
  arrow/.style={-{Stealth[length=2.5mm]}, thick, popdark}
]
\node[startstop] (start) {Solution de concentration apportée $C$ connue};
\node[decision, below=of start] (q1) {Acide fort ?};
\node[process, below left=0.9cm and -0.4cm of q1] (acide) {$[\ce{H3O+}] = C$\\[2pt] $\text{pH} = -\log C$};
\node[decision, below right=0.9cm and -0.6cm of q1] (q2) {Base forte ?};
\node[process, below=of q2] (base) {$[\ce{HO-}] = C$\\[2pt] $\text{pH} = 14 + \log C$};
\node[warn, right=1.0cm of q2] (faible) {Acide ou base \textbf{faible} :\\ réaction partielle,\\ voir leçon sur les $K_A$};
\node[warn, below=1.6cm of acide] (verif) {Vérifier : pH $< 7$ (acide) ou pH $> 7$ (base) ;\\ $C \gtrsim \SI{e-6}{mol.L^{-1}}$};
\draw[arrow] (start) -- (q1);
\draw[arrow] (q1) -| node[above left, font=\scriptsize]{oui} (acide);
\draw[arrow] (q1) -| node[above right, font=\scriptsize]{non} (q2);
\draw[arrow] (q2) -- node[right, font=\scriptsize]{oui} (base);
\draw[arrow] (q2) -- node[above, font=\scriptsize]{non} (faible);
\draw[arrow] (acide) -- (verif);
\draw[arrow] (base.south) |- (verif.east);
\end{tikzpicture}
\end{center}
\legende{Organigramme de décision : quelle formule de pH appliquer ? Toujours terminer par la vérification de cohérence.}
\end{popfigure}

\begin{methode}[Calculer le pH d'une solution en quatre étapes]
\begin{enumerate}
\item \textbf{Identifier les espèces} présentes en solution (ions du soluté, plus \ce{H3O+} et \ce{HO-} qui sont toujours là) et écrire les équations de mise en solution (totales pour acide/base forts).
\item \textbf{Écrire l'électroneutralité} de la solution et le produit ionique de l'eau $K_e = [\ce{H3O+}][\ce{HO-}]$.
\item \textbf{Faire les approximations justifiées} : négliger les espèces ultra-minoritaires (au moins 10 fois moins concentrées) pour simplifier l'électroneutralité.
\item \textbf{Calculer puis vérifier} : en déduire $[\ce{H3O+}]$, le pH, puis contrôler a posteriori que chaque approximation est valide et que le pH est cohérent ($< 7$ pour un acide, $> 7$ pour une base à \SI{25}{\celsius}).
\end{enumerate}
\end{methode}

\begin{exoresolu}[pH et bilan des espèces d'une solution d'acide chlorhydrique]
\textbf{Énoncé.} On prépare, à \SI{25}{\celsius}, un volume $V = \SI{500}{mL}$ d'une solution S d'acide chlorhydrique de concentration $C = \SI{5,0e-4}{mol.L^{-1}}$. On donne $\text{p}K_e = 14{,}0$.
\begin{enumerate}
\item Faire le bilan de toutes les espèces chimiques présentes dans S.
\item Écrire l'équation d'électroneutralité et, après approximation justifiée, calculer le pH de S.
\item Calculer $[\ce{HO-}]$ et conclure sur le caractère majoritaire ou ultra-minoritaire de chaque espèce.
\item On dilue S dix fois. Quel est le pH de la solution S$'$ obtenue ?
\end{enumerate}
\tcblower
\textbf{Solution détaillée.}
\begin{enumerate}
\item L'acide chlorhydrique est un acide fort, sa réaction avec l'eau est totale :
\[ \ce{HCl + H2O -> H3O+ + Cl-} \]
Espèces présentes : \ce{H3O+}, \ce{Cl-} (apportées, en concentration $C$ chacune) et \ce{HO-} (issue de l'autoprotolyse de l'eau : $\ce{2 H2O <=> H3O+ + HO-}$).
\item Électroneutralité : $[\ce{H3O+}] = [\ce{Cl-}] + [\ce{HO-}] = C + [\ce{HO-}]$.\\
Approximation : la solution est acide et $C = \SI{5,0e-4}{} \gg 10^{-7}$, donc $[\ce{HO-}] \ll C$ est attendu. On néglige $[\ce{HO-}]$ :
\[ [\ce{H3O+}] \simeq C = \SI{5,0e-4}{mol.L^{-1}} \quad\Longrightarrow\quad \text{pH} = -\log(5{,}0 \times 10^{-4}) = 4 - \log 5{,}0 \approx 3{,}3 \]
\item $[\ce{HO-}] = \dfrac{K_e}{[\ce{H3O+}]} = \dfrac{10^{-14}}{5{,}0 \times 10^{-4}} = \SI{2,0e-11}{mol.L^{-1}}$.\\
Vérification : $\SI{2,0e-11}{} \ll \SI{5,0e-4}{}$ (sept ordres de grandeur) : l'approximation est validée. Le pH $= 3{,}3 < 7$ : cohérent pour un acide.\\
Bilan : \ce{H3O+} et \ce{Cl-} sont \textbf{majoritaires} (\SI{5,0e-4}{mol.L^{-1}}) ; \ce{HO-} est \textbf{ultra-minoritaire}.
\item Dilution dix fois : $C' = \SI{5,0e-5}{mol.L^{-1}}$, donc pH$' = -\log(5{,}0 \times 10^{-5}) = 5 - \log 5{,}0 \approx 4{,}3$. Le pH a augmenté d'une unité, comme prévu pour un acide fort dilué dix fois.
\end{enumerate}
\end{exoresolu}

\begin{savaistu}[Le pH dans les analyses médicales]
Au Cameroun comme partout, un bilan sanguin indique le pH du sang : il doit rester entre $7{,}35$ et $7{,}45$. En dessous, on parle d'\emph{acidose} ; au-dessus, d'\emph{alcalose}. Ces valeurs correspondent à des concentrations $[\ce{H3O+}]$ comprises entre $10^{-7{,}35} \approx \SI{4,5e-8}{mol.L^{-1}}$ et $10^{-7{,}45} \approx \SI{3,5e-8}{mol.L^{-1}}$ : des variations infimes, que l'organisme régule grâce à des systèmes \emph{tampons} que tu étudieras bientôt. Le calcul de concentrations à partir du pH, que tu viens de maîtriser, est exactement celui qu'utilisent les biologistes médicaux pour interpréter ces résultats.
\end{savaistu}

\begin{exercice}[À toi de jouer]
À \SI{25}{\celsius}, on dissout \SI{0,40}{g} d'hydroxyde de sodium \ce{NaOH} ($M = \SI{40,0}{g.mol^{-1}}$) dans de l'eau distillée pour obtenir \SI{1,0}{L} de solution. La dissolution est supposée totale.
\begin{enumerate}
\item Calculer la concentration $C$ de la solution.
\item Faire le bilan des espèces présentes et écrire l'électroneutralité.
\item Calculer le pH de la solution, puis $[\ce{H3O+}]$. Vérifier la cohérence.
\item On prélève \SI{10,0}{mL} de cette solution que l'on dilue pour obtenir \SI{100,0}{mL}. Quel est le pH de la solution diluée ?
\end{enumerate}
\end{exercice}

\begin{corrige}[Exercice : solution de soude]
\begin{enumerate}
\item $n = \dfrac{m}{M} = \dfrac{0{,}40}{40{,}0} = \SI{1,0e-2}{mol}$ ; $C = \dfrac{n}{V} = \SI{1,0e-2}{mol.L^{-1}}$.
\item Dissolution totale : $\ce{NaOH ->[eau] Na+ + HO-}$. Espèces : \ce{Na+}, \ce{HO-} (majoritaires, $C$ chacune) et \ce{H3O+}. Électroneutralité : $[\ce{Na+}] + [\ce{H3O+}] = [\ce{HO-}]$.
\item On néglige $[\ce{H3O+}]$ devant $C$ : $[\ce{HO-}] \simeq C$. pH $= 14 + \log(10^{-2}) = 12{,}0$ ; $[\ce{H3O+}] = 10^{-12} = \SI{1,0e-12}{mol.L^{-1}} \ll C$ : approximation validée ; pH $> 7$ : cohérent pour une base.
\item Facteur de dilution $F = 100{,}0/10{,}0 = 10$ : $C' = \SI{1,0e-3}{mol.L^{-1}}$, pH$' = 14 + \log(10^{-3}) = 11{,}0$. Le pH a diminué d'une unité.
\end{enumerate}
\end{corrige}

\begin{aretenir}[L'essentiel de la section]
\begin{itemize}
\item Relations maîtresses : pH $= -\log[\ce{H3O+}]$, $[\ce{H3O+}] = 10^{-\text{pH}}$, $[\ce{HO-}] = 10^{\text{pH}-\text{p}K_e}$.
\item Acide fort (solution pas trop diluée) : pH $= -\log C$ ; base forte : pH $= 14 + \log C$.
\item Toute solution est électriquement neutre ; on néglige une espèce si sa concentration est au moins dix fois inférieure à celle de référence, et on vérifie toujours l'approximation a posteriori.
\item Diluer dix fois un acide fort augmente le pH de 1 ; diluer dix fois une base forte le diminue de 1 ; le pH ne franchit jamais 7 par simple dilution.
\item Contrôle final systématique : pH $< 7$ pour un acide, pH $> 7$ pour une base (à \SI{25}{\celsius}).
\item Savoir-faire : préparer une solution fille par dilution (calcul $C_m V_m = C_f V_f$, rinçage pipette/fiole, ajustement au trait de jauge, homogénéisation, étiquetage).
\end{itemize}
\end{aretenir}

\coursec{Préparation de solutions par dilution}

Dans la section précédente, tu as appris à calculer le pH d'une solution et à en déduire les concentrations des ions \ce{H3O+} et \ce{HO-}. Mais une question pratique demeure : comment le technicien du laboratoire d'analyses médicales de l'hôpital de district, ou le contrôleur qualité d'une savonnerie artisanale, obtient-il concrètement une solution de concentration précise ? Les solutions du commerce sont souvent très concentrées (l'acide chlorhydrique vendu pour déboucher les canalisations, par exemple). Il faut donc savoir les \cle{diluer} avec rigueur. C'est l'objet de cette dernière section, qui te donnera un savoir-faire directement mobilisable au travaux pratiques du Baccalauréat TI.

\courssub{Principe de la dilution}

\textbf{Intuition d'abord.} Imagine un verre de jus de bissap trop concentré, au goût presque agressif. Que fais-tu ? Tu ajoutes de l'eau. Le jus devient moins fort, mais la quantité d'extrait d'hibiscus présente dans le verre n'a pas changé d'un gramme : elle s'est simplement répartie dans un plus grand volume de liquide. La dilution en chimie repose exactement sur cette idée.

\begin{definition}[Dilution]
\cle{Diluer} une solution, c'est lui ajouter du \cle{solvant} (en général de l'eau distillée) \textbf{sans ajouter ni retirer de soluté}. La solution de départ est appelée \cle{solution mère} (concentration $C_0$, volume prélevé $V_0$) ; la solution obtenue est la \cle{solution fille} (concentration $C_f$, volume final $V_f$).
\end{definition}

Comme aucune quantité de soluté n'est ajoutée ni perdue, la quantité de matière de soluté se conserve :

\begin{propriete}[Relation de dilution]
Lors d'une dilution, la quantité de matière de soluté prélevée dans la solution mère est égale à celle contenue dans la solution fille :
\[ n_0 = n_f \quad \Longleftrightarrow \quad \boxed{C_0 \times V_0 = C_f \times V_f} \]
avec $C_0$ et $C_f$ en \si{mol.L^{-1}}, $V_0$ et $V_f$ dans la \textbf{même unité} (par exemple en \si{mL}).
\end{propriete}

\begin{attention}[Unités des volumes]
Il n'est pas obligatoire de convertir les volumes en litres, à condition que $V_0$ et $V_f$ soient exprimés \textbf{dans la même unité}. En revanche, mélanger \si{mL} et \si{L} dans la même formule est l'erreur classique qui fait perdre des points au Bac !
\end{attention}

\courssub{Le facteur de dilution}

Pour caractériser « combien de fois » on a dilué, on définit une grandeur sans unité :

\begin{definition}[Facteur de dilution]
Le \cle{facteur de dilution} $F$ est le rapport :
\[ F = \frac{C_0}{C_f} = \frac{V_f}{V_0} \qquad (F > 1) \]
On dit alors que la solution a été diluée « $F$ fois » ou « au $1/F$ ».
\end{definition}

\begin{exemplebox}[Lecture d'un facteur de dilution]
Diluer une solution « au dixième » signifie $F = 10$ : la concentration finale est dix fois plus faible, et le volume final est dix fois plus grand que le volume prélevé. Par exemple, prélever \SI{10,0}{mL} de solution mère et compléter à \SI{100,0}{mL} donne $F = \dfrac{100}{10} = 10$.
\end{exemplebox}

\courssub{Calcul du volume à prélever}

Le problème pratique type est le suivant : \emph{on dispose d'une solution mère de concentration $C_0$ ; quel volume $V_0$ faut-il prélever pour préparer un volume $V_f$ de solution fille de concentration $C_f$ ?} On isole $V_0$ dans la relation de dilution :

\[ V_0 = \frac{C_f \times V_f}{C_0} \]

\begin{exoresolu}[Préparation d'une solution d'acide chlorhydrique diluée]
Au laboratoire, on veut préparer $V_f = \SI{100,0}{mL}$ d'une solution d'acide chlorhydrique de concentration $C_f = \SI{0,010}{mol.L^{-1}}$ à partir d'une solution mère de concentration $C_0 = \SI{0,10}{mol.L^{-1}}$. Déterminer le volume $V_0$ à prélever, le facteur de dilution, puis le pH de la solution fille.
\tcblower
\textbf{1. Volume à prélever.} La conservation de la quantité de matière s'écrit $C_0 V_0 = C_f V_f$, d'où :
\[ V_0 = \frac{C_f \times V_f}{C_0} = \frac{0{,}010 \times 100{,}0}{0{,}10} = \SI{10,0}{mL} \]
Il faut donc prélever \SI{10,0}{mL} de solution mère à la pipette jaugée.

\textbf{2. Facteur de dilution.}
\[ F = \frac{C_0}{C_f} = \frac{0{,}10}{0{,}010} = 10 \]
La solution a été diluée 10 fois (au dixième).

\textbf{3. pH de la solution fille.} \ce{HCl} est un acide fort, totalement dissocié dans l'eau :
\[ \ce{HCl + H2O -> H3O+ + Cl-} \]
Donc $[\ce{H3O+}] = C_f = \SI{1,0e-2}{mol.L^{-1}}$ et :
\[ \text{pH} = -\log[\ce{H3O+}] = -\log(10^{-2}) = 2{,}0 \]
La solution fille a un pH de 2,0 (alors que la solution mère avait un pH de 1,0) : diluer un acide fort fait \textbf{augmenter} le pH.
\end{exoresolu}

\courssub{Protocole expérimental de la dilution}

La précision d'une dilution repose sur la \cle{verrerie jaugée} : la \cle{pipette jaugée} pour le prélèvement exact de $V_0$, et la \cle{fiole jaugée} pour l'ajustage exact du volume final $V_f$.

\begin{methode}[Préparer une solution diluée avec précision]
\begin{enumerate}
\item \textbf{Calculer} le volume à prélever : $V_0 = \dfrac{C_f \times V_f}{C_0}$.
\item \textbf{Prélever} $V_0$ de solution mère à l'aide d'une pipette jaugée (munie d'une poire aspirante — jamais à la bouche !). Si la pipette adaptée n'existe pas, utiliser une éprouvette graduée, moins précise.
\item \textbf{Introduire} le prélèvement dans une fiole jaugée de volume $V_f$ contenant déjà un peu d'eau distillée (environ un tiers).
\item \textbf{Compléter} avec de l'eau distillée jusqu'à ce que le bas du ménisque affleure le trait de jauge, en terminant à la pipette simple pour les dernières gouttes.
\item \textbf{Boucher et homogénéiser} en retournant la fiole une dizaine de fois. La solution fille est prête.
\end{enumerate}
\end{methode}

\begin{popfigure}
\begin{tikzpicture}[scale=0.9, every node/.style={font=\small}]
% ---- Styles ----
\tikzset{
  glass/.style={popline, thick},
  liquid/.style={popblue!40},
  liquidfinal/.style={popgreen!35},
  labelstyle/.style={popdark, align=center, font=\scriptsize},
  arrowstyle/.style={-{Stealth}, very thick, poporange},
}
% ---- 1. Pipette jaugée ----
\begin{scope}[xshift=0cm]
% Glass walls
\draw[glass] (0,0) -- (0,4);
\draw[glass] (0.5,0) -- (0.5,4);
% Tip
\draw[glass] (0,0) -- (0.25,-0.6) -- (0.5,0);
% Liquid
\fill[liquid] (0.05,0.05) rectangle (0.45,2.8);
% Jauge mark
\draw[glass] (0.05,2.8) -- (0.45,2.8) node[midway, right=3pt, popdark]{\scriptsize trait de jauge};
% Label
\node[labelstyle] at (0.25,-1.3) {Pipette jaugée\\ $V_0 = \SI{10,0}{mL}$};
\end{scope}

% ---- Transfert arrow ----
\draw[arrowstyle] (1.4,1.8) .. controls (2.2,2.8) and (3.4,2.8) .. (4.2,1.8);
\node[poporange, font=\scriptsize\bfseries] at (2.8,3.1) {Transfert};

% ---- 2. Fiole jaugée (réception + ajustage) ----
\begin{scope}[xshift=5.5cm]
% Flask outline
\draw[glass] (-1.5,0) .. controls (-1.5,1) and (-0.8,1.5) .. (-0.3,2.0); % left bottom
\draw[glass] (1.5,0) .. controls (1.5,1) and (0.8,1.5) .. (0.3,2.0);   % right bottom
\draw[glass] (-0.3,2.0) -- (-0.3,3.5); % neck left
\draw[glass] (0.3,2.0) -- (0.3,3.5);   % neck right
% Jauge mark
\draw[glass] (-0.3,3.2) -- (0.3,3.2) node[midway, right=3pt, popdark]{\scriptsize trait de jauge};
% Meniscus line (adjustment)
\draw[popline, dashed] (-0.5,3.2) .. controls (0,3.25) .. (0.5,3.2) node[midway, below=2pt, popdark]{\scriptsize ménisque};
% Initial water (1/3 volume approx)
\fill[liquid] (-1.45,0.05) .. controls (-1.45,1) and (-0.75,1.4) .. (-0.35,1.85) -- (-0.35,0.05) -- cycle;
\fill[liquid] (1.45,0.05) .. controls (1.45,1) and (0.75,1.4) .. (0.35,1.85) -- (0.35,0.05) -- cycle;
% Actually simpler fill for bottom bulb:
\fill[liquid] (-1.4,0.05) rectangle (1.4,1.2);
% Label
\node[labelstyle] at (0,-1.3) {Fiole jaugée $V_f = \SI{100,0}{mL}$\\ \scriptsize + eau distillée (1/3)};
\end{scope}

% ---- Agitation arrow ----
\draw[arrowstyle] (10.5,1.5) -- (12.0,1.5);
\node[poporange, font=\scriptsize\bfseries] at (11.25,1.9) {Agiter};

% ---- 3. Fiole jaugée (finale homogénéisée) ----
\begin{scope}[xshift=13.5cm]
% Flask outline
\draw[glass] (-1.5,0) .. controls (-1.5,1) and (-0.8,1.5) .. (-0.3,2.0);
\draw[glass] (1.5,0) .. controls (1.5,1) and (0.8,1.5) .. (0.3,2.0);
\draw[glass] (-0.3,2.0) -- (-0.3,3.5);
\draw[glass] (0.3,2.0) -- (0.3,3.5);
% Jauge mark
\draw[glass] (-0.3,3.2) -- (0.3,3.2);
% Stopper
\draw[glass, fill=popgray!20] (-0.4,3.5) rectangle (0.4,3.8);
% Liquid (full)
\fill[liquidfinal] (-1.45,0.05) .. controls (-1.45,1) and (-0.75,1.4) .. (-0.35,1.95) -- (-0.35,3.2) -- (0.35,3.2) .. controls (0.35,1.95) and (0.75,1.4) .. (1.45,1) -- (1.45,0.05) -- cycle;
% Simpler full fill:
\fill[liquidfinal] (-1.4,0.05) rectangle (1.4,1.95);
\fill[liquidfinal] (-0.3,1.95) rectangle (0.3,3.2);
% Label
\node[labelstyle] at (0,-1.3) {Solution fille\\ homogénéisée, pH $= 2{,}0$};
\end{scope}
\end{tikzpicture}
\legende{Protocole de dilution : prélèvement à la pipette jaugée, transfert dans la fiole jaugée (pré-remplie au tiers), ajustage au trait de jauge, puis homogénéisation.}
\end{popfigure}

\begin{attention}[Compléter la fiole AVANT le transfert : erreur fatale]
Une erreur très fréquente consiste à remplir la fiole jaugée d'eau jusqu'au trait de jauge \textbf{avant} d'y verser le prélèvement. En ajoutant ensuite les \SI{10,0}{mL} de solution mère, le volume dépasse $V_f$ : la concentration finale est fausse et il faut tout recommencer. Règle d'or : \textbf{le volume final s'ajuste toujours APRÈS le transfert du prélèvement}, jamais avant. On ne verse pas non plus « $V_f - V_0$ d'eau » : les volumes ne sont pas toujours additifs, seul l'ajustage au trait de jauge garantit le volume exact.
\end{attention}

\begin{attention}[Sécurité : l'acide dans l'eau, jamais l'inverse]
Lorsqu'on dilue un acide concentré (comme l'acide sulfurique ou chlorhydrique commercial), la dissolution dégage beaucoup de chaleur. Il faut \textbf{toujours verser lentement l'acide concentré dans l'eau}, en agitant. Si l'on versait l'eau dans l'acide, les premières gouttes d'eau chauffées brutalement projetteraient des gouttelettes d'acide corrosif. Port obligatoire des \cle{lunettes de protection}, des \cle{gants} et de la blouse. Moyen mnémotechnique : « \emph{l'acide plonge dans l'eau, l'eau ne plonge jamais dans l'acide} ».
\end{attention}

\courssub{Conséquence de la dilution sur le pH}

Diluer modifie la concentration en ions \ce{H3O+}, donc le pH. Retiens les deux tendances fondamentales :

\begin{propriete}[Effet de la dilution sur le pH]
\begin{itemize}
\item Diluer une solution d'\textbf{acide fort} : $[\ce{H3O+}]$ diminue, donc le \textbf{pH augmente} (il se rapproche de 7).
\item Diluer une solution de \textbf{base forte} : $[\ce{HO-}]$ diminue, donc $[\ce{H3O+}]$ augmente et le \textbf{pH diminue} (il se rapproche de 7).
\item Diluer $F = 10$ fois un acide fort de concentration $C$ (tant que $C \gg \SI{e-6}{mol.L^{-1}}$) augmente le pH d'\textbf{exactement une unité} : si $[\ce{H3O+}]$ passe de $C$ à $C/10$, alors $\text{pH} = -\log(C/10) = -\log C + 1$.
\end{itemize}
\end{propriete}

\begin{popfigure}
\begin{center}
\begin{tikzpicture}
\begin{semilogxaxis}[
  width=0.9\linewidth, height=6.5cm,
  xlabel={Facteur de dilution $F$},
  ylabel={pH},
  xmin=1, xmax=10000000, ymin=0.5, ymax=7.5,
  xtick={1,10,100,1000,10000,100000,1000000,10000000},
  xticklabels={$1$,$10$,$10^2$,$10^3$,$10^4$,$10^5$,$10^6$,$10^7$},
  ytick={1,2,3,4,5,6,7},
  minor tick num=1,
  grid=both, grid style={popline!30},
  legend style={at={(0.02,0.98)}, anchor=north west, font=\scriptsize, draw=popline, fill=popcream},
  clip=false
]
% Exact pH curve for strong acid C0 = 0.1 M diluted F times
% [H+] = (Ca + sqrt(Ca^2 + 4Kw))/2, Ca = 0.1/x
\addplot[popcurve, domain=1:1e7, samples=300]
  {-ln( (0.1/x + sqrt((0.1/x)^2 + 4e-14) ) / 2 ) / ln(10)};
\addlegendentry{Acide fort $C_0 = \SI{0,10}{mol.L^{-1}}$ (exact)}

% Asymptote pH = 7
\addplot[dashed, thick, popmuted, domain=1:1e7] {7};
\addlegendentry{pH de l'eau pure (7,0)}

% Markers for integer dilutions
\addplot[only marks, mark=*, poporange, mark size=2.5pt] coordinates {
  (1,1) (10,2) (100,3) (1000,4) (10000,5) (100000,6)
};
% Annotation
\node[popdark, font=\scriptsize, align=left, fill=popcream, rounded corners, inner sep=2pt] at (axis cs:5000,2.5) {$+1$ unité de pH\\ par dilution au $1/10$\\ (domaine approché)};
\end{semilogxaxis}
\end{tikzpicture}
\end{center}
\legende{Évolution exacte du pH d'une solution d'acide fort ($C_0 = \SI{0,10}{mol.L^{-1}}$) en fonction du facteur de dilution $F$ (échelle logarithmique). Le pH croît d'une unité par facteur 10 tant que l'acide est majoritaire, puis plafonne vers 7 : l'autoprotolyse de l'eau reprend le dessus.}
\end{popfigure}

\begin{attention}[On ne peut pas diluer un acide indéfiniment]
Croire qu'en diluant suffisamment un acide on obtiendrait un pH de 8, 9 ou 10 est une erreur conceptuelle grave. Quand la concentration de l'acide devient très faible (inférieure à environ \SI{e-6}{mol.L^{-1}}), les ions \ce{H3O+} produits par l'\cle{autoprotolyse de l'eau} ($\ce{2 H2O <=> H3O+ + HO-}$, $K_e = \SI{1,0e-14}{}$ à \SI{25}{\celsius}) ne sont plus négligeables : ils deviennent même majoritaires. Le pH d'un acide fort dilué \textbf{plafonne donc à 7,0} (il s'en approche sans jamais le dépasser), et celui d'une base forte diluée \textbf{plafonne aussi à 7,0} par le bas. Un acide dilué reste acide ; une base diluée reste basique.
\end{attention}

\begin{exemplebox}[Dilutions successives d'un acide fort]
Partons de la solution mère d'acide chlorhydrique à $C_0 = \SI{0,10}{mol.L^{-1}}$ (pH $= 1{,}0$). À chaque dilution au dixième :
\begin{center}
\begin{tabularx}{\linewidth}{|l|c|c|c|}
\hline
\rowcolor{popblueL} \textbf{Solution} & $C$ (\si{mol.L^{-1}}) & $[\ce{H3O+}]$ (\si{mol.L^{-1}}) & \textbf{pH} \\
\hline
Mère & $10^{-1}$ & $10^{-1}$ & 1,0 \\
\hline
Fille 1 ($F=10$) & $10^{-2}$ & $10^{-2}$ & 2,0 \\
\hline
Fille 2 ($F=100$) & $10^{-3}$ & $10^{-3}$ & 3,0 \\
\hline
Fille 5 ($F=10^{5}$) & $10^{-6}$ & $\approx 1,0 \times 10^{-6}$ & $\approx 6{,}0$ \\
\hline
Fille 7 ($F=10^{7}$) & $10^{-8}$ & $\approx 1,0 \times 10^{-7}$ & $\approx 7{,}0$ \\
\hline
\end{tabularx}
\end{center}
Au-delà de $F = 10^{6}$, le pH reste pratiquement égal à 7,0 : l'eau distillée elle-même impose son pH.
\end{exemplebox}

\begin{exercice}[À toi de jouer]
Un laborantin doit préparer \SI{250,0}{mL} d'une solution d'hydroxyde de sodium (soude) de concentration $C_f = \SI{0,020}{mol.L^{-1}}$ à partir d'une solution mère à $C_0 = \SI{0,50}{mol.L^{-1}}$.
\begin{enumerate}
\item Calculer le volume $V_0$ de solution mère à prélever et le facteur de dilution.
\item Calculer le pH de la solution fille à \SI{25}{\celsius} (la soude est une base forte : $[\ce{HO-}] = C_f$).
\item Citer la verrerie utilisée pour le prélèvement et pour l'ajustage.
\end{enumerate}
\end{exercice}

\begin{corrige}[Exercice : dilution de la soude]
\textbf{1.} $V_0 = \dfrac{C_f V_f}{C_0} = \dfrac{0{,}020 \times 250{,}0}{0{,}50} = \SI{10,0}{mL}$. Facteur de dilution : $F = \dfrac{C_0}{C_f} = \dfrac{0{,}50}{0{,}020} = 25$.

\textbf{2.} $[\ce{HO-}] = \SI{0,020}{mol.L^{-1}}$, donc $[\ce{H3O+}] = \dfrac{K_e}{[\ce{HO-}]} = \dfrac{10^{-14}}{0{,}020} = \SI{5,0e-13}{mol.L^{-1}}$, d'où $\text{pH} = -\log(5{,}0 \times 10^{-13}) \approx 12{,}3$. La solution fille est fortement basique.

\textbf{3.} Prélèvement à la \textbf{pipette jaugée} de \SI{10,0}{mL} ; ajustage en \textbf{fiole jaugée} de \SI{250,0}{mL}, puis homogénéisation.
\end{corrige}

\begin{savaistu}[Les solutions étalons des laboratoires camerounais]
Dans les laboratoires d'analyses médicales (hôpitaux de Yaoundé, Douala, Garoua) comme dans les usines (SONARA, CIMENCAM, brasseries), les pH-mètres doivent être \textbf{étalonnés} régulièrement à l'aide de \cle{solutions étalons} de pH connu avec précision (par exemple pH $= 4{,}00$, $7{,}00$ et $10{,}00$). Ces solutions sont préparées exactement par la méthode que tu viens d'apprendre : pesées précises, dilutions successives en fiole jaugée, eau distillée exempte de \ce{CO2} dissous. Sans cette rigueur de préparation, un dosage sanguin ou un contrôle qualité de l'huile de palme pourrait donner un résultat faux... avec des conséquences graves sur un diagnostic médical ou sur la conformité d'un produit !
\end{savaistu}

\begin{aretenir}[L'essentiel de la section]
\begin{itemize}
\item Diluer = ajouter du solvant \textbf{sans changer la quantité de matière de soluté} : $C_0 V_0 = C_f V_f$.
\item Facteur de dilution : $F = \dfrac{C_0}{C_f} = \dfrac{V_f}{V_0}$.
\item Volume à prélever : $V_0 = \dfrac{C_f V_f}{C_0}$, prélevé à la \textbf{pipette jaugée}, ajusté en \textbf{fiole jaugée} au trait de jauge, puis homogénéisé.
\item Diluer un acide fort fait \textbf{augmenter} le pH ($+1$ par dilution au $1/10$ tant que $C \gg 10^{-6}$ M) ; diluer une base forte fait \textbf{diminuer} le pH ; dans les deux cas, le pH \textbf{plafonne à 7} à cause de l'autoprotolyse de l'eau.
\item Sécurité : toujours verser \textbf{l'acide dans l'eau}, avec lunettes et gants.
\end{itemize}
\end{aretenir}

\newpage

\coursec{Travaux pratiques}

\courssub{Objectifs du TP}

À l'issue de cette séance pratique, tu seras capable de :

\begin{itemize}
\item mesurer le pH de produits usuels (jus de citron, eau minérale, eau de forage, savon dissous, vinaigre, ammoniaque diluée) à l'aide du \cle{papier pH} ;
\item confirmer le caractère acide, neutre ou basique d'une solution à l'aide des \cle{indicateurs colorés} : bleu de bromothymol (BBT), hélianthine et phénolphtaléine ;
\item calculer le volume à prélever pour préparer une solution diluée de facteur $F = 10$ ;
\item préparer \SI{100}{mL} d'une solution d'acide chlorhydrique à \SI{0,01}{mol.L^{-1}} à partir d'une solution mère à \SI{0,1}{mol.L^{-1}}, en utilisant la pipette jaugée et la fiole jaugée ;
\item vérifier expérimentalement que la dilution d'un acide fort par 10 fait augmenter le pH d'environ une unité ;
\item calculer les concentrations $[\ce{H3O+}]$ et $[\ce{HO-}]$ à partir du pH mesuré et classer les produits par acidité croissante.
\end{itemize}

\courssub{Matériel et produits}

\begin{center}
\rowcolors{2}{popblueL}{white}
\begin{tabularx}{\linewidth}{|l|X|}
\hline
\rowcolor{popblue!40} \textbf{Matériel} & \textbf{Produits} \\
\hline
6 tubes à essais + portoir & Jus de citron fraîchement pressé \\
Pipettes plastiques ou compte-gouttes & Eau minérale locale (Supermont, Tangui\ldots) \\
Papier pH + échelle de teintes & Eau de forage ou de puits \\
Spatule propre et sèche & Solution de savon (savon local dissous) \\
Pipette jaugée de \SI{10}{mL} + poire & Vinaigre du commerce (ou vinaigre de vin de palme) \\
Fiole jaugée de \SI{100}{mL} & Ammoniaque diluée (nettoyant ménager dilué) \\
Béchers de \SI{100}{mL} et \SI{250}{mL} & Acide chlorhydrique à \SI{0,1}{mol.L^{-1}} \\
pH-mètre étalonné (si disponible) & BBT, hélianthine, phénolphtaléine \\
Pissette d'eau distillée & Eau distillée \\
\hline
\end{tabularx}
\end{center}

\begin{savaistu}[Alternatives avec du matériel local]
Si le papier pH manque, on peut utiliser du \emph{papier de curcuma} (vire au brun-rouge en milieu basique) ou un jus de \emph{bissap} (hibiscus) filtré : ce jus est rouge en milieu acide et vire au vert-bleu en milieu basique, c'est un excellent indicateur coloré naturel. À défaut de fiole jaugée, on peut utiliser une bouteille propre graduée au marqueur à l'aide d'une éprouvette, avec une précision moindre. Le vinaigre de vin de palme sure remplace avantageusement le vinaigre du commerce.
\end{savaistu}

\courssub{Sécurité}

\begin{attention}[Consignes de sécurité]
\begin{itemize}
\item \textbf{Blouse, lunettes et gants} obligatoires pendant toute la séance.
\item L'\textbf{acide chlorhydrique} à \SI{0,1}{mol.L^{-1}} est \emph{corrosif} (pictogramme : deux gouttes attaquant une main et une barre métallique). En cas de contact avec la peau, rincer abondamment à l'eau froide pendant plusieurs minutes et prévenir l'enseignant.
\item L'\textbf{ammoniaque} dégage des vapeurs irritantes (pictogramme : point d'exclamation). Ne jamais sentir directement le flacon ; travailler près d'une fenêtre ouverte.
\item \textbf{Jamais de pipetage à la bouche} : utiliser systématiquement la poire à pipeter (propipette).
\item Ne jamais goûter les produits, même alimentaires, une fois apportés au laboratoire.
\item Verser les déchets acides et basiques dans le bidon de récupération prévu, pas dans l'évier.
\end{itemize}
\end{attention}

\courssub{Schéma du dispositif}

\begin{popfigure}
\begin{center}
\begin{tikzpicture}[scale=0.95, every node/.style={font=\small}]
\% Bécher solution mère
\draw[thick, popblue] (0,0) -- (0,1.6) (1.6,0) -- (1.6,1.6);
\draw[thick, popblue] (0,0) -- (1.6,0);
\fill[popblue!25] (0.05,0.05) rectangle (1.55,1.0);
\node[popdark] at (0.8,-0.45) {Solution mère};
\node[popdark] at (0.8,-0.9) {\SI{0,1}{mol.L^{-1}}};
\% Pipette jaugée
\draw[thick, poporange] (0.8,1.0) -- (0.8,3.6);
\draw[thick, poporange] (0.95,1.0) -- (0.95,3.6);
\draw[poporange] (0.68,3.0) -- (1.07,3.0);
\node[poporange, right] at (1.1,3.0) {trait de jauge};
\node[popdark, left] at (0.6,2.4) {Pipette jaugée \SI{10}{mL}};
\% Flèche transfert
\draw[-{Stealth}, very thick, poppurple] (1.9,2.6) -- (3.6,2.6) node[midway, above]{$V_0 = \SI{10,0}{mL}$};
\% Fiole jaugée
\begin{scope}[xshift=4.2cm]
\draw[thick, popgreen] (0.55,2.6) -- (0.55,3.8) (1.05,2.6) -- (1.05,3.8);
\draw[thick, popgreen] (0.55,2.6) .. controls (0.3,2.0) and (0.0,1.6) .. (0.0,0.8) .. controls (0.0,0.1) .. (0.8,0.0) .. controls (1.6,0.1) .. (1.6,0.8) .. controls (1.6,1.6) and (1.3,2.0) .. (1.05,2.6);
\draw[popgreen] (0.5,3.2) -- (1.1,3.2);
\node[popgreen, right] at (1.15,3.2) {trait de jauge};
\fill[popgreen!25] (0.06,0.75) .. controls (0.06,0.2) .. (0.8,0.06) .. controls (1.54,0.2) .. (1.54,0.75) .. controls (1.54,1.5) and (1.25,1.95) .. (1.02,2.5) -- (0.58,2.5) .. controls (0.35,1.95) and (0.06,1.5) .. (0.06,0.75);
\node[popdark] at (0.8,-0.45) {Fiole jaugée \SI{100}{mL}};
\node[popdark] at (0.8,-0.9) {Solution fille \SI{0,01}{mol.L^{-1}}};
\end{scope}
\% Mesure papier pH
\begin{scope}[xshift=8.2cm]
\draw[thick, popgold] (0,0.4) rectangle (2.6,1.0);
\foreach \i/\c in {0/poppink, 0.65/poporange, 1.3/popgold, 1.95/popgreen}{\fill[\c] (\i+0.02,0.42) rectangle (\i+0.63,0.98);}
\node[popdark] at (1.3,0.0) {Papier pH + échelle de teintes};
\draw[thick, popblue] (0.4,1.9) -- (0.4,3.0) (1.4,1.9) -- (1.4,3.0);
\draw[thick, popblue] (0.4,1.9) -- (1.4,1.9);
\fill[popblue!25] (0.44,1.94) rectangle (1.36,2.5);
\node[popdark, above] at (0.9,3.0) {solution à tester};
\draw[-{Stealth}, thick, popmuted] (1.0,1.85) -- (1.0,1.15);
\end{scope}
\end{tikzpicture}
\end{center}
\legende{Dispositif du TP : prélèvement à la pipette jaugée, dilution en fiole jaugée, puis mesure du pH au papier pH.}
\end{popfigure}

\courssub{Protocole}

\textbf{Partie A : mesure du pH de produits usuels (séance 1)}\par

\begin{enumerate}
\item Étiqueter six tubes à essais : citron, eau minérale, eau de forage, savon, vinaigre, ammoniaque.
\item Verser environ \SI{2}{mL} de chaque produit dans le tube correspondant.
\item À l'aide d'une spatule propre et sèche, déposer une goutte de chaque solution sur une bandelette de papier pH. \textbf{Ne jamais plonger le papier pH dans le flacon} : cela contaminerait tout le produit.
\item Comparer immédiatement la teinte obtenue à l'échelle de référence et noter le pH au demi-unité près dans le tableau de résultats.
\item Tester chaque solution avec l'indicateur coloré le plus adapté : \textbf{hélianthine} pour les solutions supposées acides (citron, vinaigre), \textbf{bleu de bromothymol (BBT)} pour les solutions supposées neutres (eaux), \textbf{phénolphtaléine} pour les solutions supposées basiques (savon, ammoniaque). Ajouter 2 gouttes d'indicateur dans chaque tube, homogénéiser, noter la teinte observée et conclure sur le caractère acide, neutre ou basique.
\end{enumerate}

\textbf{Partie B : dilution d'une solution d'acide chlorhydrique (séance 2)}\par

\begin{enumerate}
\item \textbf{Calcul préalable} : la solution mère a pour concentration $C_0 = \SI{0,10}{mol.L^{-1}}$. On veut $V_1 = \SI{100}{mL}$ de solution fille à $C_1 = \SI{0,010}{mol.L^{-1}}$. Calculer le volume $V_0$ à prélever grâce à la relation de dilution $C_0 V_0 = C_1 V_1$.
\item Mesurer le pH de la solution mère au papier pH (ou au pH-mètre étalonné). Noter la valeur.
\item Verser environ \SI{30}{mL} de solution mère dans un bécher propre et sec. Rincer la pipette jaugée de \SI{10}{mL} avec un peu de cette solution.
\item Prélever $V_0 = \SI{10,0}{mL}$ de solution mère à la pipette jaugée munie de la poire, en ajustant le bas du ménisque sur le trait de jauge.
\item Introduire le prélèvement dans la fiole jaugée de \SI{100}{mL} contenant déjà environ \SI{30}{mL} d'eau distillée. Homogénéiser.
\item Compléter avec de l'eau distillée jusqu'au trait de jauge, en terminant à la pissette, le bas du ménisque tangent au trait. Boucher et retourner dix fois pour homogénéiser.
\item Mesurer le pH de la solution fille. Comparer avec celui de la solution mère.
\end{enumerate}

\courssub{Observations et résultats attendus}

\begin{center}
\rowcolors{2}{popgreenL}{white}
\begin{tabularx}{\linewidth}{|X|c|c|c|}
\hline
\rowcolor{popgreen!40} \textbf{Produit} & \textbf{pH mesuré} & \textbf{Teinte avec indicateur} & \textbf{Caractère} \\
\hline
Jus de citron & 2,5 & Hélianthine : rouge & Acide \\
Vinaigre & 3,0 & Hélianthine : rouge-orangé & Acide \\
Eau minérale locale & 7,0 & BBT : vert & Neutre \\
Eau de forage & 6,5 & BBT : vert-jaune & Légèrement acide \\
Solution de savon & 10,0 & Phénolphtaléine : rose & Basique \\
Ammoniaque diluée & 11,0 & Phénolphtaléine : rose vif & Basique \\
\hline
Solution mère \ce{HCl} \SI{0,1}{mol.L^{-1}} & 1,0 & --- & Acide \\
Solution fille \ce{HCl} \SI{0,01}{mol.L^{-1}} & 2,0 & --- & Acide \\
\hline
\end{tabularx}
\end{center}

On constate que la dilution au dixième de la solution d'acide chlorhydrique a fait passer le pH de 1,0 à 2,0 : le pH a \textbf{augmenté d'une unité}, car la concentration en ions \ce{H3O+} a été divisée par 10.

\courssub{Exploitation}

\begin{exoresolu}[Exploitation guidée des mesures]
\textbf{Questions.}
\begin{enumerate}
\item Vérifier le calcul du volume à prélever pour la dilution.
\item Calculer $[\ce{H3O+}]$ et $[\ce{HO-}]$ dans le jus de citron ($\mathrm{pH} = 2{,}5$) à \SI{25}{\celsius}.
\item Classer les six produits usuels par acidité croissante.
\item Expliquer pourquoi le pH de la solution fille vaut environ une unité de plus que celui de la solution mère.
\end{enumerate}
\tcblower
\textbf{1. Volume à prélever.} La conservation de la quantité de matière lors de la dilution s'écrit $C_0 V_0 = C_1 V_1$, d'où :
\[ V_0 = \frac{C_1 V_1}{C_0} = \frac{0{,}010 \times 100}{0{,}10} = \SI{10,0}{mL}. \]
Le facteur de dilution est $F = \dfrac{C_0}{C_1} = 10$.

\textbf{2. Concentrations dans le jus de citron.} À \SI{25}{\celsius}, $K_e = 10^{-14}$ :
\[ [\ce{H3O+}] = 10^{-\mathrm{pH}} = 10^{-2{,}5} \approx \SI{3,2e-3}{mol.L^{-1}}. \]
\[ [\ce{HO-}] = \frac{K_e}{[\ce{H3O+}]} = \frac{10^{-14}}{3{,}2 \times 10^{-3}} \approx \SI{3,2e-12}{mol.L^{-1}}. \]
Les ions \ce{H3O+} sont majoritaires ; les ions \ce{HO-} sont ultra-minoritaires, mais jamais absents.

\textbf{3. Classement par acidité croissante} (pH croissant) :
\[ \text{citron} < \text{vinaigre} < \text{eau de forage} < \text{eau minérale} < \text{savon} < \text{ammoniaque}. \]

\textbf{4. Effet de la dilution.} L'acide chlorhydrique est un acide fort, totalement dissocié : $[\ce{H3O+}] = C$. En diluant 10 fois, $[\ce{H3O+}]$ passe de \SI{0,10}{mol.L^{-1}} à \SI{0,010}{mol.L^{-1}}, donc le pH passe de $-\log 0{,}10 = 1{,}0$ à $-\log 0{,}010 = 2{,}0$ : il augmente exactement d'une unité, car $\mathrm{pH} = -\log[\ce{H3O+}]$ et $\log([\ce{H3O+}]/10) = \log[\ce{H3O+}] - 1$, d'où $\mathrm{pH}_{\text{fille}} = \mathrm{pH}_{\text{mère}} + 1$.
\end{exoresolu}

\begin{attention}[Erreurs fréquentes à éviter]
\begin{itemize}
\item Confondre \og diluer 10 fois \fg{} avec \og ajouter 10 fois plus d'eau \fg{} : on complète \emph{jusqu'à} \SI{100}{mL}, on n'ajoute pas \SI{100}{mL} d'eau.
\item Oublier que le pH augmente quand on dilue un acide (la solution devient \emph{moins} acide, jamais basique : son pH tend vers 7 sans le dépasser).
\item Lire le haut du ménisque au lieu du bas lors de l'ajustage en fiole jaugée.
\item Croire qu'une solution de pH 7 ne contient ni \ce{H3O+} ni \ce{HO-} : les deux ions sont toujours présents, en concentrations égales à $10^{-7}$ \si{mol.L^{-1}} à \SI{25}{\celsius}.
\end{itemize}
\end{attention}

\courssub{Conclusion}

Ce TP t'a permis de relier la théorie à la vie quotidienne : les produits que tu utilises chaque jour ont un pH mesurable, qui traduit leur caractère acide (citron, vinaigre), neutre (eau minérale) ou basique (savon, ammoniaque). Tu as vérifié expérimentalement la relation $\mathrm{pH} = -\log[\ce{H3O+}]$ : diviser la concentration d'un acide fort par 10 augmente son pH d'une unité. Tu maîtrises désormais les gestes de la dilution de précision — prélèvement à la pipette jaugée, ajustage en fiole jaugée — et l'utilisation du papier pH et des indicateurs colorés, compétences indispensables pour interpréter, par exemple, le résultat d'une analyse médicale ou la qualité d'une eau de boisson.

\newpage

\coursec{Méthodes et exercices résolus}

\coursec{Leçon 06 : Acides et bases en solution aqueuse : eau, autoprotolyse et pH}

\courssub{L'eau : structure et pouvoir solvant}

\begin{definition}[Molécule d'eau]
La molécule d'eau \ce{H2O} possède une géométrie \textbf{coudée} (angle \ce{H-O-H} \approx \SI{104,5}{\degree}) due aux deux doublets non liants sur l'oxygène. L'oxygène, très électronégatif ($\chi_{\ce{O}} = 3,44$), attire les nuages électroniques des liaisons \ce{O-H} : la molécule est \textbf{polaire}, porteuse d'un moment dipolaire permanent $\vec{\mu} \approx \SI{1,85}{D}$.
\end{definition}

\begin{propriete}[Liaisons hydrogène et conséquences macroscopiques]
En phase condensée, les molécules d'eau s'associent par des \cle{liaisons hydrogène} (interactions dipôle-dipôle fortes entre le \ce{H} \delta+ d'une molécule et le doublet non liant de l'\ce{O} \delta- d'une voisine).
\begin{itemize}
\item \textbf{Point d'ébullition anormalement élevé} (\SI{100}{\celsius} au lieu de \SI{-80}{\celsius} prévu sans liaisons H).
\item \textbf{Anomalie de la densité} : la glace est moins dense que l'eau liquide (maximum de densité à \SI{4}{\celsius}), ce qui permet la vie aquatique en hiver.
\item \textbf{Tension de surface élevée} : phénomène de capillarité, essentiel pour la montée de sève chez les plantes (du Mont Cameroun aux cultures du Littoral).
\end{itemize}
\end{propriete}

\begin{experience}[Modélisation du pouvoir solvant de l'eau]
\textbf{Dispositif :} Bécher d'eau, cristaux de sulfate de cuivre(II) \ce{CuSO4} (bleu), éthanol, huile.
\textbf{Protocole :} Verser \ce{CuSO4} dans l'eau $\to$ dissolution rapide, solution bleue homogène. Verser \ce{CuSO4} dans l'éthanol $\to$ dissolution lente, partielle. Verser \ce{CuSO4} dans l'huile $\to$ pas de dissolution.
\textbf{Interprétation :} L'eau solvatise les ions \ce{Cu^2+} et \ce{SO4^2-} par ses dipôles (sphères d'hydratation). L'éthanol, polaire mais moins que l'eau, solvatise mal les ions. L'huile, apolaire, ne solvatise pas les ions.
\textbf{Conclusion :} L'eau est un \textbf{solvant polaire protique} excellent pour les espèces ioniques et les molécules polaires (sucre, éthanol, acides, bases), médiocre pour les espèces apolaires (huiles, graisses).
\end{experience}

\begin{savaistu}[L'eau au Cameroun : entre resource et chimie]
Les sources du Mont Cameroun (eaux de Buea, Limbe) sont riches en ions dissous (\ce{HCO3-}, \ce{Na+}, \ce{Mg^2+}, \ce{Ca^2+}) grâce au pouvoir solvant de l'eau sur les roches volcaniques. L'eau de mer (golfe de Guinée) contient $\approx$ \SI{35}{g.L^{-1}} de sels dissous, majoritairement \ce{NaCl}. La SONARA (Société Nationale de Raffinage) utilise l'eau comme fluide caloporteur et solvant dans ses procédés de raffinage. La fabrication traditionnelle du savon (saponification de l'huile de palme par la lessive de cendres, riche en \ce{K2CO3} et \ce{Na2CO3}) repose sur la dissolution des bases dans l'eau.
\end{savaistu}

\begin{aretenir}[Bilan : l'eau, solvant de la chimie aqueuse]
\begin{itemize}
\item Molécule polaire, coudée, liaisons hydrogène.
\item Solvant privilégié des espèces ioniques (solvatation) et polaires.
\item Ampholyte : peut agir comme acide ou comme base (autoprotolyse).
\end{itemize}
\end{aretenir}

\courssub{Acides et bases selon Brønsted}

\textbf{Méthode 1 : Reconnaître un acide, une base et un ampholyte}\par
\begin{methode}[Identifier un acide, une base, un ampholyte selon Brønsted]
\begin{enumerate}
\item \textbf{Repérer les protons échangeables} : cherche un atome d'hydrogène lié à un atome très électronégatif (O, N, halogène) : c'est un \cle{proton labile}.
\item \textbf{Appliquer la définition de Brønsted} : un \cle{acide} est une espèce capable de \textbf{céder} un proton \ce{H+} ; une \cle{base} est une espèce capable de \textbf{capter} un proton \ce{H+}.
\item \textbf{Écrire le couple acide/base} : en retirant \ce{H+} à l'acide, on obtient sa base conjuguée : $\text{acide} = \text{base} + \ce{H+}$. Les deux espèces forment un \cle{couple} noté $\text{acide}/\text{base}$.
\item \textbf{Repérer les ampholytes} : une espèce qui peut à la fois céder et capter un proton (ex. \ce{H2O}, \ce{HCO3-}, \ce{H2PO4-}) est un \cle{ampholyte} (ou espèce amphotère).
\item \textbf{Écrire une réaction acido-basique} : c'est un transfert de proton de l'acide du couple 1 vers la base du couple 2 :
\[ \text{acide}_1 + \text{base}_2 \ce{->} \text{base}_1 + \text{acide}_2 \]
\end{enumerate}
\textbf{Réflexe Bac} : dans une réaction acido-basique, il y a toujours \textbf{deux couples} en jeu. Le proton ne voyage jamais seul en solution aqueuse : il est immédiatement solvaté par l'eau sous forme \ce{H3O+} (ion oxonium).
\end{methode}

\begin{exoresolu}[L'acide du jus de bissap et l'eau de Javel]
Le jus de bissap (infusion de fleurs d'hibiscus \emph{Hibiscus sabdariffa}, très consommé au Cameroun) doit sa couleur rouge et son goût acidulé à des acides organiques, dont l'acide citrique noté simplement \ce{H3A}. L'eau de Javel commerciale contient l'ion hypochlorite \ce{ClO-}.
\begin{enumerate}
\item Écrire le couple acide/base correspondant à \ce{H3A} (première acidité) et à \ce{ClO-}.
\item L'ion hydrogénocarbonate \ce{HCO3-} (présent dans les eaux minérales du Mont Cameroun) est-il un acide ou une base ? Justifier.
\item Écrire l'équation de la réaction entre \ce{H3A} et \ce{ClO-}.
\end{enumerate}
\tcblower
\textbf{1. Couples acide/base.}

D'après la définition de Brønsted, la base conjuguée s'obtient en retirant un proton \ce{H+} à l'acide :
\begin{itemize}
\item Acide citrique (premier proton) : $\ce{H3A = H2A- + H+}$, d'où le couple \ce{H3A}/\ce{H2A-}.
\item L'ion hypochlorite capte un proton : $\ce{ClO-} + \ce{H+} = \ce{HClO}$, d'où le couple \ce{HClO}/\ce{ClO-} ; \ce{ClO-} est donc une \textbf{base} (base conjuguée de l'acide hypochloreux \ce{HClO}).
\end{itemize}

\textbf{2. Nature de \ce{HCO3-}.}

L'ion \ce{HCO3-} peut \textbf{céder} un proton : $\ce{HCO3- = CO3^2- + H+}$ (couple \ce{HCO3-}/\ce{CO3^2-}), et il peut aussi \textbf{capter} un proton : $\ce{HCO3-} + \ce{H+} = \ce{H2CO3}$ (couple \ce{H2CO3}/\ce{HCO3-}).

\ce{HCO3-} est donc à la fois acide et base : c'est un \textbf{ampholyte}.

\textbf{3. Équation de la réaction.}

On écrit les deux demi-équations de transfert de proton :
\[ \ce{H3A = H2A- + H+} \qquad \ce{ClO- + H+ = HClO} \]
En additionnant membre à membre (le proton \ce{H+} s'élimine) :
\[ \ce{H3A + ClO- -> H2A- + HClO} \]
\textbf{Conclusion} : le jus de bissap (acide) réagit avec l'eau de Javel (base) par transfert de proton ; c'est pourquoi il ne faut \textbf{jamais} mélanger des produits usuels acides (détartrant, vinaigre, jus de citron/bissap) et basiques (eau de Javel, soude) : cela libère du chlore gazeux \ce{Cl2} toxique.
\end{exoresolu}

\courssub{Autoprotolyse de l'eau et produit ionique Ke}

\textbf{Méthode 2 : Exploiter l'autoprotolyse de l'eau}\par
\begin{methode}[Exploiter l'autoprotolyse de l'eau]
\begin{enumerate}
\item \textbf{Écrire l'équation-bilan} : l'eau, ampholyte, réagit sur elle-même par transfert de proton d'une molécule sur l'autre :
\[ \ce{2 H2O <=> H3O+ + HO-} \]
Cette réaction, très limitée, est l'\cle{autoprotolyse de l'eau}.
\item \textbf{Écrire le produit ionique} : $K_e = [\ce{H3O+}][\ce{HO-}]$, constante d'équilibre qui ne dépend \textbf{que de la température} (activités approximées par les concentrations molaires pour les solutions diluées).
\item \textbf{Retenir la valeur de référence} : à \SI{25}{\celsius}, $K_e = 1{,}0 \times 10^{-14}$, soit $pK_e = -\log K_e = 14{,}00$.
\item \textbf{Exploiter la neutralité de l'eau pure} : dans l'eau pure, l'électroneutralité impose $[\ce{H3O+}] = [\ce{HO-}] = \sqrt{K_e} = 1{,}0 \times 10^{-7}\ \si{mol.L^{-1}}$ à \SI{25}{\celsius}.
\item \textbf{Généraliser} : dans \textbf{toute} solution aqueuse, $[\ce{H3O+}][\ce{HO-}] = K_e$ : connaître l'une des deux concentrations donne immédiatement l'autre.
\end{enumerate}
\textbf{Réflexe Bac} : $K_e$ augmente avec la température (autoprotolyse endothermique) : à \SI{37}{\celsius} (température du corps humain), $pK_e \approx 13{,}6$ et la neutralité correspond à pH $\approx 6{,}8$ (et non 7) !
\end{methode}

\begin{exoresolu}[L'eau du robinet à différentes températures]
On mesure à \SI{25}{\celsius} la concentration en ions oxonium d'une eau minérale : $[\ce{H3O+}] = 4{,}0 \times 10^{-8}\ \si{mol.L^{-1}}$.
\begin{enumerate}
\item Écrire l'équation de l'autoprotolyse de l'eau.
\item Calculer la concentration en ions hydroxyde de cette eau à \SI{25}{\celsius}.
\item Cette eau est-elle acide, neutre ou basique ?
\item À \SI{60}{\celsius}, $pK_e = 13{,}0$. Quelle est la concentration en ions \ce{H3O+} de l'eau pure à cette température ? L'eau pure est-elle devenue acide ?
\end{enumerate}
\tcblower
\textbf{1. Équation de l'autoprotolyse.}
\[ \ce{2 H2O <=> H3O+ + HO-} \]

\textbf{2. Concentration en \ce{HO-}.}

À \SI{25}{\celsius}, $K_e = 1{,}0 \times 10^{-14}$. D'après $K_e = [\ce{H3O+}][\ce{HO-}]$ :
\[ [\ce{HO-}] = \frac{K_e}{[\ce{H3O+}]} = \frac{1{,}0 \times 10^{-14}}{4{,}0 \times 10^{-8}} = 2{,}5 \times 10^{-7}\ \si{mol.L^{-1}} \]

\textbf{3. Caractère de l'eau.}

On compare : $[\ce{HO-}] = 2{,}5 \times 10^{-7} > [\ce{H3O+}] = 4{,}0 \times 10^{-8}\ \si{mol.L^{-1}}$.

La solution contient plus d'ions hydroxyde que d'ions oxonium : elle est \textbf{basique} (pH $= -\log(4{,}0 \times 10^{-8}) = 7{,}40 > 7$).

\textbf{4. Eau pure à \SI{60}{\celsius}.}

Dans l'eau pure, l'électroneutralité impose $[\ce{H3O+}] = [\ce{HO-}]$, donc :
\[ [\ce{H3O+}]^2 = K_e \implies [\ce{H3O+}] = \sqrt{K_e} = \sqrt{10^{-13}} = 10^{-6{,}5} \approx 3{,}16 \times 10^{-7}\ \si{mol.L^{-1}} \]
Le pH de l'eau pure vaut alors $-\log(3{,}16 \times 10^{-7}) = 6{,}50 < 7$.

\textbf{Conclusion} : l'eau pure à \SI{60}{\celsius} a un pH de 6,50 mais reste \textbf{neutre}, car $[\ce{H3O+}] = [\ce{HO-}]$. Le caractère neutre se définit par l'égalité des concentrations, pas par pH = 7 (valable seulement à \SI{25}{\celsius}).
\end{exoresolu}

\courssub{Le pH : définition et mesure}

\begin{definition}[pH d'une solution aqueuse]
Le \cle{pH} (potentiel hydrogène) est défini par :
\[ \text{pH} = -\log_{10}\left( \frac{[\ce{H3O+}]}{\SI{1}{mol.L^{-1}}} \right) \]
C'est une grandeur \textbf{sans unité}. L'argument du logarithme est sans dimension (concentration divisée par l'unité standard).
Réciproquement : $[\ce{H3O+}] = 10^{-\text{pH}}\ \si{mol.L^{-1}}$.
\end{definition}

\begin{propriete}[Échelle de pH à \SI{25}{\celsius}]
\begin{center}
\begin{tabularx}{\linewidth}{|c|X|X|X|}
\hline
\rowcolor{popblueL}
\textbf{Domaine} & \textbf{pH} & \textbf{Comparaison $[\ce{H3O+}]$ vs $[\ce{HO-}]$} & \textbf{Exemples camerounais} \\
\hline
Acide & $< 7$ & $[\ce{H3O+}] > [\ce{HO-}]$ & Jus de bissap (pH 2,5--3,5), vin de palme fermenté (pH 3,5--4,5), Coca-Cola (pH 2,5), acide de batterie (pH $\approx$ 1) \\
\hline
Neutre & $= 7$ & $[\ce{H3O+}] = [\ce{HO-}] = 10^{-7}$ & Eau pure distillée à \SI{25}{\celsius}, sang (pH 7,4 à \SI{37}{\celsius} $\to$ neutre car $pK_e=13,6$) \\
\hline
Basique & $> 7$ & $[\ce{H3O+}] < [\ce{HO-}]$ & Eau de Javel (pH 11--13), lessive de cendres (pH 10--11), savon (pH 9--10), eau de mer (pH 8,1) \\
\hline
\end{tabularx}
\end{center}
\end{propriete}

\textbf{Méthode 3 : Choisir et utiliser une méthode de mesure du pH}\par
\begin{methode}[Choisir et utiliser une méthode de mesure du pH]
\begin{enumerate}
\item \textbf{Mesure précise au pH-mètre} (laboratoire, industrie : SONARA, brasseries, laiteries) :
\begin{itemize}
\item \textbf{Étalonner} l'appareil avec deux solutions tampons étalons (généralement pH 4,01 et 7,00 ou 7,00 et 10,01) à la température de la solution à mesurer ;
\item Rincer l'électrode à l'eau distillée, l'essuyer délicatement avec du papier absorbant (ne pas frotter la bulle de verre) ;
\item Plonger l'électrode dans la solution, agiter légèrement, attendre la stabilisation (drift $< \SI{0,01}{pH/min}$) ;
\item Lire le pH (précision : \SI{0,05}{pH} à \SI{0,1}{pH}).
\end{itemize}
\item \textbf{Mesure rapide au papier pH} (terrain, contrôle rapide) : déposer une goutte de solution sur le papier (ne \textbf{jamais} plonger le papier dans le flacon pour ne pas contaminer la solution), comparer la couleur à l'échelle de référence (précision : \approx 1 unité de pH).
\item \textbf{Estimation par indicateurs colorés} (dosage, repérage d'équivalence) : un \cle{indicateur coloré} (ex. phénolphtaléine, vert de bromocrésol, orange de méthyle) change de teinte dans une \cle{zone de virage} (intervalle de pH $\approx pK_a \pm 1$) ; quelques gouttes dans la solution permettent d'encadrer le pH.
\item \textbf{Choisir la méthode adaptée} : contrôle qualité industriel $\to$ pH-mètre ; test rapide sur le terrain (source, rivière) $\to$ papier pH ; repérage d'équivalence en dosage $\to$ indicateur coloré.
\end{enumerate}
\textbf{Réflexe Bac} : une mesure de pH donne $[\ce{H3O+}]$ avec au mieux 2 chiffres significatifs (ex. pH = 4,30 $\to$ $[\ce{H3O+}] = 5,0 \times 10^{-5}\ \si{mol.L^{-1}}$).
\end{methode}

\begin{exoresolu}[Contrôle du pH du vin de palme]
Le vin de palme frais (prélevé le matin sur le raphia ou le palmier à huile) a un pH voisin de 7 ; en s'acidifiant (fermentation lactique puis acétique), son pH chute. Un producteur de la région du Littoral mesure le pH de trois échantillons avec du papier pH, puis vérifie l'échantillon 2 au pH-mètre : il lit pH $= 4{,}30$.
\begin{enumerate}
\item Citer les deux étapes indispensables avant toute mesure au pH-mètre.
\item Calculer $[\ce{H3O+}]$ et $[\ce{HO-}]$ dans l'échantillon 2 à \SI{25}{\celsius}.
\item Le producteur affirme : « l'acidité a été multipliée par 1000 par rapport au vin frais ». Vérifier.
\end{enumerate}
Donnée : $10^{-4{,}30} = 5{,}0 \times 10^{-5}$.
\tcblower
\textbf{1. Avant la mesure.}

Il faut : (a) \textbf{étalonner} le pH-mètre avec des solutions tampons étalons (pH 4 et 7 par exemple) ; (b) \textbf{rincer et sécher} l'électrode à l'eau distillée avant de la plonger dans l'échantillon.

\textbf{2. Concentrations ioniques.}
\[ [\ce{H3O+}] = 10^{-\text{pH}} = 10^{-4{,}30} = 5{,}0 \times 10^{-5}\ \si{mol.L^{-1}} \]
\[ [\ce{HO-}] = \frac{K_e}{[\ce{H3O+}]} = \frac{1{,}0 \times 10^{-14}}{5{,}0 \times 10^{-5}} = 2{,}0 \times 10^{-10}\ \si{mol.L^{-1}} \]

\textbf{3. Vérification de l'affirmation.}

Vin frais : pH $\approx 7{,}00$, donc $[\ce{H3O+}]_{\text{frais}} = 1{,}0 \times 10^{-7}\ \si{mol.L^{-1}}$.
\[ \frac{[\ce{H3O+}]_{\text{échantillon 2}}}{[\ce{H3O+}]_{\text{frais}}} = \frac{5{,}0 \times 10^{-5}}{1{,}0 \times 10^{-7}} = 5{,}0 \times 10^{2} = 500 \]
\textbf{Conclusion} : l'acidité (concentration en \ce{H3O+}) a été multipliée par 500, et non par 1000 : l'affirmation est \textbf{fausse}. Rappel : quand le pH diminue de 1 unité, $[\ce{H3O+}]$ est multipliée par 10 ; ici $\Delta\text{pH} = 2{,}70$, soit un facteur $10^{2{,}70} \approx 500$.
\end{exoresolu}

\courssub{Calculs de pH et concentrations ioniques}

\textbf{Méthode 4 : Calculer le pH d'une solution et en déduire son caractère}\par
\begin{methode}[Calculer un pH et conclure (acides/bases forts, solutions diluées)]
\begin{enumerate}
\item \textbf{Identifier la nature du soluté} : acide fort (réaction totale avec l'eau : \ce{HCl}, \ce{HNO3}, \ce{H2SO4} 1er proton), base forte (hydroxydes alcalins \ce{NaOH}, \ce{KOH}, \ce{Ba(OH)2}), ou autre (acide/base faible, sel).
\item \textbf{Pour un acide fort} de concentration molaire $C$ (avec $10^{-6} \ll C \ll 1\ \si{mol.L^{-1}}$ pour négliger l'autoprotolyse et les effets d'activité) :
\[ \ce{HA + H2O -> H3O+ + A-} \quad (\text{réaction totale}) \implies [\ce{H3O+}] = C \]
\[ \boxed{\text{pH} = -\log C} \]
\item \textbf{Pour une base forte} de concentration molaire $C$ (mêmes conditions de validité) :
\[ \ce{BOH -> B+ + HO-} \quad (\text{dissociation totale}) \implies [\ce{HO-}] = C \]
\[ [\ce{H3O+}] = \frac{K_e}{[\ce{HO-}]} = \frac{K_e}{C} \implies \boxed{\text{pH} = pK_e + \log C = 14 + \log C \ \text{à}\ \SI{25}{\celsius}} \]
\item \textbf{Conclure sur le caractère} (à \SI{25}{\celsius}) :
\begin{itemize}
\item pH $< 7$ : solution \textbf{acide} ($[\ce{H3O+}] > [\ce{HO-}]$) ;
\item pH $= 7$ : solution \textbf{neutre} ;
\item pH $> 7$ : solution \textbf{basique} ($[\ce{H3O+}] < [\ce{HO-}]$).
\end{itemize}
\item \textbf{Vérifier la cohérence} : le pH d'un acide dilué ne peut pas dépasser 7 (il tend vers 7 par dilution), celui d'une base ne peut pas descendre sous 7.
\end{enumerate}
\end{methode}

\begin{exoresolu}[pH d'un détartrant et d'une solution de soude]
Un détartrant ménager contient une solution d'acide chlorhydrique (\ce{H3O+} + \ce{Cl-}) de concentration $C_1 = 2{,}0 \times 10^{-2}\ \si{mol.L^{-1}}$. Une solution de soude \ce{NaOH} (base forte) a pour concentration $C_2 = 5{,}0 \times 10^{-3}\ \si{mol.L^{-1}}$. On se place à \SI{25}{\celsius}.
Données : $\log 2 = 0{,}30$ ; $\log 5 = 0{,}70$.
\begin{enumerate}
\item Calculer le pH du détartrant.
\item Calculer le pH de la solution de soude.
\item Conclure sur le caractère de chaque solution.
\end{enumerate}
\tcblower
\textbf{1. pH du détartrant (acide fort).}

L'acide chlorhydrique est un acide fort, totalement dissocié dans l'eau :
\[ \ce{HCl + H2O -> H3O+ + Cl-} \]
donc $[\ce{H3O+}] = C_1 = 2{,}0 \times 10^{-2}\ \si{mol.L^{-1}}$.
\[ \text{pH}_1 = -\log[\ce{H3O+}] = -\log(2{,}0 \times 10^{-2}) = 2 - \log 2 = 2 - 0{,}30 = 1{,}70 \]

\textbf{2. pH de la soude (base forte).}

La soude est totalement dissociée : $\ce{NaOH -> Na+ + HO-}$, donc $[\ce{HO-}] = C_2 = 5{,}0 \times 10^{-3}\ \si{mol.L^{-1}}$.
\[ [\ce{H3O+}] = \frac{K_e}{[\ce{HO-}]} = \frac{1{,}0 \times 10^{-14}}{5{,}0 \times 10^{-3}} = 2{,}0 \times 10^{-12}\ \si{mol.L^{-1}} \]
\[ \text{pH}_2 = -\log(2{,}0 \times 10^{-12}) = 12 - \log 2 = 11{,}70 \]
Vérification par la formule directe : $\text{pH}_2 = 14 + \log(5{,}0 \times 10^{-3}) = 14 - 3 + 0{,}70 = 11{,}70$. \checkmark

\textbf{3. Conclusion.}

$\text{pH}_1 = 1{,}70 < 7$ : le détartrant est \textbf{fortement acide}. $\text{pH}_2 = 11{,}70 > 7$ : la soude est \textbf{fortement basique}. Ces deux produits sont corrosifs : port de gants, lunettes et blouse obligatoires en TP.
\end{exoresolu}

\textbf{Méthode 5 : Trouver les concentrations ioniques à partir du pH}\par
\begin{methode}[Remonter du pH aux concentrations et exploiter l'électroneutralité]
\begin{enumerate}
\item \textbf{Calculer $[\ce{H3O+}]$} par la définition inverse :
\[ [\ce{H3O+}] = 10^{-\text{pH}}\ \si{mol.L^{-1}} \]
\item \textbf{Calculer $[\ce{HO-}]$} grâce au produit ionique :
\[ [\ce{HO-}] = \frac{K_e}{[\ce{H3O+}]} = 10^{\text{pH} - pK_e}\ \si{mol.L^{-1}} \]
\item \textbf{Faire les approximations} (hiérarchie des espèces, pour concentrations usuelles $> 10^{-3}\ \si{mol.L^{-1}}$) :
\begin{itemize}
\item solution acide (pH $< 6$) : \ce{HO-} \textbf{ultra-minoritaire} ($< 10^{-8}\ \si{mol.L^{-1}}$), on le néglige dans l'électroneutralité ;
\item solution basique (pH $> 8$) : \ce{H3O+} \textbf{ultra-minoritaire} ($< 10^{-8}\ \si{mol.L^{-1}}$), on le néglige ;
\item entre pH 6 et 8 (proche neutralité) : les deux ions de l'eau peuvent être négligeables \textbf{devant les ions du soluté} si ce dernier est concentré ($C \gg 10^{-6}\ \si{mol.L^{-1}}$).
\end{itemize}
\item \textbf{Écrire l'électroneutralité} de la solution : somme des charges positives = somme des charges négatives, pour retrouver les concentrations des ions spectateurs (ions du sel).
\item \textbf{Vérifier} : $[\ce{H3O+}][\ce{HO-}] = 10^{-14}$ à \SI{25}{\celsius}.
\end{enumerate}
\end{methode}

\begin{exoresolu}[Analyse d'une eau de source du Mont Cameroun]
Une eau de source prélevée sur les flancs du Mont Cameroun a un pH de 6,20 à \SI{25}{\celsius}. Elle contient notamment des ions sodium \ce{Na+} et chlorure \ce{Cl-}, avec $[\ce{Na+}] = 1{,}0 \times 10^{-4}\ \si{mol.L^{-1}}$.
\begin{enumerate}
\item Calculer $[\ce{H3O+}]$ et $[\ce{HO-}]$.
\item Montrer que les ions \ce{HO-} sont négligeables devant \ce{Na+}.
\item En supposant que seuls \ce{Na+}, \ce{Cl-}, \ce{H3O+} et \ce{HO-} sont présents, calculer $[\ce{Cl-}]$.
\end{enumerate}
Donnée : $10^{-6{,}20} = 6{,}3 \times 10^{-7}$.
\tcblower
\textbf{1. Concentrations des ions de l'eau.}
\[ [\ce{H3O+}] = 10^{-\text{pH}} = 10^{-6{,}20} = 6{,}3 \times 10^{-7}\ \si{mol.L^{-1}} \]
\[ [\ce{HO-}] = \frac{K_e}{[\ce{H3O+}]} = \frac{1{,}0 \times 10^{-14}}{6{,}3 \times 10^{-7}} = 1{,}6 \times 10^{-8}\ \si{mol.L^{-1}} \]

\textbf{2. Approximation.}

Comparons : $\dfrac{[\ce{Na+}]}{[\ce{HO-}]} = \dfrac{1{,}0 \times 10^{-4}}{1{,}6 \times 10^{-8}} \approx 6\,250$.

$[\ce{HO-}]$ est environ 6000 fois plus petite que $[\ce{Na+}]$ : les ions hydroxyde sont \textbf{ultra-minoritaires} et négligeables dans l'électroneutralité. (De même $[\ce{H3O+}] = 6{,}3 \times 10^{-7} \ll [\ce{Na+}]$, facteur \approx 160, minoritaire).

\textbf{3. Concentration en chlorure.}

Électroneutralité de la solution :
\[ [\ce{Na+}] + [\ce{H3O+}] = [\ce{Cl-}] + [\ce{HO-}] \]
En négligeant \ce{HO-} (ultra-minoritaire) et \ce{H3O+} (minoritaire) devant \ce{Na+} et \ce{Cl-} :
\[ [\ce{Cl-}] \approx [\ce{Na+}] + [\ce{H3O+}] = 1{,}0 \times 10^{-4} + 6{,}3 \times 10^{-7} = 1{,}0063 \times 10^{-4}\ \si{mol.L^{-1}} \]
Soit $[\ce{Cl-}] \approx 1{,}0 \times 10^{-4}\ \si{mol.L^{-1}}$ à la précision des mesures (2 chiffres significatifs sur $[\ce{Na+}]$).

\textbf{Conclusion} : cette eau légèrement acide (pH 6,20) contient $[\ce{H3O+}] = 6{,}3 \times 10^{-7}\ \si{mol.L^{-1}}$ et $[\ce{HO-}] = 1{,}6 \times 10^{-8}\ \si{mol.L^{-1}}$ ; les ions du soluté (\ce{Na+}, \ce{Cl-}) dominent largement.
\end{exoresolu}

\courssub{Préparation de solutions par dilution}

\textbf{Méthode 6 : Réaliser une dilution}\par
\begin{methode}[Réaliser une dilution (solution fille à partir d'une solution mère)]
\begin{enumerate}
\item \textbf{Écrire la conservation de la quantité de matière} du soluté prélevé :
\[ \boxed{C_0 V_0 = C_1 V_1} \]
où $C_0$, $V_0$ : concentration et volume de la solution mère (prélevée) ; $C_1$, $V_1$ : concentration et volume de la solution fille (préparée).
\item \textbf{Calculer le volume à prélever} : $V_0 = \dfrac{C_1 V_1}{C_0}$. Définir le \cle{facteur de dilution} $F = \dfrac{C_0}{C_1} = \dfrac{V_1}{V_0}$.
\item \textbf{Protocole rigoureux} :
\begin{itemize}
\item Prélever $V_0$ à la \textbf{pipette jaugée} (munie d'une \textbf{propipette} ou poire à pipeter, \textbf{jamais à la bouche} : risque chimique) après rinçage de la pipette avec la solution mère ;
\item Verser dans une \textbf{fiole jaugée} de volume $V_1$ remplie au tiers d'eau distillée ;
\item Compléter avec de l'eau distillée jusqu'au \textbf{trait de jauge} : le bas du ménisque doit être tangent au trait, l'œil à hauteur du trait (éviter l'erreur de parallaxe) ;
\item Boucher et \textbf{homogénéiser} en retournant la fiole une dizaine de fois.
\end{itemize}
\item \textbf{Effet sur le pH} : diluer 10 fois un acide fort augmente le pH de 1 unité ; diluer 10 fois une base forte le diminue de 1 unité. \textbf{Jamais} une dilution ne peut rendre un acide basique ni une base acide (le pH tend vers 7).
\end{enumerate}
\end{methode}

\begin{exoresolu}[Préparation d'une solution pour le laboratoire du lycée]
Au laboratoire d'un lycée de Douala, on dispose d'une solution mère d'acide chlorhydrique de concentration $C_0 = 1{,}0 \times 10^{-1}\ \si{mol.L^{-1}}$. On veut préparer $V_1 = \SI{250,0}{mL}$ de solution fille de concentration $C_1 = 2{,}0 \times 10^{-3}\ \si{mol.L^{-1}}$.
\begin{enumerate}
\item Calculer le volume $V_0$ de solution mère à prélever et le facteur de dilution.
\item Décrire le protocole en précisant la verrerie.
\item Calculer le pH de la solution mère et celui de la solution fille à \SI{25}{\celsius}. Conclure.
\end{enumerate}
\tcblower
\textbf{1. Volume à prélever et facteur de dilution.}

Conservation de la quantité de matière : $C_0 V_0 = C_1 V_1$, d'où :
\[ V_0 = \frac{C_1 V_1}{C_0} = \frac{2{,}0 \times 10^{-3} \times 250{,}0}{1{,}0 \times 10^{-1}} = 5{,}0\ \si{mL} \]
Facteur de dilution : $F = \dfrac{C_0}{C_1} = \dfrac{1{,}0 \times 10^{-1}}{2{,}0 \times 10^{-3}} = 50$. Vérification : $F = \dfrac{V_1}{V_0} = \dfrac{250{,}0}{5{,}0} = 50$. \checkmark

\textbf{2. Protocole.}
\begin{itemize}
\item Prélever $V_0 = \SI{5,0}{mL}$ de solution mère avec une \textbf{pipette jaugée de \SI{5}{mL}} munie d'une propipette (jamais à la bouche : acide corrosif).
\item Les verser dans une \textbf{fiole jaugée de \SI{250,0}{mL}} contenant déjà un tiers d'eau distillée.
\item Compléter à l'eau distillée jusqu'au trait de jauge (ménisque tangent au trait, œil à hauteur).
\item Boucher, retourner une dizaine de fois pour homogénéiser.
\end{itemize}

\textbf{3. pH des deux solutions.}

Acide chlorhydrique = acide fort : pH $= -\log C$.
\begin{itemize}
\item Solution mère : $\text{pH}_0 = -\log(1{,}0 \times 10^{-1}) = 1{,}00$.
\item Solution fille : $\text{pH}_1 = -\log(2{,}0 \times 10^{-3}) = 3 - \log 2 = 2{,}70$.
\end{itemize}
\textbf{Conclusion} : la dilution par 50 a fait passer le pH de 1,00 à 2,70 ($\Delta\text{pH} = \log 50 = 1{,}70$). La solution fille reste acide : diluer un acide ne peut jamais le rendre neutre ou basique.
\end{exoresolu}

\begin{attention}[Les 5 erreurs qui coûtent des points au Bac]
\begin{enumerate}
\item \textbf{Confondre neutralité et pH = 7.} pH = 7 n'est la neutralité qu'à \SI{25}{\celsius}. La vraie définition : $[\ce{H3O+}] = [\ce{HO-}]$. À \SI{37}{\celsius}, le sang (pH 7,40) est légèrement basique car $pK_e \approx 13{,}6$ (neutralité à pH 6,80).
\item \textbf{Oublier $K_e$ pour les bases.} Pour une base forte, écrire directement pH $= -\log C$ est faux : il faut passer par $[\ce{HO-}] = C$ puis $[\ce{H3O+}] = K_e/[\ce{HO-}]$, soit pH $= 14 + \log C$ à \SI{25}{\celsius}.
\item \textbf{Écrire \ce{H+} libre en solution aqueuse.} Le proton n'existe pas seul dans l'eau : il faut écrire \ce{H3O+}. L'équation de l'autoprotolyse est $\ce{2 H2O <=> H3O+ + HO-}$, pas $\ce{H2O <=> H+ + HO-}$ (tolérée en Terminale mais moins rigoureuse).
\item \textbf{Se tromper de sens dans la dilution.} Diluer un acide fait \emph{augmenter} son pH (vers 7), diluer une base fait \emph{diminuer} son pH (vers 7). Un pH d'acide dilué ne dépasse jamais 7 !
\item \textbf{Négliger les unités et les conversions dans $C_0V_0 = C_1V_1$.} Les volumes doivent être dans la \textbf{même unité} (mL et mL, ou L et L) ; les concentrations en \si{mol.L^{-1}}. Et une concentration issue d'un pH se donne avec au plus 2 chiffres significatifs.
\end{enumerate}
\end{attention}

\newpage

\coursec{Exercices}

\coursec{Acides et bases en solution aqueuse : eau, autoprotolyse et pH}

\courssub{L'eau : structure et pouvoir solvant}

\begin{definition}[Molécule d'eau]
La molécule d'eau \ce{H2O} possède une géométrie \textbf{coudée} (angle \ce{H-O-H} $\approx \SI{104,5}{\degree}$). L'atome d'oxygène, plus \emph{électronégatif} que l'hydrogène, attire le nuage électronique des liaisons \ce{O-H} : la molécule est \textbf{polaire}, porteuse d'un \textbf{moment dipolaire} permanent. L'oxygène porte une charge partielle négative ($\delta^-$), les hydrogènes des charges partielles positives ($\delta^+$).
\end{definition}

\begin{popfigure}
\begin{tikzpicture}[scale=1.2, every node/.style={font=\small}]
% Molécule d'eau
\coordinate (O) at (0,0);
\coordinate (H1) at (0.95,0.55);
\coordinate (H2) at (0.95,-0.55);
% Liaisons
\draw[thick, popdark] (O) -- (H1) node[midway, above left, popdark] {$\delta^+$};
\draw[thick, popdark] (O) -- (H2) node[midway, below left, popdark] {$\delta^+$};
% Atomes
\node[circle, draw=popblue, thick, fill=popblue!20, minimum size=1.2cm] at (O) {\textbf{O}};
\node[circle, draw=poporange, thick, fill=poporange!20, minimum size=0.7cm] at (H1) {\textbf{H}};
\node[circle, draw=poporange, thick, fill=poporange!20, minimum size=0.7cm] at (H2) {\textbf{H}};
% Dipôle
\draw[->, very thick, poppurple] (0,-1.2) -- (0,0.8) node[midway, right] {$\vec{\mu}$ (moment dipolaire)};
% Angle
\draw[popmuted, <->] (0.3,0.2) arc (30:-30:0.3) node[midway, right=2pt] {$104,5^\circ$};
% Nuage électronique
\draw[popblue, dashed, decorate, decoration={snake, amplitude=1pt, segment length=4pt}] (-0.6,0) -- (-1.2,0) node[left] {Densité électronique};
\end{tikzpicture}
\legende{Structure de la molécule d'eau : géométrie coudée et polarité.}
\end{popfigure}

\begin{propriete}[Liaisons hydrogène et pouvoir solvant]
Dans l'eau liquide, les molécules s'organisent en un \textbf{réseau de liaisons hydrogène} : l'atome d'hydrogène ($\delta^+$) d'une molécule est attiré par l'atome d'oxygène ($\delta^-$) d'une voisine. Ce réseau explique :
\begin{itemize}
    \item Les points d'ébullition et de fusion anormalement élevés.
    \item Le \textbf{pouvoir solvant exceptionnel} : l'eau solubilise les composés ioniques (solvatation des ions par orientation des dipôles) et les molécules polaires (alcools, sucres, acides). C'est le \emph{solvant universel} de la chimie aqueuse.
\end{itemize}
\end{propriete}

\begin{savaistu}[L'eau au Cameroun]
L'eau de source (Mont Cameroun), l'eau de forage (Ngaoundéré) ou l'eau de rivière contiennent des ions dissous (\ce{Ca^2+}, \ce{Mg^2+}, \ce{HCO3-}, \ce{Cl-}) qui modifient leur conductivité et leur pH. L'analyse de ces paramètres est cruciale pour la potabilité (MINSANTÉ) et l'irrigation (MINADER).
\end{savaistu}

\courssub{Acides et bases selon Brønsted}

\begin{definition}[Acide et base de Brønsted]
\begin{itemize}
    \item Un \textbf{acide de Brønsted} est une espèce chimique capable de \textbf{céder un proton} \ce{H+}.
    \item Une \textbf{base de Brønsted} est une espèce chimique capable de \textbf{capter un proton} \ce{H+}.
\end{itemize}
La réaction acide-base est un \textbf{transfert de proton} :
\[ \ce{AH + B <=> A- + BH+} \]
\ce{AH} (acide) et \ce{A-} (base conjuguée) forment un \textbf{couple acide/base} noté \ce{AH}/\ce{A-}. \ce{B} (base) et \ce{BH+} (acide conjugué) forment le couple \ce{BH+}/\ce{B}.
\end{definition}

\begin{exemplebox}[Exemples de couples acide/base]
\begin{center}
\begin{tabularx}{\linewidth}{|l|X|X|}
\hline
\rowcolor{popblueL} \textbf{Acide} & \textbf{Base conjuguée} & \textbf{Couple} \\
\hline
\ce{HCl} (acide chlorhydrique) & \ce{Cl-} (ion chlorure) & \ce{HCl}/\ce{Cl-} \\
\ce{CH3COOH} (acide éthanoïque) & \ce{CH3COO-} (ion éthanoate) & \ce{CH3COOH}/\ce{CH3COO-} \\
\ce{H2O} (eau) & \ce{HO-} (ion hydroxyde) & \ce{H2O}/\ce{HO-} \\
\ce{H3O+} (ion hydronium/oxonium) & \ce{H2O} (eau) & \ce{H3O+}/\ce{H2O} \\
\ce{H2CO3} (acide carbonique) & \ce{HCO3-} (ion hydrogénocarbonate) & \ce{H2CO3}/\ce{HCO3-} \\
\ce{HCO3-} (ion hydrogénocarbonate) & \ce{CO3^2-} (ion carbonate) & \ce{HCO3-}/\ce{CO3^2-} \\
\hline
\end{tabularx}
\end{center}
\end{exemplebox}

\begin{definition}[Ampholyte (espèce amphotère)]
Une espèce qui peut agir \textbf{à la fois comme acide et comme base} est un \textbf{ampholyte}. L'eau est l'exemple type : elle cède un proton à une base forte (rôle d'acide) et capte un proton d'un acide fort (rôle de base).
\[ \ce{H2O + H2O <=> H3O+ + HO-} \]
D'autres ampholytes courants : \ce{HCO3-}, \ce{HSO4-}, \ce{H2PO4-}, \ce{NH3} (dans l'eau liquide).
\end{definition}

\begin{attention}[Piège : \ce{H3O+} vs \ce{H+}]
En solution aqueuse, le proton \ce{H+} n'existe pas \emph{nu} : il est instantanément solvaté par l'eau pour former l'\textbf{ion hydronium} (ou oxonium) \ce{H3O+}. On écrit \ce{H3O+} dans les équations-bilan. L'écriture \ce{H+} est une abréviation commode pour les calculs de pH.
\end{attention}

\courssub{Autoprotolyse de l'eau et produit ionique $K_e$}

\begin{experience}[Autoprotolyse de l'eau pure]
\textbf{Dispositif :} Bécher d'eau distillée, électrodes de platine, générateur basse tension, ampoule.
\textbf{Observation :} L'ampoule ne s'allume pas (ou très faiblement). L'eau pure conduit très mal le courant.
\textbf{Interprétation :} L'eau pure contient très peu d'ions. Elle est le siège d'une réaction limitée : l'\textbf{autoprotolyse}. Une molécule d'eau (acide) cède un proton à une autre molécule d'eau (base).
\end{experience}

\begin{propriete}[Équation-bilan et produit ionique de l'eau]
À la température $\theta$ (généralement \SI{25}{\celsius}) :
\[ \ce{2 H2O(l) <=> H3O+(aq) + HO-(aq)} \]
Le \textbf{produit ionique de l'eau} $K_e$ (sans unité) est :
\[ K_e = [\ce{H3O+}] \times [\ce{HO-}] \]
À \SI{25}{\celsius} : \quad $K_e = \SI{1,0e-14}{}$ \quad ($pK_e = 14$).
\end{propriete}

\begin{aretenir}[Variation de $K_e$ avec la température]
L'autoprotolyse est \textbf{endothermique} ($\Delta H > 0$). Selon le principe de Le Chatelier, \textbf{$K_e$ augmente avec la température}.
\begin{center}
\begin{tabularx}{\linewidth}{|c|c|c|c|}
\hline
\rowcolor{popblueL} $\theta$ (\si{\celsius}) & $K_e$ & $pK_e$ & pH neutre ($=\frac{pK_e}{2}$) \\
\hline
$0$ & $0,11 \times 10^{-14}$ & $14,9$ & $7,47$ \\
$25$ & $1,0 \times 10^{-14}$ & $14,0$ & $7,00$ \\
$37$ (corps humain) & $2,5 \times 10^{-14}$ & $13,6$ & $6,80$ \\
$100$ & $51 \times 10^{-14}$ & $12,3$ & $6,15$ \\
\hline
\end{tabularx}
\end{center}
\end{aretenir}

\begin{propriete}[Solution neutre, acide, basique]
Dans \textbf{toute} solution aqueuse (même contenant d'autres solutés), le produit $[\ce{H3O+}][\ce{HO-}]$ vaut $K_e$ à la température considérée.
\begin{itemize}
    \item \textbf{Solution neutre :} $[\ce{H3O+}] = [\ce{HO-}] = \sqrt{K_e}$.
    \item \textbf{Solution acide :} $[\ce{H3O+}] > [\ce{HO-}]$ (donc $[\ce{H3O+}] > \sqrt{K_e}$).
    \item \textbf{Solution basique :} $[\ce{H3O+}] < [\ce{HO-}]$ (donc $[\ce{H3O+}] < \sqrt{K_e}$).
\end{itemize}
\end{propriete}

\courssub{Le pH : définition et mesure}

\begin{definition}[pH d'une solution aqueuse]
Le \textbf{pH} (puissance hydrogène) est défini par :
\[ \mathrm{pH} = -\log_{10}\left(\frac{[\ce{H3O+}]}{c^\circ}\right) \quad \text{avec } c^\circ = \SI{1}{mol.L^{-1}} \]
En pratique, on utilise la concentration numérique : $\mathrm{pH} = -\log_{10}[\ce{H3O+}]$ (avec $[\ce{H3O+}]$ en \si{mol.L^{-1}}).
Réciproquement : \quad $[\ce{H3O+}] = 10^{-\mathrm{pH}}$ \quad (en \si{mol.L^{-1}}).
\end{definition}

\begin{methode}[Calculer le pH ou [\ce{H3O+}]]
\begin{enumerate}
    \item Identifier la donnée : $[\ce{H3O+}]$ ou pH.
    \item Appliquer la formule : $\mathrm{pH} = -\log[\ce{H3O+}]$ ou $[\ce{H3O+}] = 10^{-\mathrm{pH}}$.
    \item \textbf{Vérifier l'ordre de grandeur :} pH entre 0 et 14 (généralement) ; $[\ce{H3O+}]$ entre 1 et $10^{-14}$.
    \item Arrondir le pH à \textbf{deux décimales} (précision du pH-mètre) ; $[\ce{H3O+}]$ à \textbf{deux chiffres significatifs}.
\end{enumerate}
\end{methode}

\begin{exemplebox}[Calculs directs]
\begin{itemize}
    \item $[\ce{H3O+}] = \SI{2,5e-4}{mol.L^{-1}} \Rightarrow \mathrm{pH} = -\log(2,5\times10^{-4}) = 4 - \log 2,5 \approx 4 - 0,40 = \mathbf{3,60}$.
    \item $\mathrm{pH} = 9,20 \Rightarrow [\ce{H3O+}] = 10^{-9,20} = 10^{-10} \times 10^{0,80} \approx \mathbf{6,3 \times 10^{-10}}\ \si{mol.L^{-1}}$.
\end{itemize}
\end{exemplebox}

\begin{propriete}[Échelle de pH à \SI{25}{\celsius}]
\begin{center}
\begin{tikzpicture}
\draw[popline, thick] (0,0) -- (14,0);
\foreach \x in {0,1,...,14} \draw (\x,0.1) -- (\x,-0.1) node[below] {\x};
% Zones
\fill[popred!30] (0,0) rectangle (7,-0.6) node[midway, popdark] {ACIDE};
\fill[popgreen!30] (7,0) rectangle (14,-0.6) node[midway, popdark] {BASIQUE};
\draw[thick, popdark] (7,0.1) -- (7,-0.7) node[midway, right] {NEUTRE (7,0)};
% Exemples Cameroun
\node[above=0.5cm, popblue] at (3,0.3) {Jus de bissap $\approx 3,2$};
\node[above=0.5cm, popblue] at (4.5,0.3) {Bière $\approx 4,5$};
\node[above=0.5cm, popblue] at (6.8,0.3) {Eau de pluie $\approx 5,8$};
\node[above=0.5cm, poporange] at (7.4,0.3) {Sang $\approx 7,4$};
\node[above=0.5cm, poporange] at (8.5,0.3) {Eau de mer $\approx 8,1$};
\node[above=0.5cm, poporange] at (11,0.3) {Lessive (soude) $\approx 11-12$};
\node[above=0.5cm, poporange] at (13,0.3) {Eau de Javel $\approx 12,5$};
\end{tikzpicture}
\end{center}
\end{propriete}

\begin{propriete}[Méthodes de mesure du pH]
\begin{center}
\begin{tabularx}{\linewidth}{|l|X|X|}
\hline
\rowcolor{popblueL} \textbf{Méthode} & \textbf{Principe} & \textbf{Précision / Usage} \\
\hline
\textbf{Papier pH} (bandelette) & Changement de couleur d'un mélange d'indicateurs imprégné sur papier. & $\pm 0,5$ à $1$ unité. \emph{Usage :} contrôle rapide, terrain, TP. \\
\hline
\textbf{Indicateurs colorés} (solution) & Couleur acide $\neq$ couleur basique. Zone de virage $\mathrm{pH} = pK_a \pm 1$. & $\pm 0,2$ à $0,5$ unité. \emph{Usage :} repérage, dosage (équivalence). \\
\hline
\textbf{pH-mètre} (électrode de verre) & Mesure de la différence de potentiel $E = E^\circ - \frac{RT}{F}\ln[\ce{H3O+}]$. & $\pm 0,01$ à $0,05$ unité. \emph{Usage :} analyses précises, médicales, industrielles (SONARA, CIMENCAM, brasseries). \\
\hline
\end{tabularx}
\end{center}
\end{propriete}

\begin{attention}[Choix de l'indicateur]
Pour un dosage acide-base, l'indicateur doit avoir sa \textbf{zone de virage incluse dans le saut de pH à l'équivalence}. Exemple : Phénolphtaléine (virage 8,2--10, rose incolore) pour dosage acide fort/base forte ; Hélianthine (3,1--4,4) pour acide faible/base forte.
\end{attention}

\courssub{Calculs de pH et concentrations ioniques : approximations}

\begin{methode}[Détermination systématique des concentrations ioniques]
Pour une solution aqueuse à température connue ($K_e$ connu) :
\begin{enumerate}
    \item \textbf{Inventaire des espèces :} Lister toutes les entités (molécules, ions) présentes.
    \item \textbf{Équation d'électroneutralité :} Somme des charges positives = Somme des charges négatives.
    \item \textbf{Relations de conservation / dissociation :} $K_e$, formules de dissociation, stœchiométrie.
    \item \textbf{Approximations (ordre de grandeur) :}
    \begin{itemize}
        \item \textbf{Majoritaire :} $C > \SI{e-3}{mol.L^{-1}}$.
        \item \textbf{Minoritaire :} $\SI{e-6}{} < C < \SI{e-3}{mol.L^{-1}}$.
        \item \textbf{Ultra-minoritaire :} $C < \SI{e-6}{mol.L^{-1}}$.
    \end{itemize}
    \item \textbf{Résolution :} Négocier les termes négligeables dans l'électroneutralité.
    \item \textbf{Vérification :} Contrôler que les approximations sont valides a posteriori.
\end{enumerate}
\end{methode}

\begin{exoresolu}[Inventaire et approximations : Solution de NaCl $\SI{1,0e-2}{M}$]
\textbf{Énoncé :} On dissout du chlorure de sodium \ce{NaCl} dans l'eau distillée ; on obtient une solution de concentration $C = \SI{1,0e-2}{mol.L^{-1}}$ à \SI{25}{\celsius}.
\begin{enumerate}
    \item \textbf{Inventaire :} \ce{Na+}, \ce{Cl-}, \ce{H2O}, \ce{H3O+}, \ce{HO-}. (\ce{NaCl} dissocié totalement).
    \item \textbf{Électroneutralité :} $[\ce{Na+}] + [\ce{H3O+}] = [\ce{Cl-}] + [\ce{HO-}]$.
    \item \textbf{Justification $[\ce{H3O+}] = [\ce{HO-}]$ :} $[\ce{Na+}] = [\ce{Cl-}] = C = \SI{1,0e-2}{M}$. L'électroneutralité donne $[\ce{H3O+}] = [\ce{HO-}]$. La solution est \textbf{neutre}.
    \item \textbf{Valeurs :} $[\ce{H3O+}] = [\ce{HO-}] = \sqrt{K_e} = \SI{1,0e-7}{mol.L^{-1}}$. $\mathrm{pH} = 7,00$.
    \item \textbf{Classification :}
    \begin{itemize}
        \item Majoritaires : \ce{H2O} ($\approx \SI{55}{M}$), \ce{Na+}, \ce{Cl-} ($10^{-2}$ M).
        \item Ultra-minoritaires : \ce{H3O+}, \ce{HO-} ($10^{-7}$ M).
    \end{itemize}
\end{enumerate}
\end{exoresolu}

\courssub{Préparation de solutions par dilution}

\begin{propriete}[Loi de dilution]
Lors d'une dilution, la \textbf{quantité de matière de soluté est conservée} :
\[ n = C_0 V_0 = C_1 V_1 \]
$C_0, V_0$ : concentration et volume de la solution mère (prélevée).
$C_1, V_1$ : concentration et volume de la solution fille (préparée).
Le \textbf{facteur de dilution} est $f = \frac{C_0}{C_1} = \frac{V_1}{V_0}$.
\end{propriete}

\begin{methode}[Protocole de dilution au dixième ($f=10$)]
\textbf{Matériel :} Pipette jaugée \SI{10,0}{mL} (ou \SI{20,0}{mL}, \SI{25,0}{mL}...), fiole jaugée \SI{100}{mL} (ou \SI{200}{mL}, \SI{250}{mL}...), propipette, eau distillée, bécher.
\begin{enumerate}
    \item Rincer la pipette avec la solution mère $S_0$ (2--3 fois).
    \item Prélever le volume $V_0$ (ex: \SI{10,0}{mL}) avec la propipette. \textbf{Ne pas souffler} la dernière goutte (calibrée "à livrer").
    \item Verser dans la fiole jaugée $V_1$ (ex: \SI{100}{mL}) rincée à l'eau distillée.
    \item Rincer la pipette au-dessus de la fiole avec un peu d'eau distillée (récupérer les gouttes).
    \item Compléter à l'eau distillée jusqu'au trait de jauge (bas du ménisque au niveau du trait).
    \item Bouchonner et homogénéiser (retourner la fiole plusieurs fois).
\end{enumerate}
\end{methode}

\begin{attention}[Sécurité : Acides et bases concentrés]
\begin{itemize}
    \item \textbf{Toujours verser l'acide (ou la base) dans l'eau}, jamais l'inverse (éclaboussures, projection par chaleur de dilution).
    \item Porter \textbf{lunettes de protection}, \textbf{gants}, \textbf{blouse}.
    \item En cas de projection : rincer abondamment à l'eau \textbf{pendant 15 min minimum}, consulter.
    \item Manipuler les solutions concentrées (acide chlorhydrique \SI{12}{M}, soude \SI{10}{M}) sous \textbf{hotte aspirante} (vapeurs irritantes).
\end{itemize}
\end{attention}

\begin{savaistu}[Dilution industrielle à la SONARA / Brasseries]
À la SONARA (Société Nationale de Raffinage), les solutions étalons pour étalonner les analyseurs (pH-mètres, conductimètres) sont préparées par dilutions successives rigoureuses à partir de solutions mères certifiées (NIST). Une erreur de dilution fausse tout le contrôle qualité des carburants. De même, dans les brasseries (Isangui, Ngok), le contrôle du pH du moût et de la bière nécessite des solutions tampons étalonnées par dilution précise.
\end{savaistu}

% ============================================================
% EXERCICES
% ============================================================

\courssub{Je vérifie mes connaissances}

\begin{exercice}[QCM : à toi de choisir]
Pour chaque question, une seule réponse est exacte. Recopie la lettre correspondante et la réponse complète.
\begin{enumerate}
\item À \SI{25}{\celsius}, le produit ionique de l'eau vaut :
\begin{enumerate}
\item $K_e = 10^{-7}$ \quad b) $K_e = 10^{-14}$ \quad c) $K_e = 14$ \quad d) $K_e = 1$
\end{enumerate}
\item Une solution dont le pH vaut $9{,}0$ à \SI{25}{\celsius} est :
\begin{enumerate}
\item acide \quad b) neutre \quad c) basique \quad d) impossible à classer
\end{enumerate}
\item Si $[\ce{H3O+}] = \SI{2,0e-4}{mol.L^{-1}}$, alors le pH vaut environ :
\begin{enumerate}
\item $3{,}0$ \quad b) $3{,}7$ \quad c) $4{,}0$ \quad d) $4{,}3$
\end{enumerate}
\item Dans l'autoprotolyse de l'eau, la molécule d'eau qui capte un proton joue le rôle de :
\begin{enumerate}
\item acide de Brønsted \quad b) base de Brønsted \quad c) oxydant \quad d) solvant inerte
\end{enumerate}
\item Pour diluer une solution au dixième, on prélève \SI{10,0}{mL} de solution mère que l'on introduit dans une fiole jaugée de :
\begin{enumerate}
\item \SI{50}{mL} \quad b) \SI{100}{mL} \quad c) \SI{250}{mL} \quad d) \SI{1000}{mL}
\end{enumerate}
\end{enumerate}
\end{exercice}

\begin{exercice}[Vrai ou faux ? Justifie]
Réponds par VRAI ou FAUX, puis justifie ta réponse en une ou deux phrases.
\begin{enumerate}
\item Une solution aqueuse neutre a toujours un pH égal à $7$, quelle que soit la température.
\item Dans toute solution aqueuse, on a toujours $[\ce{H3O+}] \times [\ce{HO-}] = K_e$.
\item L'eau est un ampholyte (espèce amphotère).
\item Quand on dilue dix fois une solution d'acide chlorhydrique, son pH diminue d'une unité.
\item Le papier pH permet une mesure précise du pH au centième d'unité près.
\end{enumerate}
\end{exercice}

\begin{exercice}[Texte à trous]
Recopie et complète le texte suivant avec les mots ou expressions de la liste : \emph{proton ; autoprotolyse ; $10^{-14}$ ; hydronium ; hydroxyde ; augmente ; diminue ; ampholyte ; électroneutralité}.

L'eau pure conduit très faiblement le courant électrique car elle est le siège d'une réaction limitée appelée \dotfill de l'eau. Au cours de cette réaction, une molécule d'eau cède un \dotfill à une autre molécule d'eau : il se forme alors l'ion \dotfill{} \ce{H3O+} et l'ion \dotfill{} \ce{HO-}. L'eau se comporte à la fois comme un acide et comme une base : on dit que c'est un \dotfill. À \SI{25}{\celsius}, le produit ionique de l'eau vaut $K_e = $ \dotfill. Quand on ajoute un acide à de l'eau, la concentration $[\ce{H3O+}]$ \dotfill{} et le pH \dotfill. Toute solution aqueuse vérifie la loi d'\dotfill : la somme des charges positives est égale à la somme des charges négatives.
\end{exercice}

\begin{exercice}[Définitions et formules]
\begin{enumerate}
\item Définis un acide et une base selon Brønsted. Donne un exemple de chaque.
\item Écris l'équation-bilan de l'autoprotolyse de l'eau et donne l'expression du produit ionique $K_e$.
\item Donne la définition du pH d'une solution, puis la formule permettant de calculer $[\ce{H3O+}]$ à partir du pH.
\item Cite deux méthodes de mesure du pH en précisant, pour chacune, la précision obtenue.
\end{enumerate}
\end{exercice}

\courssub{Je m'entraîne}

\begin{exercice}[Calculs de pH et de concentrations]
On donne à \SI{25}{\celsius} : $K_e = 10^{-14}$.
\begin{enumerate}
\item Le jus de bissap préparé par une vendeuse du marché de Bafoussam a un pH de $3{,}2$. Calcule $[\ce{H3O+}]$ et $[\ce{HO-}]$ dans ce jus.
\item Une solution de soude a une concentration en ions hydroxyde $[\ce{HO-}] = \SI{5,0e-3}{mol.L^{-1}}$. Calcule $[\ce{H3O+}]$ puis le pH de la solution.
\item Une eau minérale a une concentration $[\ce{H3O+}] = \SI{4,0e-8}{mol.L^{-1}}$. Calcule son pH et précise son caractère acide, neutre ou basique.
\end{enumerate}
\end{exercice}

\begin{exercice}[Dilution d'une solution d'acide chlorhydrique]
On dispose d'une solution $S_0$ d'acide chlorhydrique de pH $= 2{,}0$ à \SI{25}{\celsius}. L'acide chlorhydrique est un acide fort, totalement dissocié dans l'eau : \ce{HCl + H2O -> H3O+ + Cl-}.
\begin{enumerate}
\item Calcule la concentration $C_0$ de la solution $S_0$.
\item On veut préparer $V_1 = \SI{500}{mL}$ d'une solution $S_1$ de pH $= 3{,}0$. Calcule la concentration $C_1$ de $S_1$.
\item Calcule le volume $V_0$ de solution $S_0$ à prélever. Décris précisément le protocole de cette dilution (verrerie utilisée, étapes, sécurité).
\end{enumerate}
\end{exercice}

\begin{exercice}[Indicateurs colorés]
Un technicien de laboratoire à Douala teste une solution incolore $S$ avec deux indicateurs colorés. La solution devient \textbf{jaune} en présence d'hélianthine et \textbf{bleue} en présence de bleu de bromothymol (BBT). Zones de virage :
\begin{center}
\begin{tabularx}{\linewidth}{|l|X|X|X|}
\hline
\rowcolor{popblueL} Indicateur & Couleur acide & Zone de virage & Couleur basique \\
\hline
Hélianthine & rouge & $3{,}1$ -- $4{,}4$ & jaune \\
\hline
BBT & jaune & $6{,}0$ -- $7{,}6$ & bleu \\
\hline
\end{tabularx}
\end{center}
\begin{enumerate}
\item À l'aide de la teinte observée avec l'hélianthine, donne une inégalité portant sur le pH de $S$.
\item À l'aide de la teinte observée avec le BBT, donne une seconde inégalité portant sur le pH de $S$.
\item Déduis-en un encadrement du pH de $S$. La solution est-elle acide ou basique ?
\item Calcule l'encadrement correspondant de $[\ce{H3O+}]$.
\end{enumerate}
\end{exercice}

\begin{exercice}[Inventaire des espèces et approximations]
On dissout du chlorure de sodium \ce{NaCl} dans l'eau distillée ; on obtient une solution de concentration $C = \SI{1,0e-2}{mol.L^{-1}}$ à \SI{25}{\celsius}.
\begin{enumerate}
\item Fais l'inventaire de toutes les espèces chimiques (molécules et ions) présentes en solution.
\item Écris l'équation d'électroneutralité de la solution.
\item Justifie que $[\ce{H3O+}] = [\ce{HO-}]$. En déduire leur valeur et le pH de la solution.
\item Classe les espèces en espèces majoritaires, minoritaires et ultra-minoritaires.
\end{enumerate}
\end{exercice}

\courssub{Je me prépare au Bac}

\begin{exercice}[Type Bac : l'eau de forage de Ngaoundéré]
\textit{(D'après Baccalauréat TI, Cameroun)}

Une coopérative agricole de Ngaoundéré fait analyser l'eau d'un forage destiné à l'abreuvage du bétail. À \SI{25}{\celsius}, le laboratoire mesure un pH de $6{,}2$. On donne : $K_e = 10^{-14}$ à \SI{25}{\celsius} ; $10^{-6{,}2} \approx 6{,}3 \times 10^{-7}$.
\begin{enumerate}
\item Rappelle la définition du pH et écris l'équation-bilan de l'autoprotolyse de l'eau.
\item Calcule les concentrations $[\ce{H3O+}]$ et $[\ce{HO-}]$ de cette eau de forage.
\item Cette eau est-elle acide, neutre ou basique ? Justifie par comparaison des deux concentrations.
\item L'eau du forage contient en outre des ions calcium \ce{Ca^2+} ($\SI{2,0e-4}{mol.L^{-1}}$), hydrogénocarbonate \ce{HCO3-} ($\SI{4,0e-4}{mol.L^{-1}}$), chlorure \ce{Cl-} ($\SI{1,0e-4}{mol.L^{-1}}$) et sodium \ce{Na+} ($\SI{1,0e-4}{mol.L^{-1}}$).
\begin{enumerate}
\item Fais l'inventaire complet des espèces ioniques présentes.
\item Identifie les espèces majoritaires et les espèces ultra-minoritaires. Justifie.
\item Vérifie approximativement l'électroneutralité de la solution.
\end{enumerate}
\item Le technicien hésite entre le papier pH et un pH-mètre pour contrôler quotidiennement le forage. Indique, en justifiant, l'appareil le plus adapté à un suivi précis.
\end{enumerate}
\end{exercice}

\begin{exercice}[Type Bac : préparation d'une solution de soude à la SONARA]
\textit{(D'après Baccalauréat TI, Cameroun)}

Dans un laboratoire de contrôle qualité, un technicien doit préparer $V = \SI{250}{mL}$ d'une solution $S$ d'hydroxyde de sodium (soude) de concentration $C = \SI{2,0e-2}{mol.L^{-1}}$. La soude est une base forte, totalement dissociée dans l'eau : \ce{NaOH -> Na+ + HO-}. On donne : $M(\ce{NaOH}) = \SI{40,0}{g.mol^{-1}}$ ; $K_e = 10^{-14}$ à \SI{25}{\celsius}.
\begin{enumerate}
\item Calcule la quantité de matière $n$ d'hydroxyde de sodium nécessaire, puis la masse $m$ à peser.
\item Décris le protocole de préparation de la solution $S$ (verrerie, étapes) et cite deux règles de sécurité liées à la manipulation de la soude.
\item Calcule $[\ce{HO-}]$ dans $S$, puis $[\ce{H3O+}]$ et le pH de la solution $S$.
\item On prélève un volume de $S$ que l'on dilue dix fois pour obtenir une solution $S'$.
\begin{enumerate}
\item Calcule la concentration $C'$ de $S'$.
\item Calcule le pH de $S'$. Compare au pH de $S$ et conclus sur l'effet de la dilution.
\end{enumerate}
\item On ajoute ensuite à \SI{20,0}{mL} de $S'$ quelques gouttes de bleu de bromothymol (zone de virage $6{,}0$ -- $7{,}6$, teinte acide jaune, teinte basique bleue). Quelle est la teinte observée ? Justifie.
\end{enumerate}
\end{exercice}

\begin{exercice}[Type Bac : analyse médicale et acidose]
\textit{(D'après Baccalauréat TI, Cameroun)}

À l'hôpital régional de Bafoussam, un laborantin mesure le pH sanguin d'un patient. À la température du corps humain, $\theta = \SI{37}{\celsius}$, le produit ionique de l'eau vaut $pK_e = 13{,}6$. Le résultat d'analyse indique : pH sanguin $= 7{,}30$. On rappelle que chez un individu sain, le pH du sang est compris entre $7{,}38$ et $7{,}42$ ; on parle d'\emph{acidose} lorsque le pH sanguin est inférieur à $7{,}38$.
\begin{enumerate}
\item Donne la valeur de $K_e$ à \SI{37}{\celsius}.
\item À \SI{37}{\celsius}, calcule le pH d'une solution neutre. En déduire les valeurs de $[\ce{H3O+}]$ et $[\ce{HO-}]$ dans une solution neutre à cette température.
\item Calcule $[\ce{H3O+}]$ et $[\ce{HO-}]$ dans le sang du patient. On donne : $10^{-7{,}30} \approx 5{,}0 \times 10^{-8}$.
\item Le sang du patient est-il acide, neutre ou basique \emph{au sens chimique} à \SI{37}{\celsius} ? Justifie soigneusement en tenant compte de la température.
\item Le patient est-il en acidose \emph{au sens médical} ? Justifie. Explique pourquoi les notions de neutralité chimique et de neutralité physiologique ne coïncident pas ici.
\item Le pH sanguin est régulé notamment par le couple \ce{CO2,H2O}/\ce{HCO3-} ($pK_a = 6{,}1$ à \SI{37}{\celsius}). Écris l'équation de la réaction de l'acide de ce couple avec l'eau et indique le rôle (acide ou base) joué par l'ion \ce{HCO3-} dans ce couple.
\end{enumerate}
\end{exercice}

\newpage

\coursec{Corrigés des exercices}

\begin{corrige}[Exercice 1 — QCM : à toi de choisir]
\begin{enumerate}
\item \textbf{Réponse b)} : à \SI{25}{\celsius}, $K_e = [\ce{H3O+}][\ce{HO-}] = 10^{-14}$. (La valeur $10^{-7}$ est la concentration en \ce{H3O+} dans l'eau pure, pas le produit ionique.)
\item \textbf{Réponse c)} : à \SI{25}{\celsius}, pH $= 9{,}0 > 7{,}0$, donc $[\ce{H3O+}] < [\ce{HO-}]$ : la solution est \textbf{basique}.
\item \textbf{Réponse b)} : $\mathrm{pH} = -\log(2{,}0\times10^{-4}) = 4 - \log 2{,}0 = 4 - 0{,}30 = 3{,}70$.
\item \textbf{Réponse b)} : capter un proton est la définition d'une \textbf{base de Brønsted}. Dans \ce{2 H2O <=> H3O+ + HO-}, l'eau qui devient \ce{H3O+} a capté \ce{H+}.
\item \textbf{Réponse b)} : diluer au dixième signifie $f = \dfrac{V_1}{V_0} = 10$, donc $V_1 = 10 \times V_0 = 10 \times \SI{10,0}{mL} = \SI{100}{mL}$. On utilise une fiole jaugée de \SI{100}{mL}.
\end{enumerate}
\textbf{Point de vigilance :} ne pas confondre $K_e = 10^{-14}$ (produit de deux concentrations) et $pK_e = 14$ (son cologarithme).
\end{corrige}

\begin{corrige}[Exercice 2 — Vrai ou faux ? Justifie]
\begin{enumerate}
\item \textbf{FAUX.} Le pH de neutralité dépend de la température : $\mathrm{pH}_{\text{neutre}} = \dfrac{pK_e}{2}$. À \SI{25}{\celsius} il vaut $7{,}0$, mais à \SI{37}{\celsius} il vaut $6{,}8$ et à \SI{100}{\celsius} environ $6{,}15$, car l'autoprotolyse est endothermique et $K_e$ augmente avec la température.
\item \textbf{VRAI.} L'équilibre d'autoprotolyse existe dans toute solution aqueuse ; la relation $[\ce{H3O+}][\ce{HO-}] = K_e$ est donc toujours vérifiée, quels que soient les solutés présents, à température donnée.
\item \textbf{VRAI.} L'eau peut céder un proton (couple \ce{H2O}/\ce{HO-}, rôle d'acide) ou en capter un (couple \ce{H3O+}/\ce{H2O}, rôle de base) : c'est un \cle{ampholyte}.
\item \textbf{FAUX.} Quand on dilue dix fois un acide fort, $[\ce{H3O+}]$ est divisée par 10, donc le pH \emph{augmente} d'une unité ($\mathrm{pH} = -\log[\ce{H3O+}]$) : la solution devient moins acide, pas plus.
\item \textbf{FAUX.} Le papier pH donne une précision de l'ordre de $\pm 0{,}5$ à $1$ unité de pH. Seul le pH-mètre atteint une précision de $\pm 0{,}01$ à $0{,}05$ unité.
\end{enumerate}
\end{corrige}

\begin{corrige}[Exercice 3 — Texte à trous]
L'eau pure conduit très faiblement le courant électrique car elle est le siège d'une réaction limitée appelée \textbf{autoprotolyse} de l'eau. Au cours de cette réaction, une molécule d'eau cède un \textbf{proton} à une autre molécule d'eau : il se forme alors l'ion \textbf{hydronium} \ce{H3O+} et l'ion \textbf{hydroxyde} \ce{HO-}. L'eau se comporte à la fois comme un acide et comme une base : on dit que c'est un \textbf{ampholyte}. À \SI{25}{\celsius}, le produit ionique de l'eau vaut $K_e = \mathbf{10^{-14}}$. Quand on ajoute un acide à de l'eau, la concentration $[\ce{H3O+}]$ \textbf{augmente} et le pH \textbf{diminue}. Toute solution aqueuse vérifie la loi d'\textbf{électroneutralité} : la somme des charges positives est égale à la somme des charges négatives.
\end{corrige}

\begin{corrige}[Exercice 4 — Définitions et formules]
\begin{enumerate}
\item Un \textbf{acide de Brønsted} est une espèce chimique capable de céder un proton \ce{H+} ; exemple : \ce{CH3COOH} (acide éthanoïque du vinaigre), qui cède \ce{H+} pour donner \ce{CH3COO-}. Une \textbf{base de Brønsted} est une espèce capable de capter un proton \ce{H+} ; exemple : \ce{HO-} (ion hydroxyde), qui capte \ce{H+} pour donner \ce{H2O}.
\item Équation-bilan de l'autoprotolyse de l'eau :
\[ \ce{2 H2O(l) <=> H3O+(aq) + HO-(aq)} \]
Produit ionique de l'eau : $K_e = [\ce{H3O+}] \times [\ce{HO-}]$ (sans unité), avec $K_e = 10^{-14}$ à \SI{25}{\celsius}.
\item Le pH est défini par $\mathrm{pH} = -\log_{10}[\ce{H3O+}]$, avec $[\ce{H3O+}]$ exprimée en \si{mol.L^{-1}}. Réciproquement : $[\ce{H3O+}] = 10^{-\mathrm{pH}}$ (en \si{mol.L^{-1}}).
\item Deux méthodes de mesure :
\begin{itemize}
\item Le \textbf{papier pH} : comparaison de la couleur prise par la bandelette à une échelle de teintes ; précision de $\pm 0{,}5$ à $1$ unité.
\item Le \textbf{pH-mètre} : mesure de la différence de potentiel entre une électrode de verre et une électrode de référence ; précision de $\pm 0{,}01$ à $0{,}05$ unité.
\end{itemize}
(On pouvait aussi citer les indicateurs colorés, précision $\pm 0{,}2$ à $0{,}5$ unité.)
\end{enumerate}
\end{corrige}

\begin{corrige}[Exercice 5 — Calculs de pH et de concentrations]
\begin{enumerate}
\item Jus de bissap, $\mathrm{pH} = 3{,}2$ :
\[ [\ce{H3O+}] = 10^{-\mathrm{pH}} = 10^{-3{,}2} = 10^{-4} \times 10^{0{,}8} \approx \SI{6,3e-4}{mol.L^{-1}}. \]
\[ [\ce{HO-}] = \frac{K_e}{[\ce{H3O+}]} = \frac{10^{-14}}{6{,}3\times10^{-4}} \approx \SI{1,6e-11}{mol.L^{-1}}. \]
Vérification : $6{,}3\times10^{-4} \times 1{,}6\times10^{-11} \approx 1{,}0\times10^{-14} = K_e$. Le jus est acide, ce qui est cohérent avec le goût acidulé du bissap.
\item Solution de soude, $[\ce{HO-}] = \SI{5,0e-3}{mol.L^{-1}}$ :
\[ [\ce{H3O+}] = \frac{K_e}{[\ce{HO-}]} = \frac{10^{-14}}{5{,}0\times10^{-3}} = \SI{2,0e-12}{mol.L^{-1}}. \]
\[ \mathrm{pH} = -\log(2{,}0\times10^{-12}) = 12 - \log 2{,}0 = 12 - 0{,}30 = \mathbf{11{,}70}. \]
Solution fortement basique, cohérent avec une solution de soude.
\item Eau minérale, $[\ce{H3O+}] = \SI{4,0e-8}{mol.L^{-1}}$ :
\[ \mathrm{pH} = -\log(4{,}0\times10^{-8}) = 8 - \log 4{,}0 = 8 - 0{,}60 = \mathbf{7{,}40}. \]
Comme $[\ce{H3O+}] = 4{,}0\times10^{-8} < \sqrt{K_e} = 10^{-7}$ \si{mol.L^{-1}}, on a $[\ce{H3O+}] < [\ce{HO-}]$ : l'eau est \textbf{légèrement basique} (pH $> 7$).
\end{enumerate}
\end{corrige}

\begin{corrige}[Exercice 6 — Dilution d'une solution d'acide chlorhydrique]
\begin{enumerate}
\item \ce{HCl} est un acide fort totalement dissocié : $[\ce{H3O+}]_0 = C_0$. Or $\mathrm{pH}_0 = 2{,}0$, donc :
\[ C_0 = [\ce{H3O+}]_0 = 10^{-2{,}0} = \SI{1,0e-2}{mol.L^{-1}}. \]
\item Pour $S_1$ de pH $= 3{,}0$ :
\[ C_1 = [\ce{H3O+}]_1 = 10^{-3{,}0} = \SI{1,0e-3}{mol.L^{-1}}. \]
Le facteur de dilution est $f = \dfrac{C_0}{C_1} = \dfrac{10^{-2}}{10^{-3}} = 10$ : dilution au dixième.
\item Conservation de la quantité de matière : $C_0 V_0 = C_1 V_1$, donc :
\[ V_0 = \frac{C_1 V_1}{C_0} = \frac{1{,}0\times10^{-3} \times 500}{1{,}0\times10^{-2}} = \SI{50,0}{mL}. \]
\textbf{Protocole :}
\begin{enumerate}
\item Verser un peu de solution $S_0$ dans un bécher propre ; rincer une \textbf{pipette jaugée de \SI{50,0}{mL}} avec $S_0$.
\item Prélever $V_0 = \SI{50,0}{mL}$ de $S_0$ à l'aide de la pipette munie d'une propipette (\textbf{jamais à la bouche}).
\item Introduire le prélèvement dans une \textbf{fiole jaugée de \SI{500}{mL}} contenant déjà un peu d'eau distillée.
\item Compléter avec de l'eau distillée jusqu'au trait de jauge, en ajustant le bas du ménisque au niveau du trait.
\item Boucher et homogénéiser en retournant la fiole plusieurs fois.
\end{enumerate}
\textbf{Sécurité :} porter blouse, gants et lunettes ; verser toujours l'acide dans l'eau.
\end{enumerate}
\end{corrige}

\begin{corrige}[Exercice 7 — Indicateurs colorés]
\begin{enumerate}
\item L'hélianthine est \textbf{jaune}, c'est-à-dire sous sa teinte basique : le pH est au-delà de sa zone de virage :
\[ \mathrm{pH} > 4{,}4. \]
\item Le BBT est \textbf{bleu}, c'est-à-dire sous sa teinte basique : le pH est au-delà de sa zone de virage :
\[ \mathrm{pH} > 7{,}6. \]
\item Les deux conditions doivent être satisfaites simultanément ; la plus restrictive est $\mathrm{pH} > 7{,}6$. On obtient donc l'encadrement :
\[ \boxed{7{,}6 < \mathrm{pH}} \quad \text{(en pratique } 7{,}6 < \mathrm{pH} < 14 \text{)}. \]
Comme pH $> 7$, la solution $S$ est \textbf{basique}.
\item Comme $[\ce{H3O+}] = 10^{-\mathrm{pH}}$ est une fonction décroissante du pH :
\[ \mathrm{pH} > 7{,}6 \iff [\ce{H3O+}] < 10^{-7{,}6} \approx \SI{2,5e-8}{mol.L^{-1}}. \]
Donc $[\ce{H3O+}] < \SI{2,5e-8}{mol.L^{-1}}$, ce qui confirme le caractère basique ($[\ce{H3O+}] < 10^{-7}$ \si{mol.L^{-1}}).
\end{enumerate}
\textbf{Point de vigilance :} quand un indicateur prend sa teinte basique, le pH est \emph{supérieur} à la borne haute de la zone de virage ; attention au sens des inégalités lors du passage aux concentrations (la fonction $10^{-x}$ inverse l'ordre).
\end{corrige}

\begin{corrige}[Exercice 8 — Inventaire des espèces et approximations]
\begin{enumerate}
\item \ce{NaCl} est totalement dissocié : \ce{NaCl -> Na+ + Cl-}. L'eau est le siège de son autoprotolyse. Espèces présentes :
\begin{itemize}
\item Molécule : \ce{H2O} ;
\item Ions : \ce{Na+}, \ce{Cl-}, \ce{H3O+}, \ce{HO-}.
\end{itemize}
\item Équation d'électroneutralité :
\[ [\ce{Na+}] + [\ce{H3O+}] = [\ce{Cl-}] + [\ce{HO-}]. \]
\item La dissociation du sel impose $[\ce{Na+}] = [\ce{Cl-}] = C = \SI{1,0e-2}{mol.L^{-1}}$. En reportant dans l'électroneutralité, il reste $[\ce{H3O+}] = [\ce{HO-}]$ : la solution est \textbf{neutre}. D'où :
\[ [\ce{H3O+}] = [\ce{HO-}] = \sqrt{K_e} = \SI{1,0e-7}{mol.L^{-1}}, \qquad \mathrm{pH} = -\log(10^{-7}) = \mathbf{7{,}00}. \]
\item Classification des espèces :
\begin{center}
\begin{tabularx}{\linewidth}{|l|X|}
\hline
\rowcolor{popblueL} \textbf{Catégorie} & \textbf{Espèces} \\
\hline
Majoritaires ($C > 10^{-3}$ \si{mol.L^{-1}}) & \ce{H2O} ($\approx \SI{55,5}{mol.L^{-1}}$), \ce{Na+} et \ce{Cl-} ($10^{-2}$ \si{mol.L^{-1}}) \\
\hline
Minoritaires ($10^{-6} < C < 10^{-3}$) & aucune \\
\hline
Ultra-minoritaires ($C < 10^{-6}$ \si{mol.L^{-1}}) & \ce{H3O+} et \ce{HO-} ($10^{-7}$ \si{mol.L^{-1}}) \\
\hline
\end{tabularx}
\end{center}
\end{enumerate}
\end{corrige}

\begin{corrige}[Exercice 9 — Type Bac : l'eau de forage de Ngaoundéré]
\begin{enumerate}
\item Le pH est défini par $\mathrm{pH} = -\log_{10}[\ce{H3O+}]$ ($[\ce{H3O+}]$ en \si{mol.L^{-1}}). Autoprotolyse de l'eau :
\[ \ce{2 H2O(l) <=> H3O+(aq) + HO-(aq)}. \]
\item Concentrations ioniques :
\[ [\ce{H3O+}] = 10^{-\mathrm{pH}} = 10^{-6{,}2} \approx \SI{6,3e-7}{mol.L^{-1}}. \]
\[ [\ce{HO-}] = \frac{K_e}{[\ce{H3O+}]} = \frac{10^{-14}}{6{,}3\times10^{-7}} \approx \SI{1,6e-8}{mol.L^{-1}}. \]
\item On compare : $[\ce{H3O+}] = 6{,}3\times10^{-7} > [\ce{HO-}] = 1{,}6\times10^{-8}$ \si{mol.L^{-1}}. De plus pH $= 6{,}2 < 7{,}0$ à \SI{25}{\celsius}. L'eau du forage est donc \textbf{légèrement acide}.
\item Ions minéraux de l'eau :
\begin{enumerate}
\item Inventaire complet des espèces ioniques : \ce{Ca^2+}, \ce{HCO3-}, \ce{Cl-}, \ce{Na+}, \ce{H3O+}, \ce{HO-} (et la molécule \ce{H2O}).
\item Espèces \textbf{majoritaires} ($C > 10^{-3}$ \si{mol.L^{-1}}) : \ce{H2O} ($\approx \SI{55,5}{mol.L^{-1}}$). Espèces \textbf{minoritaires} ($10^{-6} < C < 10^{-3}$ \si{mol.L^{-1}}) : \ce{Ca^2+}, \ce{HCO3-}, \ce{Cl-}, \ce{Na+} (concentrations de l'ordre de $10^{-4}$ \si{mol.L^{-1}}). Espèces \textbf{ultra-minoritaires} ($C < 10^{-6}$ \si{mol.L^{-1}}) : \ce{H3O+} ($6{,}3\times10^{-7}$) et \ce{HO-} ($1{,}6\times10^{-8}$ \si{mol.L^{-1}}), dont les concentrations sont environ 1000 fois plus faibles que celles des ions minéraux.
\item Électroneutralité : charges positives $=$ charges négatives.
\begin{align*}
\sum (+) &= 2[\ce{Ca^2+}] + [\ce{Na+}] + [\ce{H3O+}] = 2(2{,}0\times10^{-4}) + 1{,}0\times10^{-4} + 6{,}3\times10^{-7}\\
&= 4{,}0\times10^{-4} + 1{,}0\times10^{-4} + 0{,}0063\times10^{-4} \approx 5{,}0006\times10^{-4}\ \si{mol.L^{-1}}.\\
\sum (-) &= [\ce{HCO3-}] + [\ce{Cl-}] + [\ce{HO-}] = 4{,}0\times10^{-4} + 1{,}0\times10^{-4} + 1{,}6\times10^{-8}\\
&= 5{,}0\times10^{-4} + 0{,}00016\times10^{-4} \approx 5{,}0002\times10^{-4}\ \si{mol.L^{-1}}.
\end{align*}
Les deux sommes sont égales à mieux que $0{,}1\,\%$ près : l'électroneutralité est \textbf{vérifiée} (l'écart provient des arrondis des analyses). Attention : l'ion \ce{Ca^2+} compte pour \emph{deux} charges positives.
\end{enumerate}
\item Pour un suivi quotidien précis (détection de variations de $0{,}1$ unité de pH), le \textbf{pH-mètre} est le plus adapté : précision de $\pm 0{,}01$ à $0{,}05$ unité, contre $\pm 0{,}5$ à $1$ pour le papier pH. Le papier pH reste utile pour un contrôle rapide de dépannage sur le terrain.
\end{enumerate}
\textbf{Barème indicatif (sur 10) :} 1. définition + équation : 2 pts ; 2. calculs des concentrations : 2 pts ; 3. caractère acide justifié : 1 pt ; 4a. inventaire : 1 pt ; 4b. classification : 1 pt ; 4c. électroneutralité : 2 pts ; 5. choix justifié : 1 pt.\\
\textbf{Points de vigilance :} ne pas oublier le facteur 2 devant $[\ce{Ca^2+}]$ dans l'électroneutralité ; utiliser la valeur approchée $10^{-6{,}2} \approx 6{,}3\times10^{-7}$ fournie par l'énoncé ; toujours conclure sur le caractère acido-basique en comparant les \emph{deux} concentrations.
\end{corrige}

\begin{corrige}[Exercice 10 — Type Bac : préparation d'une solution de soude à la SONARA]
\begin{enumerate}
\item Quantité de matière et masse :
\[ n = C \times V = 2{,}0\times10^{-2} \times 0{,}250 = \SI{5,0e-3}{mol}. \]
\[ m = n \times M = 5{,}0\times10^{-3} \times 40{,}0 = \SI{0,20}{g}. \]
\item \textbf{Protocole :}
\begin{enumerate}
\item Peser précisément $m = \SI{0,20}{g}$ de pastilles de soude sur une balance au centième, dans une coupelle ou un verre de montre.
\item Introduire le solide dans une \textbf{fiole jaugée de \SI{250}{mL}} à l'aide d'un entonnoir ; rincer la coupelle et l'entonnoir à l'eau distillée.
\item Remplir la fiole aux deux tiers d'eau distillée et agiter jusqu'à dissolution complète (la dissolution est exothermique : laisser tiédir).
\item Compléter à l'eau distillée jusqu'au trait de jauge, boucher et homogénéiser.
\end{enumerate}
\textbf{Règles de sécurité (deux au choix) :} la soude est corrosive : porter gants, lunettes et blouse ; ne jamais verser d'eau sur la soude concentrée mais dissoudre progressivement ; en cas de contact cutané, rincer abondamment à l'eau pendant au moins 15 min.
\item La soude est une base forte totalement dissociée : $[\ce{HO-}] = C = \SI{2,0e-2}{mol.L^{-1}}$.
\[ [\ce{H3O+}] = \frac{K_e}{[\ce{HO-}]} = \frac{10^{-14}}{2{,}0\times10^{-2}} = \SI{5,0e-13}{mol.L^{-1}}. \]
\[ \mathrm{pH} = -\log(5{,}0\times10^{-13}) = 13 - \log 5{,}0 = 13 - 0{,}70 = \mathbf{12{,}30}. \]
\item Dilution au dixième :
\begin{enumerate}
\item $C' = \dfrac{C}{10} = \SI{2,0e-3}{mol.L^{-1}}$.
\item $[\ce{HO-}]' = \SI{2,0e-3}{mol.L^{-1}}$, donc $[\ce{H3O+}]' = \dfrac{10^{-14}}{2{,}0\times10^{-3}} = \SI{5,0e-12}{mol.L^{-1}}$ et :
\[ \mathrm{pH}' = -\log(5{,}0\times10^{-12}) = 12 - \log 5{,}0 = \mathbf{11{,}30}. \]
\textbf{Conclusion :} en diluant dix fois une base forte, le pH \emph{diminue} d'une unité ($12{,}30 \to 11{,}30$) : la solution devient moins basique, mais reste fortement basique.
\end{enumerate}
\item Dans $S'$, pH $= 11{,}30 > 7{,}6$ : le pH est au-dessus de la zone de virage du BBT, qui prend donc sa teinte basique : la solution devient \textbf{bleue}.
\end{enumerate}
\textbf{Barème indicatif (sur 10) :} 1. $n$ et $m$ : 2 pts ; 2. protocole + sécurité : 2 pts ; 3. $[\ce{HO-}]$, $[\ce{H3O+}]$, pH : 2 pts ; 4a. $C'$ : 1 pt ; 4b. pH$'$ + conclusion : 2 pts ; 5. teinte justifiée : 1 pt.\\
\textbf{Points de vigilance :} convertir \SI{250}{mL} en \SI{0,250}{L} ; pour une \emph{base}, calculer d'abord $[\ce{HO-}]$ puis passer par $K_e$ (ne jamais écrire pH $= -\log C$ pour une base !) ; la dilution d'une base forte fait \emph{baisser} le pH, contrairement à celle d'un acide.
\end{corrige}

\begin{corrige}[Exercice 11 — Type Bac : analyse médicale et acidose]
\begin{enumerate}
\item $pK_e = -\log K_e = 13{,}6$, donc :
\[ K_e = 10^{-13{,}6} = 10^{0{,}4}\times10^{-14} \approx \SI{2,5e-14}{} \quad \text{à } \SI{37}{\celsius}. \]
\item Pour une solution neutre, $[\ce{H3O+}] = [\ce{HO-}]$, donc $K_e = [\ce{H3O+}]^2$ et :
\[ \mathrm{pH}_{\text{neutre}} = \frac{pK_e}{2} = \frac{13{,}6}{2} = \mathbf{6{,}80}. \]
\[ [\ce{H3O+}] = [\ce{HO-}] = \sqrt{K_e} = 10^{-6{,}8} \approx \SI{1,6e-7}{mol.L^{-1}}. \]
\item Dans le sang du patient, pH $= 7{,}30$ :
\[ [\ce{H3O+}] = 10^{-7{,}30} \approx \SI{5,0e-8}{mol.L^{-1}}. \]
\[ [\ce{HO-}] = \frac{K_e}{[\ce{H3O+}]} = \frac{2{,}5\times10^{-14}}{5{,}0\times10^{-8}} = \SI{5,0e-7}{mol.L^{-1}}. \]
\item \textbf{Au sens chimique}, à \SI{37}{\celsius} la neutralité correspond à pH $= 6{,}80$. Ici :
\[ [\ce{H3O+}] = 5{,}0\times10^{-8} < [\ce{HO-}] = 5{,}0\times10^{-7}\ \si{mol.L^{-1}}, \]
et pH $= 7{,}30 > 6{,}80$ : le sang du patient est \textbf{basique au sens chimique} à cette température. C'est le cas de tout sang humain normal.
\item \textbf{Au sens médical}, le pH du patient ($7{,}30$) est inférieur à la borne normale $7{,}38$ : le patient est bien en \textbf{acidose}. Les deux notions ne coïncident pas car la neutralité \emph{chimique} est définie par $[\ce{H3O+}] = [\ce{HO-}]$ (soit pH $= 6{,}80$ à \SI{37}{\celsius}), tandis que la neutralité \emph{physiologique} est définie par rapport à la plage de fonctionnement normal de l'organisme ($7{,}38$ à $7{,}42$), qui est déjà légèrement basique au sens chimique. Une chute de pH, même si le sang reste chimiquement basique, est dangereuse pour l'organisme.
\item Couple \ce{CO2,H2O}/\ce{HCO3-} : le dioxyde de carbone hydraté joue le rôle d'acide :
\[ \ce{CO2,H2O + H2O <=> HCO3- + H3O+}. \]
Dans ce couple, l'ion \ce{HCO3-} est la \textbf{base conjuguée} (il peut capter un proton pour régénérer \ce{CO2,H2O}). C'est ainsi qu'il neutralise un excès d'ions \ce{H3O+} dans le sang : c'est le principal \emph{tampon} sanguin.
\end{enumerate}
\textbf{Barème indicatif (sur 12) :} 1. $K_e$ : 1 pt ; 2. pH neutre + concentrations : 2 pts ; 3. concentrations sanguines : 2 pts ; 4. caractère chimique justifié : 2 pts ; 5. acidose + explication : 3 pts ; 6. équation + rôle de \ce{HCO3-} : 2 pts.\\
\textbf{Points de vigilance :} le piège majeur de cet exercice est d'affirmer « pH $= 7{,}30 > 7$ donc basique » sans tenir compte de la température : à \SI{37}{\celsius}, la référence de neutralité est $6{,}80$, pas $7{,}00$. Toujours comparer pH et $\dfrac{pK_e}{2}$ \emph{à la température de l'expérience}. Autre subtilité : un sang « basique au sens chimique » peut correspondre à une « acidose au sens médical » : les deux vocabulaires ne se recouvrent pas.
\end{corrige}

\newpage

\coursec{Fiche bilan}

\begin{popfigure}
\begin{center}
\begin{center}\tcbox[colback=popdark,colframe=popdark,coltext=white,arc=9pt,boxsep=4pt,left=6pt,right=6pt]{\bfseries\small Eau, autoprotolyse et pH}\end{center}
\vspace{-2pt}
\begin{tcbraster}[raster columns=3,raster equal height=rows,raster column skip=6pt,raster row skip=6pt,enhanced,blankest]
\begin{tcolorbox}[enhanced,colback=white,colframe=popblue,boxrule=0.9pt,arc=3pt,title={L'eau solvant},fonttitle=\bfseries\scriptsize,coltitle=white,colbacktitle=popblue,left=4pt,right=4pt,top=2pt,bottom=2pt,boxsep=1pt,toptitle=1.5pt,bottomtitle=1.5pt,halign=left]
\begin{itemize}[nosep,leftmargin=*,itemsep=1pt,font=\scriptsize]
\item Molécule polaire
\item Liaisons H
\item Dissocie les ions
\end{itemize}
\end{tcolorbox}
\begin{tcolorbox}[enhanced,colback=white,colframe=popgreen,boxrule=0.9pt,arc=3pt,title={Bronsted},fonttitle=\bfseries\scriptsize,coltitle=white,colbacktitle=popgreen,left=4pt,right=4pt,top=2pt,bottom=2pt,boxsep=1pt,toptitle=1.5pt,bottomtitle=1.5pt,halign=left]
\begin{itemize}[nosep,leftmargin=*,itemsep=1pt,font=\scriptsize]
\item Acide : donneur \ce{H+}
\item Base : accepteur \ce{H+}
\item Couples conjugués
\end{itemize}
\end{tcolorbox}
\begin{tcolorbox}[enhanced,colback=white,colframe=poporange,boxrule=0.9pt,arc=3pt,title={Autoprotolyse},fonttitle=\bfseries\scriptsize,coltitle=white,colbacktitle=poporange,left=4pt,right=4pt,top=2pt,bottom=2pt,boxsep=1pt,toptitle=1.5pt,bottomtitle=1.5pt,halign=left]
\begin{itemize}[nosep,leftmargin=*,itemsep=1pt,font=\scriptsize]
\item \ce{2H2O <=> H3O+ + HO-}
\item $K_e = 10^{-14}$
\item p$K_e = 14$ (\SI{25}{\celsius})
\end{itemize}
\end{tcolorbox}
\begin{tcolorbox}[enhanced,colback=white,colframe=popred,boxrule=0.9pt,arc=3pt,title={Le pH},fonttitle=\bfseries\scriptsize,coltitle=white,colbacktitle=popred,left=4pt,right=4pt,top=2pt,bottom=2pt,boxsep=1pt,toptitle=1.5pt,bottomtitle=1.5pt,halign=left]
\begin{itemize}[nosep,leftmargin=*,itemsep=1pt,font=\scriptsize]
\item pH $= -\log[\ce{H3O+}]$
\item Acide $<$ 7, Neutre $=$ 7
\item Basique $>$ 7
\end{itemize}
\end{tcolorbox}
\begin{tcolorbox}[enhanced,colback=white,colframe=poppurple,boxrule=0.9pt,arc=3pt,title={Calculs},fonttitle=\bfseries\scriptsize,coltitle=white,colbacktitle=poppurple,left=4pt,right=4pt,top=2pt,bottom=2pt,boxsep=1pt,toptitle=1.5pt,bottomtitle=1.5pt,halign=left]
\begin{itemize}[nosep,leftmargin=*,itemsep=1pt,font=\scriptsize]
\item $[\ce{H3O+}] = 10^{-\mathrm{pH}}$
\item $[\ce{HO-}] = K_e/[\ce{H3O+}]$
\item Approximations
\end{itemize}
\end{tcolorbox}
\begin{tcolorbox}[enhanced,colback=white,colframe=popgold,boxrule=0.9pt,arc=3pt,title={Dilution},fonttitle=\bfseries\scriptsize,coltitle=white,colbacktitle=popgold,left=4pt,right=4pt,top=2pt,bottom=2pt,boxsep=1pt,toptitle=1.5pt,bottomtitle=1.5pt,halign=left]
\begin{itemize}[nosep,leftmargin=*,itemsep=1pt,font=\scriptsize]
\item $C_i V_i = C_f V_f$
\item Pipette + fiole jaugée
\item Facteur $F = C_i/C_f$
\end{itemize}
\end{tcolorbox}
\end{tcbraster}

\end{center}
\legende{Carte mentale de la leçon : les six idées forces à retenir pour le Bac TI.}
\end{popfigure}

\coursec{L'eau : structure et pouvoir solvant}

\begin{definition}[Molécule d'eau \ce{H2O}]
La molécule d'eau a une géométrie \textbf{coudée} (angle \ce{H-O-H} $\approx \SI{104,5}{\degree}$). L'oxygène, plus électronégatif que l'hydrogène, attire le nuage électronique des liaisons \ce{O-H} : la molécule est \cle{polaire} (dipôle permanent $\mu = \SI{1,85}{D}$). Les deux doublets non liants sur l'oxygène renforcent la polarité.
\end{definition}

\begin{propriete}[Liaisons hydrogène et pouvoir solvant]
En phase liquide, les molécules d'eau s'organisent en un réseau tridimensionnel maintenu par des \cle{liaisons hydrogène} (force $\approx \SI{20}{kJ.mol^{-1}}$). Ce réseau explique :
\begin{itemize}
  \item Les points d'ébullition et de fusion anormalement élevés.
  \item L'excellent pouvoir solvant pour les \textbf{espèces ioniques} (solvatation des cations par les doublets de l'O, des anions par les H $\delta^+$) et les \textbf{molécules polaires} (misibilité avec l'éthanol, le vinaigre, le bili-bili).
  \item La faible solubilité des corps gras apolaires (huile de palme, hydrocarbures SONARA).
\end{itemize}
\end{propriete}

\begin{savaistu}[L'eau au Cameroun]
L'eau de source (ex. \emph{Volcanic}, \emph{Tangui}) contient des ions minéraux (\ce{Ca^2+}, \ce{Mg^2+}, \ce{HCO3-}) dissous par érosion des roches volcaniques (Mont Cameroun). L'eau de mer (Kribi, Limbé) contient $\approx \SI{0,6}{mol.L^{-1}}$ de \ce{NaCl}. Le traitement de l'eau (CAMWATER, SNEC) utilise la coagulation-floculation (sulfate d'alumine) puis la chloration (\ce{Cl2}/\ce{HOCl}) pour la potabilisation.
\end{savaistu}

\coursec{Acides et bases selon Brønsted}

\begin{definition}[Acide et base de Brønsted (1923)]
\begin{itemize}
  \item Un \cle{acide} (noté \ce{AH}) est une espèce chimique capable de \textbf{céder} un proton \ce{H+} (ou ion hydrogène).
  \item Une \cle{base} (notée \ce{B}) est une espèce chimique capable de \textbf{capter} un proton \ce{H+}.
\end{itemize}
La réaction acide-base est un \textbf{transfert de proton} : $\ce{AH + B <=> A- + BH+}$.
\end{definition}

\begin{definition}[Couple acide/base conjugués]
Deux espèces qui ne diffèrent que par un proton forment un \cle{couple acide/base}, noté \ce{AH}/\ce{A-} (ou \ce{BH+}/\ce{B}). L'acide du couple a \textbf{un proton de plus} que sa base conjuguée.
\begin{center}
\begin{tabularx}{\linewidth}{|c|X|X|}\hline
\rowcolor{popblueL} \textbf{Couple} & \textbf{Acide} & \textbf{Base} \\\hline
\ce{H3O+}/\ce{H2O} & \ce{H3O+} (ion oxonium) & \ce{H2O} \\\hline
\ce{H2O}/\ce{HO-} & \ce{H2O} & \ce{HO-} (ion hydroxyde) \\\hline
\ce{HCl}/\ce{Cl-} & \ce{HCl} (acide chlorhydrique) & \ce{Cl-} (ion chlorure) \\\hline
\ce{CH3COOH}/\ce{CH3COO-} & \ce{CH3COOH} (acide éthanoïque) & \ce{CH3COO-} (ion éthanoate) \\\hline
\ce{NH4+}/\ce{NH3} & \ce{NH4+} (ion ammonium) & \ce{NH3} (ammoniac) \\\hline
\end{tabularx}
\end{center}
\end{definition}

\begin{propriete}[Amphotérie de l'eau]
L'eau peut agir comme un acide (donne \ce{H+} à une base plus forte qu'elle, ex. \ce{NH3}) ou comme une base (capte \ce{H+} d'un acide plus fort qu'elle, ex. \ce{HCl}). Elle est \cle{amphotère} (ou ampholyte).
\[ \ce{H2O + NH3 <=> HO- + NH4+} \quad \text{(eau = acide)} \]
\[ \ce{H2O + HCl -> H3O+ + Cl-} \quad \text{(eau = base)} \]
\end{propriete}

\begin{attention}[Pièges à éviter]
\begin{itemize}
  \item Ne pas confondre \ce{H+} (proton nu, n'existe pas seul en solution aqueuse) et \ce{H3O+} (ion oxonium, forme hydratée réelle). On écrit \ce{H3O+} dans les équations-bilan.
  \item Un acide fort a une base conjuguée \textbf{très faible} (négligeable comme base). Une base forte a un acide conjugué \textbf{très faible}.
  \item La force d'un acide/base dépend du solvant (niveau d'acidité/basicité de l'eau). En eau, \ce{HCl} est fort ; dans l'acide acétique glacial, il est faible.
\end{itemize}
\end{attention}

\coursec{Autoprotolyse de l'eau et produit ionique $K_e$}

\begin{experience}[Mise en évidence de l'autoprotolyse : conductimétrie]
\begin{itemize}
  \item \textbf{Matériel} : Conductimètre, eau distillée, eau du robinet, solution de \ce{NaCl} \SI{0,01}{mol.L^{-1}}, bécher.
  \item \textbf{Protocole} : Mesurer la conductivité $\sigma$ de l'eau distillée, de l'eau du robinet, de la solution saline.
  \item \textbf{Observations} : Eau distillée : $\sigma \approx \SI{5}{\micro S.cm^{-1}}$ (faible mais non nulle). Eau du robinet : $\sigma \approx \SI{200}{\micro S.cm^{-1}}$. Solution saline : $\sigma \approx \SI{1400}{\micro S.cm^{-1}}$.
  \item \textbf{Interprétation} : L'eau pure conduit faiblement le courant $\Rightarrow$ présence de quelques ions \ce{H3O+} et \ce{HO-} issus de la réaction de l'eau sur elle-même : \textbf{autoprotolyse}.
\end{itemize}
\end{experience}

\begin{propriete}[Équation-bilan et produit ionique de l'eau]
L'autoprotolyse est une réaction \textbf{limitée} (faible rendement) :
\[ \ce{2H2O(l) <=> H3O+(aq) + HO-(aq)} \]
La constante d'équilibre (produit ionique de l'eau) s'écrit :
\[ K_e = \frac{[\ce{H3O+}][\ce{HO-}]}{a_{\ce{H2O}}^2} \]
L'activité de l'eau $a_{\ce{H2O}} \approx 1$ (solvant en grande quantité), d'où :
\[ \boxed{K_e = [\ce{H3O+}][\ce{HO-}]} \]
$K_e$ est \textbf{constant à température donnée} pour \textbf{toute solution aqueuse} (pure, acide, basique, salée).
\end{propriete}

\begin{aretenir}[Valeurs de référence à \SI{25}{\celsius}]
\begin{itemize}
  \item $K_e = 1{,}0 \times 10^{-14}$ (sans unité, car activités sans unité).
  \item $\mathrm{p}K_e = -\log K_e = 14{,}00$.
  \item Dans l'eau pure : $[\ce{H3O+}] = [\ce{HO-}] = \sqrt{K_e} = 10^{-7}\ \si{mol.L^{-1}}$.
  \item \textbf{Variation avec la température} : $K_e$ augmente avec $T$ (réaction endothermique). À \SI{37}{\celsius} (température corps humain), $K_e \approx 2,5 \times 10^{-14}$, pH neutre $\approx 6,8$.
\end{itemize}
\end{aretenir}

\coursec{Le pH : définition, mesure et caractère de la solution}

\begin{definition}[pH d'une solution aqueuse]
Le \cle{pH} (potentiel hydrogène) est défini par :
\[ \boxed{\mathrm{pH} = -\log_{10}\left( \frac{[\ce{H3O+}]}{\SI{1}{mol.L^{-1}}} \right)} \]
Le pH est une \textbf{grandeur sans unité}. La concentration $[\ce{H3O+}]$ s'exprime en \si{mol.L^{-1}}.
\end{definition}

\begin{propriete}[Échelle de pH à \SI{25}{\celsius}]
\begin{center}
\begin{tabularx}{\linewidth}{|c|c|c|c|}\hline
\rowcolor{popblueL} \textbf{Caractère} & \textbf{Condition sur $[\ce{H3O+}]$} & \textbf{Condition sur $[\ce{HO-}]$} & \textbf{pH} \\\hline
Acide & $> 10^{-7}$ & $< 10^{-7}$ & $< 7$ \\\hline
Neutre & $= 10^{-7}$ & $= 10^{-7}$ & $= 7$ \\\hline
Basique & $< 10^{-7}$ & $> 10^{-7}$ & $> 7$ \\\hline
\end{tabularx}
\end{center}
\emph{Attention : ces seuils (7) ne valent qu'à \SI{25}{\celsius}.}
\end{propriete}

\begin{methode}[Mesure du pH]
\begin{enumerate}
  \item \textbf{pH-mètre (électrode de verre + référence)} : Précision $\pm \SI{0,01}{pH}$. \textbf{Étalonnage obligatoire} avant usage avec deux solutions tampons (ex. pH 4,01 et 7,00 ou 7,00 et 9,21). Rincer l'électrode à l'eau distillée, essuyer délicatement (pas frotter), plonger dans la solution, attendre stabilisation.
  \item \textbf{Papier pH / Bandelettes} : Rapide, précision $\pm \SI{0,5}{pH}$ à $\pm \SI{1}{pH}$. Tremper, comparer la couleur à l'échelle fournie.
  \item \textbf{Indicateurs colorés} (phénolphtaléine, bleu de bromothymol, héliantine) : Changement de couleur sur un intervalle $\mathrm{pH} = \mathrm{p}K_a \pm 1$. Utiles pour les dosages (équivalence).
\end{enumerate}
\end{methode}

\begin{savaistu}[pH des produits usuels au Cameroun]
\begin{itemize}
  \item \textbf{Bili-bili} (bière de mil/maïs) : pH $\approx 3,5-4,0$ (acide lactique, acide acétique).
  \item \textbf{Vin de palme} frais : pH $\approx 4,5-5,5$ ; après fermentation : pH $\approx 3,5$.
  \item \textbf{Jus de bissap} (\emph{Hibiscus sabdariffa}) : pH $\approx 2,5-3,0$ (acides organiques, anthocyanes = indicateur naturel rouge/rose en acide, bleu en base).
  \item \textbf{Savon artisanal} (huile de palme + soude) : pH $\approx 9-10$ (excès de soude, hydrolysis).
  \item \textbf{Sang humain} : pH $= 7,40 \pm 0,05$ (tampon bicarbonate/\ce{CO2}, vital).
  \item \textbf{Acide de batterie} (plomb) : $\ce{H2SO4} \approx \SI{4,5}{mol.L^{-1}}$, pH $< 0$.
\end{itemize}
\end{savaistu}

\coursec{Calculs de pH et concentrations ioniques}

\begin{methode}[Passer du pH aux concentrations (et inversement)]
\begin{enumerate}
  \item Identifier la donnée : pH ou $[\ce{H3O+}]$.
  \item \textbf{Si pH connu} : $[\ce{H3O+}] = 10^{-\mathrm{pH}}\ \si{mol.L^{-1}}$ (touche \texttt{10\textasciicircum x} ou \texttt{SHIFT log} sur calculatrice).
  \item Calculer $[\ce{HO-}] = \dfrac{K_e}{[\ce{H3O+}]} = 10^{\mathrm{pH} - \mathrm{p}K_e}\ \si{mol.L^{-1}}$.
  \item \textbf{Si $[\ce{H3O+}]$ connu} : $\mathrm{pH} = -\log[\ce{H3O+}]$ (touche \texttt{log}).
  \item Vérifier la cohérence : solution acide $\Rightarrow$ pH $< 7$ et $[\ce{H3O+}] > 10^{-7}$.
\end{enumerate}
\end{methode}

\begin{exoresolu}[Calcul des concentrations ioniques]
\textbf{Énoncé :} Une solution aqueuse a un pH $= 4,20$ à \SI{25}{\celsius}. Déterminer $[\ce{H3O+}]$ et $[\ce{HO-}]$. Conclure sur son caractère.
\textbf{Solution détaillée :}
\begin{align*}
[\ce{H3O+}] &= 10^{-4,20} = \SI{6,31 \times 10^{-5}}{mol.L^{-1}} \\
[\ce{HO-}] &= \frac{K_e}{[\ce{H3O+}]} = \frac{1,0 \times 10^{-14}}{6,31 \times 10^{-5}} = \SI{1,58 \times 10^{-10}}{mol.L^{-1}} \\
\text{Caractère :} &\quad [\ce{H3O+}] > [\ce{HO-}] \quad (\text{ou pH } 4,20 < 7) \Rightarrow \textbf{Solution acide}.
\end{align*}
\end{exoresolu}

\begin{exoresolu}[pH de solutions d'acide/base forts]
\textbf{Énoncé :} Calculer le pH à \SI{25}{\celsius} de :
\begin{enumerate}
  \item \SI{10,0}{mL} de solution mère d'acide chlorhydrique \ce{HCl} \SI{1,00}{mol.L^{-1}} diluée à \SI{100,0}{mL}.
  \item \SI{20,0}{mL} de solution de soude \ce{NaOH} \SI{0,100}{mol.L^{-1}}.
\end{enumerate}
\textbf{Solution détaillée :}
\begin{enumerate}
  \item \textbf{Dilution HCl} : $C_f = C_i \frac{V_i}{V_f} = 1,00 \times \frac{10,0}{100,0} = \SI{0,100}{mol.L^{-1}}$.
  \ce{HCl} acide fort $\Rightarrow$ réaction totale : $\ce{HCl + H2O -> H3O+ + Cl-}$.
  $[\ce{H3O+}] = C_f = \SI{0,100}{mol.L^{-1}}$.
  $\mathrm{pH} = -\log(0,100) = 1,00$.
  \item \textbf{NaOH} base forte $\Rightarrow$ dissociation totale : $\ce{NaOH -> Na+ + HO-}$.
  $[\ce{HO-}] = \SI{0,100}{mol.L^{-1}}$.
  $[\ce{H3O+}] = \frac{K_e}{[\ce{HO-}]} = \frac{10^{-14}}{0,100} = 10^{-13}\ \si{mol.L^{-1}}$.
  $\mathrm{pH} = -\log(10^{-13}) = 13,00$. \quad (Ou $\mathrm{pH} = \mathrm{p}K_e + \log[\ce{HO-}] = 14 + \log(0,1) = 13$).
\end{enumerate}
\end{exoresolu}

\begin{propriete}[Approximations : espèces majoritaires, minoritaires, ultra-minoritaires]
Dans une solution aqueuse, on classe les espèces selon leur concentration molaire $c$ :
\begin{center}
\begin{tabular*}{\linewidth}{@{\extracolsep{\fill}}|>{\centering\arraybackslash}p{0.12\linewidth}|>{\centering\arraybackslash}p{0.25\linewidth}|p{0.45\linewidth}|@{}}
\hline
\rowcolor{poporangeL} \textbf{Statut} & \textbf{Domaine de concentration} & \textbf{Conséquence pour les calculs} \\\hline
\cle{Majoritaire} & $c > 10^{-2}$ \si{mol.L^{-1}} & Détermine le pH (acide/base fort) ou l'équilibre (faible). \\\hline
\cle{Minoritaire} & $10^{-6} < c < 10^{-2}$ \si{mol.L^{-1}} & Présente, participe aux équilibres, non négligeable. \\\hline
\cle{Ultra-minoritaire} & $c < 10^{-6}$ \si{mol.L^{-1}} & Souvent négligeable devant les autres (ex. \ce{HO-} en milieu acide pH $< 6$). \\\hline
\end{tabular*}
\end{center}
\textbf{Règle pratique} : En solution acide (pH $< 6$), $[\ce{HO-}] < 10^{-8}$ \si{mol.L^{-1}} (ultra-minoritaire) $\Rightarrow$ on néglige $[\ce{HO-}]$ dans les bilans de matière/charge. En solution basique (pH $> 8$), $[\ce{H3O+}]$ est ultra-minoritaire.
\end{propriete}

\begin{attention}[Erreurs fréquentes en calcul de pH]
\begin{itemize}
  \item Oublier la dilution avant de calculer le pH (utiliser $C_i$ au lieu de $C_f$).
  \item Confondre $\log$ (base 10) et $\ln$ (base $e$) sur la calculatrice.
  \item Ne pas mettre de parenthèses : $-\log(10^{-3}) \neq -\log 10^{-3}$ (interprété comme $(-\log 10)^{-3}$ par certaines machines).
  \item Donner une unité au pH (pH est sans unité).
  \item Appliquer $\mathrm{pH} = 7$ pour la neutralité à toute température.
  \item Écrire $[\ce{H+}]$ au lieu de $[\ce{H3O+}]$ dans les formules de $K_e$ et pH.
\end{itemize}
\end{attention}

\coursec{Préparation de solutions par dilution}

\begin{definition}[Dilution]
La dilution consiste à augmenter le volume d'une solution (solution fille) en ajoutant du solvant (eau) à une solution plus concentrée (solution mère). La \textbf{quantité de matière de soluté est conservée} :
\[ n_{\text{soluté}} = C_i V_i = C_f V_f \]
$C_i, C_f$ en \si{mol.L^{-1}} ; $V_i, V_f$ en \si{L} (ou même unité de volume).
Le \cle{facteur de dilution} $F = \dfrac{C_i}{C_f} = \dfrac{V_f}{V_i} > 1$.
\end{definition}

\begin{methode}[Protocole rigoureux de dilution (verrerie de classe A)]
\begin{enumerate}
  \item \textbf{Calculer} $V_i = \dfrac{C_f V_f}{C_i}$ à prélever.
  \item \textbf{Choisir} une \textbf{pipette jaugée} de volume $V_i$ (ou pipette graduée si $V_i$ non standard, moins précis).
  \item \textbf{Rincer} la pipette : 3 fois avec l'eau distillée, puis 2 fois avec la solution mère (évite la dilution par l'eau résiduelle).
  \item \textbf{Prélever} $V_i$ : Aspirer au-dessus du trait de jauge (poire à pipeter, \textbf{jamais à la bouche}), abaisser le ménisque \textbf{au trait de jauge} (œil au niveau du ménisque, bas du ménisque sur le trait), laisser s'écouler par gravité dans une \textbf{fiole jaugée} de volume $V_f$ (ne pas souffler la dernière goutte).
  \item \textbf{Compléter} : Ajouter l'eau distillée au goulot, puis au compte-gouttes pour amener le bas du ménisque \textbf{exactement sur le trait de jauge} de la fiole.
  \item \textbf{Homogénéiser} : Boucher, retourner la fiole plusieurs fois (mouvements de rotation + inversion).
  \item \textbf{Étiqueter} : Nom, concentration, date, nom du préparateur, danger (pictogrammes CLP).
\end{enumerate}
\end{methode}

\begin{experience}[TP : Préparation de \SI{100,0}{mL} de \ce{HCl} \SI{0,100}{mol.L^{-1}} à partir de \ce{HCl} \SI{1,00}{mol.L^{-1}}]
\begin{itemize}
  \item \textbf{Calcul} : $V_i = \frac{0,100 \times 100,0}{1,00} = \SI{10,0}{mL}$.
  \item \textbf{Matériel} : Pipette jaugée \SI{10,0}{mL} (classe A), fiole jaugée \SI{100,0}{mL} (classe A), poire à pipeter, bécher \SI{100}{mL} pour la solution mère, eau distillée, compte-gouttes.
  \item \textbf{Contrôle} : Mesurer le pH de la solution fille. Théorique : pH $= 1,00$. Mesure attendue : $1,00 \pm 0,02$ (incertitude verrerie + température).
  \item \textbf{Sécurité} : \ce{HCl} \SI{1}{mol.L^{-1}} est corrosif (H314). Porter lunettes, blouse, gants. Sous hotte si concentré. En cas de projection : rincer \SI{15}{min} à l'eau.
\end{itemize}
\end{experience}

\begin{exoresolu}[Exercice type Bac : Dilution et pH]
\textbf{Énoncé :} On souhaite préparer \SI{250,0}{mL} d'une solution de soude \ce{NaOH} de pH $= 12,00$ à \SI{25}{\celsius} à partir d'une solution mère de concentration molaire $C_m = \SI{0,500}{mol.L^{-1}}$. Calculer le volume $V_m$ à prélever.
\textbf{Solution détaillée :}
\begin{enumerate}
  \item pH $= 12,00 \Rightarrow \mathrm{pOH} = \mathrm{p}K_e - \mathrm{pH} = 14,00 - 12,00 = 2,00$.
  \item $[\ce{HO-}]_f = 10^{-\mathrm{pOH}} = 10^{-2,00} = \SI{0,0100}{mol.L^{-1}} = C_f$ (base forte).
  \item Relation de dilution : $C_m V_m = C_f V_f$.
  \item $V_m = \dfrac{C_f V_f}{C_m} = \dfrac{0,0100 \times 250,0}{0,500} = \SI{5,00}{mL}$.
  \item \textbf{Réponse} : Prélever \SI{5,00}{mL} de solution mère à la pipette jaugée \SI{5}{mL}, verser dans fiole jaugée \SI{250}{mL}, compléter à l'eau distillée.
\end{enumerate}
\end{exoresolu}

\coursec{Synthèse et Entraînement Bac}

\begin{aretenir}[Checklist finale : Suis-je prêt pour l'épreuve ?]
\begin{itemize}
  \item[$\square$] Je dessine la molécule d'eau (géométrie, dipôle, liaisons H) et j'explique son pouvoir solvant.
  \item[$\square$] Je définis acide/base Brønsted, j'identifie les couples conjugués (acide = base + \ce{H+}).
  \item[$\square$] J'écris l'autoprotolyse : \ce{2H2O <=> H3O+ + HO-} et $K_e = [\ce{H3O+}][\ce{HO-}] = 10^{-14}$ (\SI{25}{\celsius}).
  \item[$\square$] Je calcule pH, pOH, $[\ce{H3O+}]$, $[\ce{HO-}]$ dans les deux sens sans erreur de calculatrice.
  \item[$\square$] Je classe les espèces (majoritaire/minoritaire/ultra-minoritaire) et justifie les approximations.
  \item[$\square$] Je maîtrise la dilution $C_iV_i = C_fV_f$ : calcul, choix verrerie (pipette/fiole jaugées), protocole (rinçage, ménisque, homogénéisation).
  \item[$\square$] Je décris la mesure du pH (pH-mètre : étalonnage 2 points ; papier/indicateurs : limites).
  \item[$\square$] Je cite des exemples concrets camerounais (bili-bili, bissap, savon, sang, eau de source).
  \item[$\square$] Je vérifie toujours : pH sans unité, concentrations en \si{mol.L^{-1}}, cohérence acide/basique.
\end{itemize}
\end{aretenir}

\begin{exercice}[Entraînement Bac TI - 20 min]
\begin{enumerate}
  \item L'eau de mer prélevée à Kribi a une concentration en ions hydroxyde $[\ce{HO-}] = \SI{1,2 \times 10^{-6}}{mol.L^{-1}}$ à \SI{25}{\celsius}. Calculer son pH et conclure sur son caractère.
  \item On dilue \SI{20,0}{mL} d'une solution d'acide nitrique \ce{HNO3} \SI{0,200}{mol.L^{-1}} à \SI{200,0}{mL}. Calculer le pH de la solution fille.
  \item Une solution de pH $= 10,50$ est mélangée à un volume égal d'eau. Quel est le nouveau pH ? (On néglige la variation de $K_e$).
  \item \textbf{QCM} : À \SI{25}{\celsius}, une solution aqueuse a $[\ce{H3O+}] = \SI{2,0 \times 10^{-4}}{mol.L^{-1}}$. Quelle affirmation est fausse ?
  \begin{itemize}
    \item A. $[\ce{HO-}] = \SI{5,0 \times 10^{-11}}{mol.L^{-1}}$
    \item B. Le pH vaut 3,70
    \item C. La solution est basique
    \item D. $[\ce{H3O+}] > [\ce{HO-}]$
  \end{itemize}
\end{enumerate}
\end{exercice}

\begin{corrige}[Exercice Entraînement Bac]
\begin{enumerate}
  \item $[\ce{H3O+}] = \frac{K_e}{[\ce{HO-}]} = \frac{10^{-14}}{1,2 \times 10^{-6}} = \SI{8,33 \times 10^{-9}}{mol.L^{-1}}$.
  $\mathrm{pH} = -\log(8,33 \times 10^{-9}) = 8,08$.
  pH $> 7 \Rightarrow$ \textbf{Solution basique}. (Normal : eau de mer riche en \ce{HCO3-}/\ce{CO3^2-}).
  \item Dilution : $C_f = C_i \frac{V_i}{V_f} = 0,200 \times \frac{20,0}{200,0} = \SI{0,0200}{mol.L^{-1}}$.
  \ce{HNO3} acide fort $\Rightarrow [\ce{H3O+}] = C_f = \SI{0,0200}{mol.L^{-1}}$.
  $\mathrm{pH} = -\log(0,0200) = 1,70$.
  \item Mélange volume égal $\Rightarrow$ dilution facteur 2. $[\ce{H3O+}]$ divisé par 2.
  Initial : $[\ce{H3O+}]_i = 10^{-10,50}$.
  Final : $[\ce{H3O+}]_f = 0,5 \times 10^{-10,50} = 5,0 \times 10^{-11}$.
  $\mathrm{pH}_f = -\log(5,0 \times 10^{-11}) = 10,30$.
  \emph{Astuce} : $\Delta\mathrm{pH} = \log 2 \approx 0,30$. pH augmente de 0,30 (dilution base).
  \item \textbf{Réponse C}.
  Vérif : $[\ce{HO-}] = 10^{-14}/(2,0 \times 10^{-4}) = 5,0 \times 10^{-11}$ (A vrai).
  $\mathrm{pH} = -\log(2,0 \times 10^{-4}) = 4 - \log 2 = 4 - 0,30 = 3,70$ (B vrai).
  $[\ce{H3O+}] > [\ce{HO-}] \Rightarrow$ Acide (C faux : dit basique).
  D vrai.
\end{enumerate}
\end{corrige}

\findecours

\end{document}
